\documentclass[11pt,a4paper]{article}
\usepackage[a4paper,margin=20mm]{geometry}
\usepackage{fontspec}
\setmainfont{Noto Serif}
\setsansfont{Noto Sans}
\usepackage{amsmath,amssymb}
\usepackage{booktabs,tabularx,array}
\usepackage[table]{xcolor}
\usepackage{hyperref}
\definecolor{primary}{RGB}{27,79,114}
\definecolor{accent}{RGB}{40,116,166}
\definecolor{rowtint}{RGB}{235,245,251}
\hypersetup{colorlinks=true,linkcolor=primary,urlcolor=primary}
\setlength{\parskip}{5pt plus 2pt minus 1pt}
\setlength{\parindent}{0pt}
\begin{document}
\begin{center}
  {\LARGE\bfseries\color{primary} PHIẾU HỌC TẬP TRẮC NGHIỆM TOÁN 11}\\[4pt]
  {\normalsize\itshape\color{accent} Bộ sách: Kết nối tri thức với cuộc sống (Tập 1 \& Tập 2 – 9 Chương)}\\[4pt]
  {\small Cấu trúc mỗi chương: 20 câu Dễ (Câu 1--20) $\bullet$ 20 câu Trung bình (Câu 21--40) $\bullet$ 10 câu Khó (Câu 41--50) -- Tổng 50 câu/chương}
\end{center}
\vspace{6pt}\hrule\vspace{10pt}
\section*{\color{primary} CHƯƠNG I: HÀM SỐ LƯỢNG GIÁC VÀ PHƯƠNG TRÌNH LƯỢNG GIÁC}
\textit{Trọng tâm kiến thức SGK: Giá trị lượng giác của góc lượng giác, công thức lượng giác, hàm số lượng giác và phương trình lượng giác cơ bản.}\par\vspace{4pt}
\subsection*{\color{accent} Phần I. Mức độ Dễ (Nhận biết -- Câu 1 đến Câu 20)}
\textbf{\color{primary} Câu 1.} Đổi số đo của góc 30° sang đơn vị radian ta được:\par
\hspace*{1.5em}\textbf{A.} $\pi$/6\quad\quad \textbf{B.} 2$\pi$/3\quad\quad \textbf{C.} $\pi$/5\quad\quad \textbf{D.} 3$\pi$/4\par\vspace{2pt}
\textbf{\color{primary} Câu 2.} Thực hiện yêu cầu nhận biết sau: Tập giá trị T của hàm số y = sin x trên $\mathbb{R}$ là.\par
\hspace*{1.5em}\textbf{A.} T = [0; 1]\quad\quad \textbf{B.} T = $\mathbb{R}$\quad\quad \textbf{C.} T = (-1; 1)\quad\quad \textbf{D.} T = [-1; 1]\par\vspace{2pt}
\textbf{\color{primary} Câu 3.} Trong bài học thuộc Chương I, xét nội dung: Chu kì tuần hoàn nhỏ nhất T của hàm số y = cos x là. Đáp án đúng là:\par
\hspace*{1.5em}\textbf{A.} T = $\pi$\quad\quad \textbf{B.} T = $\pi$/2\quad\quad \textbf{C.} T = 2$\pi$\quad\quad \textbf{D.} T = 4$\pi$\par\vspace{2pt}
\textbf{\color{primary} Câu 4.} Dựa vào kiến thức chuẩn SGK của Chương I (Hàm số lượng giác và phương trình lượng giác), đổi số đo góc lượng giác $\alpha$ = 5$\pi$/6 (rad) sang đơn vị độ ta được:\par
\hspace*{1.5em}\textbf{A.} 300°\quad\quad \textbf{B.} 75°\quad\quad \textbf{C.} 225°\quad\quad \textbf{D.} 150°\par\vspace{2pt}
\textbf{\color{primary} Câu 5.} Chọn kết quả đúng nhất khi xét: Đổi số đo của góc 90° sang đơn vị radian ta được:\par
\hspace*{1.5em}\textbf{A.} 2$\pi$/3\quad\quad \textbf{B.} $\pi$/5\quad\quad \textbf{C.} 3$\pi$/4\quad\quad \textbf{D.} $\pi$/2\par\vspace{2pt}
\textbf{\color{primary} Câu 6.} Khi làm bài tập cơ bản về hàm số lượng giác và phương trình lượng giác, một học sinh gặp câu hỏi: Tập giá trị T của hàm số y = sin x trên $\mathbb{R}$ là. Kết quả nào dưới đây đúng?\par
\hspace*{1.5em}\textbf{A.} T = (-1; 1)\quad\quad \textbf{B.} T = $\mathbb{R}$\quad\quad \textbf{C.} T = [0; 1]\quad\quad \textbf{D.} T = [-1; 1]\par\vspace{2pt}
\textbf{\color{primary} Câu 7.} Xác định giá trị hoặc khẳng định đúng cho bài toán sau: Chu kì tuần hoàn nhỏ nhất T của hàm số y = cos x là:\par
\hspace*{1.5em}\textbf{A.} T = $\pi$/2\quad\quad \textbf{B.} T = 4$\pi$\quad\quad \textbf{C.} T = $\pi$\quad\quad \textbf{D.} T = 2$\pi$\par\vspace{2pt}
\textbf{\color{primary} Câu 8.} Áp dụng trực tiếp định nghĩa và quy tắc trong Chương I: Đổi số đo góc lượng giác $\alpha$ = 9$\pi$/6 (rad) sang đơn vị độ ta được:\par
\hspace*{1.5em}\textbf{A.} 270°\quad\quad \textbf{B.} 540°\quad\quad \textbf{C.} 405°\quad\quad \textbf{D.} 135°\par\vspace{2pt}
\textbf{\color{primary} Câu 9.} Kiểm tra nhanh kiến thức nền tảng: Đổi số đo của góc 60° sang đơn vị radian ta được:\par
\hspace*{1.5em}\textbf{A.} $\pi$/5\quad\quad \textbf{B.} 2$\pi$/3\quad\quad \textbf{C.} 3$\pi$/4\quad\quad \textbf{D.} $\pi$/3\par\vspace{2pt}
\textbf{\color{primary} Câu 10.} Cho dữ kiện nhận biết sau đây: Tập giá trị T của hàm số y = sin x trên $\mathbb{R}$ là. Phương án chính xác là:\par
\hspace*{1.5em}\textbf{A.} T = [0; 1]\quad\quad \textbf{B.} T = [-1; 1]\quad\quad \textbf{C.} T = $\mathbb{R}$\quad\quad \textbf{D.} T = (-1; 1)\par\vspace{2pt}
\textbf{\color{primary} Câu 11.} Hãy chỉ ra kết quả đúng trong bốn lựa chọn A, B, C, D đối với yêu cầu: Chu kì tuần hoàn nhỏ nhất T của hàm số y = cos x là:\par
\hspace*{1.5em}\textbf{A.} T = $\pi$/2\quad\quad \textbf{B.} T = $\pi$\quad\quad \textbf{C.} T = 2$\pi$\quad\quad \textbf{D.} T = 4$\pi$\par\vspace{2pt}
\textbf{\color{primary} Câu 12.} Câu hỏi ôn tập lý thuyết và kỹ năng cơ bản (Chương I): Đổi số đo góc lượng giác $\alpha$ = 13$\pi$/6 (rad) sang đơn vị độ ta được:\par
\hspace*{1.5em}\textbf{A.} 390°\quad\quad \textbf{B.} 585°\quad\quad \textbf{C.} 195°\quad\quad \textbf{D.} 780°\par\vspace{2pt}
\textbf{\color{primary} Câu 13.} Theo nội dung sách giáo khoa Kết nối tri thức với cuộc sống (Chương I), đổi số đo của góc 30° sang đơn vị radian ta được:\par
\hspace*{1.5em}\textbf{A.} $\pi$/5\quad\quad \textbf{B.} 2$\pi$/3\quad\quad \textbf{C.} 3$\pi$/4\quad\quad \textbf{D.} $\pi$/6\par\vspace{2pt}
\textbf{\color{primary} Câu 14.} Tìm đáp số đúng của bài toán cơ bản sau: Tập giá trị T của hàm số y = sin x trên $\mathbb{R}$ là:\par
\hspace*{1.5em}\textbf{A.} T = $\mathbb{R}$\quad\quad \textbf{B.} T = [0; 1]\quad\quad \textbf{C.} T = [-1; 1]\quad\quad \textbf{D.} T = (-1; 1)\par\vspace{2pt}
\textbf{\color{primary} Câu 15.} Xét đối tượng toán học trong Chương I: Chu kì tuần hoàn nhỏ nhất T của hàm số y = cos x là. Giá trị hoặc tính chất tương ứng là:\par
\hspace*{1.5em}\textbf{A.} T = $\pi$/2\quad\quad \textbf{B.} T = 4$\pi$\quad\quad \textbf{C.} T = $\pi$\quad\quad \textbf{D.} T = 2$\pi$\par\vspace{2pt}
\textbf{\color{primary} Câu 16.} Hoàn thành phép xác định cơ bản sau đây: Đổi số đo góc lượng giác $\alpha$ = 17$\pi$/6 (rad) sang đơn vị độ ta được:\par
\hspace*{1.5em}\textbf{A.} 1020°\quad\quad \textbf{B.} 510°\quad\quad \textbf{C.} 255°\quad\quad \textbf{D.} 765°\par\vspace{2pt}
\textbf{\color{primary} Câu 17.} Trong các khẳng định hoặc giá trị dưới đây, đâu là đáp án đúng cho câu hỏi: Đổi số đo của góc 90° sang đơn vị radian ta được?\par
\hspace*{1.5em}\textbf{A.} 3$\pi$/4\quad\quad \textbf{B.} 2$\pi$/3\quad\quad \textbf{C.} $\pi$/5\quad\quad \textbf{D.} $\pi$/2\par\vspace{2pt}
\textbf{\color{primary} Câu 18.} Bài tập khởi động Chương I: Tập giá trị T của hàm số y = sin x trên $\mathbb{R}$ là:\par
\hspace*{1.5em}\textbf{A.} T = $\mathbb{R}$\quad\quad \textbf{B.} T = [0; 1]\quad\quad \textbf{C.} T = (-1; 1)\quad\quad \textbf{D.} T = [-1; 1]\par\vspace{2pt}
\textbf{\color{primary} Câu 19.} Quan sát dữ kiện đề bài: Chu kì tuần hoàn nhỏ nhất T của hàm số y = cos x là. Hãy chọn phương án đúng:\par
\hspace*{1.5em}\textbf{A.} T = 2$\pi$\quad\quad \textbf{B.} T = 4$\pi$\quad\quad \textbf{C.} T = $\pi$\quad\quad \textbf{D.} T = $\pi$/2\par\vspace{2pt}
\textbf{\color{primary} Câu 20.} Sử dụng công thức hoặc tính chất đầu bài của Chương I để giải quyết yêu cầu: Đổi số đo góc lượng giác $\alpha$ = 21$\pi$/6 (rad) sang đơn vị độ ta được:\par
\hspace*{1.5em}\textbf{A.} 945°\quad\quad \textbf{B.} 315°\quad\quad \textbf{C.} 630°\quad\quad \textbf{D.} 1260°\par\vspace{2pt}
\subsection*{\color{accent} Phần II. Mức độ Trung bình (Thông hiểu -- Câu 21 đến Câu 40)}
\textbf{\color{primary} Câu 21.} Tìm giá trị lớn nhất M và giá trị nhỏ nhất m của hàm số y = 2$\cdot$cos 2x + 5:\par
\hspace*{1.5em}\textbf{A.} M = 9; m = 5\quad\quad \textbf{B.} M = 5; m = 3\quad\quad \textbf{C.} M = 7; m = 3\quad\quad \textbf{D.} M = 2; m = -2\par\vspace{2pt}
\textbf{\color{primary} Câu 22.} Vận dụng quy tắc trong Chương I (Hàm số lượng giác và phương trình lượng giác) để giải bài toán: Họ nghiệm của phương trình lượng giác cơ bản sin x = 1/2 là:\par
\begin{itemize}\setlength{\itemsep}{1pt}\setlength{\parskip}{0pt}
  \item[] \textbf{A.} x = $\pi$/6 + k$\pi$ (k $\in$ $\mathbb{Z}$)
  \item[] \textbf{B.} x = $\pi$/3 + k2$\pi$ hoặc x = 2$\pi$/3 + k2$\pi$ (k $\in$ $\mathbb{Z}$)
  \item[] \textbf{C.} x = $\pm$$\pi$/6 + k2$\pi$ (k $\in$ $\mathbb{Z}$)
  \item[] \textbf{D.} x = $\pi$/6 + k2$\pi$ hoặc x = 5$\pi$/6 + k2$\pi$ (k $\in$ $\mathbb{Z}$)
\end{itemize}\vspace{2pt}
\textbf{\color{primary} Câu 23.} Biến đổi và tính toán theo kiến thức Chương I: Cho góc $\alpha$ thỏa mãn cos $\alpha$ = 1/3. Theo công thức nhân đôi, giá trị của cos 2$\alpha$ = 2cos$^2$$\alpha$ - 1 bằng. Kết quả thu được là:\par
\hspace*{1.5em}\textbf{A.} 2/9\quad\quad \textbf{B.} -7/9\quad\quad \textbf{C.} -1/9\quad\quad \textbf{D.} 7/9\par\vspace{2pt}
\textbf{\color{primary} Câu 24.} Xét bài toán mức độ thông hiểu sau: Tìm giá trị lớn nhất M và giá trị nhỏ nhất m của hàm số y = 5$\cdot$cos 2x + 5. Đáp án đúng là:\par
\hspace*{1.5em}\textbf{A.} M = 5; m = -5\quad\quad \textbf{B.} M = 15; m = 5\quad\quad \textbf{C.} M = 10; m = 0\quad\quad \textbf{D.} M = 5; m = 0\par\vspace{2pt}
\textbf{\color{primary} Câu 25.} Để hoàn thành phiếu học tập Chương I, hãy xác định: Họ nghiệm của phương trình lượng giác cơ bản sin x = 1/2 là:\par
\begin{itemize}\setlength{\itemsep}{1pt}\setlength{\parskip}{0pt}
  \item[] \textbf{A.} x = $\pi$/3 + k2$\pi$ hoặc x = 2$\pi$/3 + k2$\pi$ (k $\in$ $\mathbb{Z}$)
  \item[] \textbf{B.} x = $\pi$/6 + k$\pi$ (k $\in$ $\mathbb{Z}$)
  \item[] \textbf{C.} x = $\pm$$\pi$/6 + k2$\pi$ (k $\in$ $\mathbb{Z}$)
  \item[] \textbf{D.} x = $\pi$/6 + k2$\pi$ hoặc x = 5$\pi$/6 + k2$\pi$ (k $\in$ $\mathbb{Z}$)
\end{itemize}\vspace{2pt}
\textbf{\color{primary} Câu 26.} Phân tích dữ kiện bài toán: Cho góc $\alpha$ thỏa mãn cos $\alpha$ = 1/3. Theo công thức nhân đôi, giá trị của cos 2$\alpha$ = 2cos$^2$$\alpha$ - 1 bằng. Giá trị cần tìm bằng:\par
\hspace*{1.5em}\textbf{A.} -1/9\quad\quad \textbf{B.} 2/9\quad\quad \textbf{C.} -7/9\quad\quad \textbf{D.} 7/9\par\vspace{2pt}
\textbf{\color{primary} Câu 27.} Áp dụng tính chất đã học ở Chương I vào bài tập sau: Tìm giá trị lớn nhất M và giá trị nhỏ nhất m của hàm số y = 8$\cdot$cos 2x + 5:\par
\hspace*{1.5em}\textbf{A.} M = 8; m = -8\quad\quad \textbf{B.} M = 21; m = 5\quad\quad \textbf{C.} M = 5; m = -3\quad\quad \textbf{D.} M = 13; m = -3\par\vspace{2pt}
\textbf{\color{primary} Câu 28.} Thực hiện các bước suy luận hợp lý để giải quyết yêu cầu: Họ nghiệm của phương trình lượng giác cơ bản sin x = 1/2 là:\par
\begin{itemize}\setlength{\itemsep}{1pt}\setlength{\parskip}{0pt}
  \item[] \textbf{A.} x = $\pi$/6 + k2$\pi$ hoặc x = 5$\pi$/6 + k2$\pi$ (k $\in$ $\mathbb{Z}$)
  \item[] \textbf{B.} x = $\pm$$\pi$/6 + k2$\pi$ (k $\in$ $\mathbb{Z}$)
  \item[] \textbf{C.} x = $\pi$/3 + k2$\pi$ hoặc x = 2$\pi$/3 + k2$\pi$ (k $\in$ $\mathbb{Z}$)
  \item[] \textbf{D.} x = $\pi$/6 + k$\pi$ (k $\in$ $\mathbb{Z}$)
\end{itemize}\vspace{2pt}
\textbf{\color{primary} Câu 29.} Cho bài toán tự luyện (Hàm số lượng giác và phương trình lượng giác): Cho góc $\alpha$ thỏa mãn cos $\alpha$ = 1/3. Theo công thức nhân đôi, giá trị của cos 2$\alpha$ = 2cos$^2$$\alpha$ - 1 bằng. Phương án nào sau đây chính xác?\par
\hspace*{1.5em}\textbf{A.} -7/9\quad\quad \textbf{B.} -1/9\quad\quad \textbf{C.} 2/9\quad\quad \textbf{D.} 7/9\par\vspace{2pt}
\textbf{\color{primary} Câu 30.} Kết hợp các phép biến đổi trong Chương I, hãy tìm kết quả của bài toán: Tìm giá trị lớn nhất M và giá trị nhỏ nhất m của hàm số y = 11$\cdot$cos 2x + 5:\par
\hspace*{1.5em}\textbf{A.} M = 16; m = -6\quad\quad \textbf{B.} M = 11; m = -11\quad\quad \textbf{C.} M = 5; m = -6\quad\quad \textbf{D.} M = 27; m = 5\par\vspace{2pt}
\textbf{\color{primary} Câu 31.} Một bài kiểm tra định kỳ nội dung Chương I nêu câu hỏi: Họ nghiệm của phương trình lượng giác cơ bản sin x = 1/2 là. Hãy chọn đáp số đúng:\par
\begin{itemize}\setlength{\itemsep}{1pt}\setlength{\parskip}{0pt}
  \item[] \textbf{A.} x = $\pi$/6 + k2$\pi$ hoặc x = 5$\pi$/6 + k2$\pi$ (k $\in$ $\mathbb{Z}$)
  \item[] \textbf{B.} x = $\pi$/3 + k2$\pi$ hoặc x = 2$\pi$/3 + k2$\pi$ (k $\in$ $\mathbb{Z}$)
  \item[] \textbf{C.} x = $\pm$$\pi$/6 + k2$\pi$ (k $\in$ $\mathbb{Z}$)
  \item[] \textbf{D.} x = $\pi$/6 + k$\pi$ (k $\in$ $\mathbb{Z}$)
\end{itemize}\vspace{2pt}
\textbf{\color{primary} Câu 32.} Tìm lời giải đúng cho bài toán thông hiểu: Cho góc $\alpha$ thỏa mãn cos $\alpha$ = 1/3. Theo công thức nhân đôi, giá trị của cos 2$\alpha$ = 2cos$^2$$\alpha$ - 1 bằng:\par
\hspace*{1.5em}\textbf{A.} -7/9\quad\quad \textbf{B.} 7/9\quad\quad \textbf{C.} -1/9\quad\quad \textbf{D.} 2/9\par\vspace{2pt}
\textbf{\color{primary} Câu 33.} Sử dụng mối liên hệ giữa các đại lượng trong Chương I: Tìm giá trị lớn nhất M và giá trị nhỏ nhất m của hàm số y = 14$\cdot$cos 2x + 5:\par
\hspace*{1.5em}\textbf{A.} M = 14; m = -14\quad\quad \textbf{B.} M = 19; m = -9\quad\quad \textbf{C.} M = 5; m = -9\quad\quad \textbf{D.} M = 33; m = 5\par\vspace{2pt}
\textbf{\color{primary} Câu 34.} Hãy tính toán và chọn phương án đúng đối với bài tập: Họ nghiệm của phương trình lượng giác cơ bản sin x = 1/2 là:\par
\begin{itemize}\setlength{\itemsep}{1pt}\setlength{\parskip}{0pt}
  \item[] \textbf{A.} x = $\pi$/3 + k2$\pi$ hoặc x = 2$\pi$/3 + k2$\pi$ (k $\in$ $\mathbb{Z}$)
  \item[] \textbf{B.} x = $\pm$$\pi$/6 + k2$\pi$ (k $\in$ $\mathbb{Z}$)
  \item[] \textbf{C.} x = $\pi$/6 + k2$\pi$ hoặc x = 5$\pi$/6 + k2$\pi$ (k $\in$ $\mathbb{Z}$)
  \item[] \textbf{D.} x = $\pi$/6 + k$\pi$ (k $\in$ $\mathbb{Z}$)
\end{itemize}\vspace{2pt}
\textbf{\color{primary} Câu 35.} Dựa trên các tính chất trọng tâm của Chương I, xét bài toán: Cho góc $\alpha$ thỏa mãn cos $\alpha$ = 1/3. Theo công thức nhân đôi, giá trị của cos 2$\alpha$ = 2cos$^2$$\alpha$ - 1 bằng:\par
\hspace*{1.5em}\textbf{A.} 7/9\quad\quad \textbf{B.} -1/9\quad\quad \textbf{C.} -7/9\quad\quad \textbf{D.} 2/9\par\vspace{2pt}
\textbf{\color{primary} Câu 36.} Thực hành giải toán (Hàm số lượng giác và phương trình lượng giác): Tìm giá trị lớn nhất M và giá trị nhỏ nhất m của hàm số y = 17$\cdot$cos 2x + 5. Đáp án nào dưới đây thỏa mãn?\par
\hspace*{1.5em}\textbf{A.} M = 17; m = -17\quad\quad \textbf{B.} M = 22; m = -12\quad\quad \textbf{C.} M = 5; m = -12\quad\quad \textbf{D.} M = 39; m = 5\par\vspace{2pt}
\textbf{\color{primary} Câu 37.} Kiểm tra kỹ năng biến đổi và suy luận qua bài toán: Họ nghiệm của phương trình lượng giác cơ bản sin x = 1/2 là:\par
\begin{itemize}\setlength{\itemsep}{1pt}\setlength{\parskip}{0pt}
  \item[] \textbf{A.} x = $\pi$/6 + k2$\pi$ hoặc x = 5$\pi$/6 + k2$\pi$ (k $\in$ $\mathbb{Z}$)
  \item[] \textbf{B.} x = $\pi$/3 + k2$\pi$ hoặc x = 2$\pi$/3 + k2$\pi$ (k $\in$ $\mathbb{Z}$)
  \item[] \textbf{C.} x = $\pm$$\pi$/6 + k2$\pi$ (k $\in$ $\mathbb{Z}$)
  \item[] \textbf{D.} x = $\pi$/6 + k$\pi$ (k $\in$ $\mathbb{Z}$)
\end{itemize}\vspace{2pt}
\textbf{\color{primary} Câu 38.} Cho giả thiết thuộc phạm vi Chương I: Cho góc $\alpha$ thỏa mãn cos $\alpha$ = 1/3. Theo công thức nhân đôi, giá trị của cos 2$\alpha$ = 2cos$^2$$\alpha$ - 1 bằng. Kết luận đúng là:\par
\hspace*{1.5em}\textbf{A.} 2/9\quad\quad \textbf{B.} -1/9\quad\quad \textbf{C.} 7/9\quad\quad \textbf{D.} -7/9\par\vspace{2pt}
\textbf{\color{primary} Câu 39.} Rút gọn hoặc tính giá trị theo yêu cầu sau của Chương I: Tìm giá trị lớn nhất M và giá trị nhỏ nhất m của hàm số y = 20$\cdot$cos 2x + 5:\par
\hspace*{1.5em}\textbf{A.} M = 25; m = -15\quad\quad \textbf{B.} M = 45; m = 5\quad\quad \textbf{C.} M = 5; m = -15\quad\quad \textbf{D.} M = 20; m = -20\par\vspace{2pt}
\textbf{\color{primary} Câu 40.} Xác định phương án đúng sau khi thực hiện phép tính / suy luận: Họ nghiệm của phương trình lượng giác cơ bản sin x = 1/2 là:\par
\begin{itemize}\setlength{\itemsep}{1pt}\setlength{\parskip}{0pt}
  \item[] \textbf{A.} x = $\pm$$\pi$/6 + k2$\pi$ (k $\in$ $\mathbb{Z}$)
  \item[] \textbf{B.} x = $\pi$/3 + k2$\pi$ hoặc x = 2$\pi$/3 + k2$\pi$ (k $\in$ $\mathbb{Z}$)
  \item[] \textbf{C.} x = $\pi$/6 + k$\pi$ (k $\in$ $\mathbb{Z}$)
  \item[] \textbf{D.} x = $\pi$/6 + k2$\pi$ hoặc x = 5$\pi$/6 + k2$\pi$ (k $\in$ $\mathbb{Z}$)
\end{itemize}\vspace{2pt}
\subsection*{\color{accent} Phần III. Mức độ Khó (Vận dụng \& Vận dụng cao -- Câu 41 đến Câu 50)}
\textbf{\color{primary} Câu 41.} Số nghiệm của phương trình cos x = 1/2 thuộc đoạn [0; 2$\pi$] là:\par
\hspace*{1.5em}\textbf{A.} 1 nghiệm\quad\quad \textbf{B.} 2 nghiệm\quad\quad \textbf{C.} 3 nghiệm\quad\quad \textbf{D.} 4 nghiệm\par\vspace{2pt}
\textbf{\color{primary} Câu 42.} Bài toán vận dụng (Chương I – Hàm số lượng giác và phương trình lượng giác): Tìm tất cả các giá trị thực của tham số m để phương trình sin x = m - 3 có nghiệm thực. Kết quả đúng là:\par
\hspace*{1.5em}\textbf{A.} -1 $\le$ m $\le$ 1\quad\quad \textbf{B.} 1 $\le$ m $\le$ 5\quad\quad \textbf{C.} 2 $\le$ m $\le$ 4\quad\quad \textbf{D.} m $\ge$ 2\par\vspace{2pt}
\textbf{\color{primary} Câu 43.} Vận dụng kiến thức nâng cao trong phạm vi Chương I để giải bài toán: Số nghiệm của phương trình cos x = 1/2 thuộc đoạn [0; 2$\pi$] là:\par
\hspace*{1.5em}\textbf{A.} 1 nghiệm\quad\quad \textbf{B.} 4 nghiệm\quad\quad \textbf{C.} 2 nghiệm\quad\quad \textbf{D.} 3 nghiệm\par\vspace{2pt}
\textbf{\color{primary} Câu 44.} Phân tích và giải bài toán tổng hợp sau của Chương I: Tìm tất cả các giá trị thực của tham số m để phương trình sin x = m - 5 có nghiệm thực:\par
\hspace*{1.5em}\textbf{A.} -1 $\le$ m $\le$ 1\quad\quad \textbf{B.} m $\ge$ 4\quad\quad \textbf{C.} 4 $\le$ m $\le$ 6\quad\quad \textbf{D.} 3 $\le$ m $\le$ 7\par\vspace{2pt}
\textbf{\color{primary} Câu 45.} Thử thách tư duy trong phạm vi kiến thức Chương I: Số nghiệm của phương trình cos x = 1/2 thuộc đoạn [0; 2$\pi$] là. Đáp số chính xác là:\par
\hspace*{1.5em}\textbf{A.} 1 nghiệm\quad\quad \textbf{B.} 3 nghiệm\quad\quad \textbf{C.} 4 nghiệm\quad\quad \textbf{D.} 2 nghiệm\par\vspace{2pt}
\textbf{\color{primary} Câu 46.} Sử dụng phương pháp giải toán của Chương I để tìm đáp án cho bài toán: Tìm tất cả các giá trị thực của tham số m để phương trình sin x = m - 7 có nghiệm thực:\par
\hspace*{1.5em}\textbf{A.} m $\ge$ 6\quad\quad \textbf{B.} 6 $\le$ m $\le$ 8\quad\quad \textbf{C.} -1 $\le$ m $\le$ 1\quad\quad \textbf{D.} 5 $\le$ m $\le$ 9\par\vspace{2pt}
\textbf{\color{primary} Câu 47.} Xét bài toán vận dụng thực tiễn / nâng cao thuộc Chương I: Số nghiệm của phương trình cos x = 1/2 thuộc đoạn [0; 2$\pi$] là:\par
\hspace*{1.5em}\textbf{A.} 2 nghiệm\quad\quad \textbf{B.} 4 nghiệm\quad\quad \textbf{C.} 1 nghiệm\quad\quad \textbf{D.} 3 nghiệm\par\vspace{2pt}
\textbf{\color{primary} Câu 48.} Tìm kết quả cuối cùng của bài toán phân loại học sinh khá giỏi (Chương I): Tìm tất cả các giá trị thực của tham số m để phương trình sin x = m - 9 có nghiệm thực:\par
\hspace*{1.5em}\textbf{A.} 7 $\le$ m $\le$ 11\quad\quad \textbf{B.} -1 $\le$ m $\le$ 1\quad\quad \textbf{C.} 8 $\le$ m $\le$ 10\quad\quad \textbf{D.} m $\ge$ 8\par\vspace{2pt}
\textbf{\color{primary} Câu 49.} Kết hợp các tính chất đã học trong Chương I, hãy giải quyết bài toán: Số nghiệm của phương trình cos x = 1/2 thuộc đoạn [0; 2$\pi$] là:\par
\hspace*{1.5em}\textbf{A.} 1 nghiệm\quad\quad \textbf{B.} 3 nghiệm\quad\quad \textbf{C.} 4 nghiệm\quad\quad \textbf{D.} 2 nghiệm\par\vspace{2pt}
\textbf{\color{primary} Câu 50.} Câu hỏi vận dụng cao cuối Chương I: Tìm tất cả các giá trị thực của tham số m để phương trình sin x = m - 11 có nghiệm thực. Phương án đúng là:\par
\hspace*{1.5em}\textbf{A.} -1 $\le$ m $\le$ 1\quad\quad \textbf{B.} 9 $\le$ m $\le$ 13\quad\quad \textbf{C.} m $\ge$ 10\quad\quad \textbf{D.} 10 $\le$ m $\le$ 12\par\vspace{2pt}
\subsection*{\color{primary} Bảng đáp án Chương I: Hàm số lượng giác và phương trình lượng giác}
\begin{center}
\renewcommand{\arraystretch}{1.25}
\begin{tabularx}{\textwidth}{*{10}{>{\centering\arraybackslash}X}}
\toprule
\textbf{\color{primary} Cột 1} & \textbf{\color{primary} Cột 2} & \textbf{\color{primary} Cột 3} & \textbf{\color{primary} Cột 4} & \textbf{\color{primary} Cột 5} & \textbf{\color{primary} Cột 6} & \textbf{\color{primary} Cột 7} & \textbf{\color{primary} Cột 8} & \textbf{\color{primary} Cột 9} & \textbf{\color{primary} Cột 10} \\
\midrule
\textbf{1.A} & \textbf{2.D} & \textbf{3.C} & \textbf{4.D} & \textbf{5.D} & \textbf{6.D} & \textbf{7.D} & \textbf{8.A} & \textbf{9.D} & \textbf{10.B} \\
\rowcolor{rowtint} \textbf{11.C} & \textbf{12.A} & \textbf{13.D} & \textbf{14.C} & \textbf{15.D} & \textbf{16.B} & \textbf{17.D} & \textbf{18.D} & \textbf{19.A} & \textbf{20.C} \\
\textbf{21.C} & \textbf{22.D} & \textbf{23.B} & \textbf{24.C} & \textbf{25.D} & \textbf{26.C} & \textbf{27.D} & \textbf{28.A} & \textbf{29.A} & \textbf{30.A} \\
\rowcolor{rowtint} \textbf{31.A} & \textbf{32.A} & \textbf{33.B} & \textbf{34.C} & \textbf{35.C} & \textbf{36.B} & \textbf{37.A} & \textbf{38.D} & \textbf{39.A} & \textbf{40.D} \\
\textbf{41.B} & \textbf{42.C} & \textbf{43.C} & \textbf{44.C} & \textbf{45.D} & \textbf{46.B} & \textbf{47.A} & \textbf{48.C} & \textbf{49.D} & \textbf{50.D} \\
\bottomrule
\end{tabularx}
\end{center}\vspace{10pt}
\section*{\color{primary} CHƯƠNG II: DÃY SỐ. CẤP SỐ CỘNG VÀ CẤP SỐ NHÂN}
\textit{Trọng tâm kiến thức SGK: Dãy số, tính tăng giảm và bị chặn của dãy số, cấp số cộng, cấp số nhân (số hạng tổng quát và tổng n số hạng đầu).}\par\vspace{4pt}
\subsection*{\color{accent} Phần I. Mức độ Dễ (Nhận biết -- Câu 1 đến Câu 20)}
\textbf{\color{primary} Câu 1.} Cho dãy số ($u_n$) có số hạng tổng quát $u_n$ = 3n + 2. Số hạng thứ 5 (u$_5$) của dãy số bằng:\par
\hspace*{1.5em}\textbf{A.} 14\quad\quad \textbf{B.} 17\quad\quad \textbf{C.} 20\quad\quad \textbf{D.} 7\par\vspace{2pt}
\textbf{\color{primary} Câu 2.} Thực hiện yêu cầu nhận biết sau: Cho cấp số cộng ($u_n$) có số hạng đầu u$_1$ = 3 và công sai d = 4. Số hạng u$_6$ của cấp số cộng bằng.\par
\hspace*{1.5em}\textbf{A.} 23\quad\quad \textbf{B.} 19\quad\quad \textbf{C.} 27\quad\quad \textbf{D.} 12\par\vspace{2pt}
\textbf{\color{primary} Câu 3.} Trong bài học thuộc Chương II, xét nội dung: Cho cấp số nhân ($u_n$) có u$_1$ = 4 và công bội q = 2. Số hạng u$_4$ của cấp số nhân bằng. Đáp án đúng là:\par
\hspace*{1.5em}\textbf{A.} 64\quad\quad \textbf{B.} 32\quad\quad \textbf{C.} 24\quad\quad \textbf{D.} 16\par\vspace{2pt}
\textbf{\color{primary} Câu 4.} Dựa vào kiến thức chuẩn SGK của Chương II (Dãy số. Cấp số cộng và cấp số nhân), cho dãy số ($u_n$) có số hạng tổng quát $u_n$ = 3n + 5. Số hạng thứ 5 (u$_5$) của dãy số bằng:\par
\hspace*{1.5em}\textbf{A.} 17\quad\quad \textbf{B.} 10\quad\quad \textbf{C.} 20\quad\quad \textbf{D.} 23\par\vspace{2pt}
\textbf{\color{primary} Câu 5.} Chọn kết quả đúng nhất khi xét: Cho cấp số cộng ($u_n$) có số hạng đầu u$_1$ = 6 và công sai d = 4. Số hạng u$_6$ của cấp số cộng bằng:\par
\hspace*{1.5em}\textbf{A.} 24\quad\quad \textbf{B.} 26\quad\quad \textbf{C.} 30\quad\quad \textbf{D.} 22\par\vspace{2pt}
\textbf{\color{primary} Câu 6.} Khi làm bài tập cơ bản về dãy số. cấp số cộng và cấp số nhân, một học sinh gặp câu hỏi: Cho cấp số nhân ($u_n$) có u$_1$ = 7 và công bội q = 2. Số hạng u$_4$ của cấp số nhân bằng. Kết quả nào dưới đây đúng?\par
\hspace*{1.5em}\textbf{A.} 56\quad\quad \textbf{B.} 112\quad\quad \textbf{C.} 28\quad\quad \textbf{D.} 42\par\vspace{2pt}
\textbf{\color{primary} Câu 7.} Xác định giá trị hoặc khẳng định đúng cho bài toán sau: Cho dãy số ($u_n$) có số hạng tổng quát $u_n$ = 3n + 8. Số hạng thứ 5 (u$_5$) của dãy số bằng:\par
\hspace*{1.5em}\textbf{A.} 23\quad\quad \textbf{B.} 26\quad\quad \textbf{C.} 13\quad\quad \textbf{D.} 20\par\vspace{2pt}
\textbf{\color{primary} Câu 8.} Áp dụng trực tiếp định nghĩa và quy tắc trong Chương II: Cho cấp số cộng ($u_n$) có số hạng đầu u$_1$ = 9 và công sai d = 4. Số hạng u$_6$ của cấp số cộng bằng:\par
\hspace*{1.5em}\textbf{A.} 33\quad\quad \textbf{B.} 29\quad\quad \textbf{C.} 25\quad\quad \textbf{D.} 36\par\vspace{2pt}
\textbf{\color{primary} Câu 9.} Kiểm tra nhanh kiến thức nền tảng: Cho cấp số nhân ($u_n$) có u$_1$ = 10 và công bội q = 2. Số hạng u$_4$ của cấp số nhân bằng:\par
\hspace*{1.5em}\textbf{A.} 60\quad\quad \textbf{B.} 40\quad\quad \textbf{C.} 80\quad\quad \textbf{D.} 160\par\vspace{2pt}
\textbf{\color{primary} Câu 10.} Cho dữ kiện nhận biết sau đây: Cho dãy số ($u_n$) có số hạng tổng quát $u_n$ = 3n + 11. Số hạng thứ 5 (u$_5$) của dãy số bằng. Phương án chính xác là:\par
\hspace*{1.5em}\textbf{A.} 29\quad\quad \textbf{B.} 26\quad\quad \textbf{C.} 16\quad\quad \textbf{D.} 23\par\vspace{2pt}
\textbf{\color{primary} Câu 11.} Hãy chỉ ra kết quả đúng trong bốn lựa chọn A, B, C, D đối với yêu cầu: Cho cấp số cộng ($u_n$) có số hạng đầu u$_1$ = 12 và công sai d = 4. Số hạng u$_6$ của cấp số cộng bằng:\par
\hspace*{1.5em}\textbf{A.} 28\quad\quad \textbf{B.} 32\quad\quad \textbf{C.} 36\quad\quad \textbf{D.} 48\par\vspace{2pt}
\textbf{\color{primary} Câu 12.} Câu hỏi ôn tập lý thuyết và kỹ năng cơ bản (Chương II): Cho cấp số nhân ($u_n$) có u$_1$ = 13 và công bội q = 2. Số hạng u$_4$ của cấp số nhân bằng:\par
\hspace*{1.5em}\textbf{A.} 78\quad\quad \textbf{B.} 52\quad\quad \textbf{C.} 104\quad\quad \textbf{D.} 208\par\vspace{2pt}
\textbf{\color{primary} Câu 13.} Theo nội dung sách giáo khoa Kết nối tri thức với cuộc sống (Chương II), cho dãy số ($u_n$) có số hạng tổng quát $u_n$ = 3n + 14. Số hạng thứ 5 (u$_5$) của dãy số bằng:\par
\hspace*{1.5em}\textbf{A.} 29\quad\quad \textbf{B.} 26\quad\quad \textbf{C.} 19\quad\quad \textbf{D.} 32\par\vspace{2pt}
\textbf{\color{primary} Câu 14.} Tìm đáp số đúng của bài toán cơ bản sau: Cho cấp số cộng ($u_n$) có số hạng đầu u$_1$ = 15 và công sai d = 4. Số hạng u$_6$ của cấp số cộng bằng:\par
\hspace*{1.5em}\textbf{A.} 31\quad\quad \textbf{B.} 35\quad\quad \textbf{C.} 39\quad\quad \textbf{D.} 60\par\vspace{2pt}
\textbf{\color{primary} Câu 15.} Xét đối tượng toán học trong Chương II: Cho cấp số nhân ($u_n$) có u$_1$ = 16 và công bội q = 2. Số hạng u$_4$ của cấp số nhân bằng. Giá trị hoặc tính chất tương ứng là:\par
\hspace*{1.5em}\textbf{A.} 256\quad\quad \textbf{B.} 128\quad\quad \textbf{C.} 64\quad\quad \textbf{D.} 96\par\vspace{2pt}
\textbf{\color{primary} Câu 16.} Hoàn thành phép xác định cơ bản sau đây: Cho dãy số ($u_n$) có số hạng tổng quát $u_n$ = 3n + 17. Số hạng thứ 5 (u$_5$) của dãy số bằng:\par
\hspace*{1.5em}\textbf{A.} 32\quad\quad \textbf{B.} 35\quad\quad \textbf{C.} 29\quad\quad \textbf{D.} 22\par\vspace{2pt}
\textbf{\color{primary} Câu 17.} Trong các khẳng định hoặc giá trị dưới đây, đâu là đáp án đúng cho câu hỏi: Cho cấp số cộng ($u_n$) có số hạng đầu u$_1$ = 18 và công sai d = 4. Số hạng u$_6$ của cấp số cộng bằng?\par
\hspace*{1.5em}\textbf{A.} 42\quad\quad \textbf{B.} 38\quad\quad \textbf{C.} 34\quad\quad \textbf{D.} 72\par\vspace{2pt}
\textbf{\color{primary} Câu 18.} Bài tập khởi động Chương II: Cho cấp số nhân ($u_n$) có u$_1$ = 19 và công bội q = 2. Số hạng u$_4$ của cấp số nhân bằng:\par
\hspace*{1.5em}\textbf{A.} 114\quad\quad \textbf{B.} 152\quad\quad \textbf{C.} 76\quad\quad \textbf{D.} 304\par\vspace{2pt}
\textbf{\color{primary} Câu 19.} Quan sát dữ kiện đề bài: Cho dãy số ($u_n$) có số hạng tổng quát $u_n$ = 3n + 20. Số hạng thứ 5 (u$_5$) của dãy số bằng. Hãy chọn phương án đúng:\par
\hspace*{1.5em}\textbf{A.} 25\quad\quad \textbf{B.} 32\quad\quad \textbf{C.} 35\quad\quad \textbf{D.} 38\par\vspace{2pt}
\textbf{\color{primary} Câu 20.} Sử dụng công thức hoặc tính chất đầu bài của Chương II để giải quyết yêu cầu: Cho cấp số cộng ($u_n$) có số hạng đầu u$_1$ = 21 và công sai d = 4. Số hạng u$_6$ của cấp số cộng bằng:\par
\hspace*{1.5em}\textbf{A.} 41\quad\quad \textbf{B.} 45\quad\quad \textbf{C.} 84\quad\quad \textbf{D.} 37\par\vspace{2pt}
\subsection*{\color{accent} Phần II. Mức độ Trung bình (Thông hiểu -- Câu 21 đến Câu 40)}
\textbf{\color{primary} Câu 21.} Cho cấp số cộng ($u_n$) có u$_1$ = 3 và công sai d = 2. Tổng 10 số hạng đầu tiên S$_1$$_0$ của cấp số cộng bằng:\par
\hspace*{1.5em}\textbf{A.} 75\quad\quad \textbf{B.} 120\quad\quad \textbf{C.} 130\quad\quad \textbf{D.} 105\par\vspace{2pt}
\textbf{\color{primary} Câu 22.} Vận dụng quy tắc trong Chương II (Dãy số. Cấp số cộng và cấp số nhân) để giải bài toán: Cho cấp số cộng ($u_n$) biết u$_2$ = 6 và u$_6$ = 18. Công sai d của cấp số cộng đó bằng:\par
\hspace*{1.5em}\textbf{A.} d = 4\quad\quad \textbf{B.} d = 2\quad\quad \textbf{C.} d = 6\quad\quad \textbf{D.} d = 3\par\vspace{2pt}
\textbf{\color{primary} Câu 23.} Biến đổi và tính toán theo kiến thức Chương II: Cho cấp số nhân ($u_n$) có u$_1$ = 3 và công bội q = 2. Tổng 6 số hạng đầu tiên S$_6$ của cấp số nhân bằng. Kết quả thu được là:\par
\hspace*{1.5em}\textbf{A.} 381\quad\quad \textbf{B.} 93\quad\quad \textbf{C.} 189\quad\quad \textbf{D.} 192\par\vspace{2pt}
\textbf{\color{primary} Câu 24.} Xét bài toán mức độ thông hiểu sau: Cho cấp số cộng ($u_n$) có u$_1$ = 6 và công sai d = 2. Tổng 10 số hạng đầu tiên S$_1$$_0$ của cấp số cộng bằng. Đáp án đúng là:\par
\hspace*{1.5em}\textbf{A.} 160\quad\quad \textbf{B.} 105\quad\quad \textbf{C.} 120\quad\quad \textbf{D.} 150\par\vspace{2pt}
\textbf{\color{primary} Câu 25.} Để hoàn thành phiếu học tập Chương II, hãy xác định: Cho cấp số cộng ($u_n$) biết u$_2$ = 9 và u$_6$ = 21. Công sai d của cấp số cộng đó bằng:\par
\hspace*{1.5em}\textbf{A.} d = 2\quad\quad \textbf{B.} d = 6\quad\quad \textbf{C.} d = 3\quad\quad \textbf{D.} d = 4\par\vspace{2pt}
\textbf{\color{primary} Câu 26.} Phân tích dữ kiện bài toán: Cho cấp số nhân ($u_n$) có u$_1$ = 3 và công bội q = 2. Tổng 6 số hạng đầu tiên S$_6$ của cấp số nhân bằng. Giá trị cần tìm bằng:\par
\hspace*{1.5em}\textbf{A.} 93\quad\quad \textbf{B.} 381\quad\quad \textbf{C.} 192\quad\quad \textbf{D.} 189\par\vspace{2pt}
\textbf{\color{primary} Câu 27.} Áp dụng tính chất đã học ở Chương II vào bài tập sau: Cho cấp số cộng ($u_n$) có u$_1$ = 9 và công sai d = 2. Tổng 10 số hạng đầu tiên S$_1$$_0$ của cấp số cộng bằng:\par
\hspace*{1.5em}\textbf{A.} 190\quad\quad \textbf{B.} 135\quad\quad \textbf{C.} 180 (1)\quad\quad \textbf{D.} 180\par\vspace{2pt}
\textbf{\color{primary} Câu 28.} Thực hiện các bước suy luận hợp lý để giải quyết yêu cầu: Cho cấp số cộng ($u_n$) biết u$_2$ = 12 và u$_6$ = 24. Công sai d của cấp số cộng đó bằng:\par
\hspace*{1.5em}\textbf{A.} d = 4\quad\quad \textbf{B.} d = 6\quad\quad \textbf{C.} d = 2\quad\quad \textbf{D.} d = 3\par\vspace{2pt}
\textbf{\color{primary} Câu 29.} Cho bài toán tự luyện (Dãy số. Cấp số cộng và cấp số nhân): Cho cấp số nhân ($u_n$) có u$_1$ = 3 và công bội q = 2. Tổng 6 số hạng đầu tiên S$_6$ của cấp số nhân bằng. Phương án nào sau đây chính xác?\par
\hspace*{1.5em}\textbf{A.} 189\quad\quad \textbf{B.} 192\quad\quad \textbf{C.} 93\quad\quad \textbf{D.} 381\par\vspace{2pt}
\textbf{\color{primary} Câu 30.} Kết hợp các phép biến đổi trong Chương II, hãy tìm kết quả của bài toán: Cho cấp số cộng ($u_n$) có u$_1$ = 12 và công sai d = 2. Tổng 10 số hạng đầu tiên S$_1$$_0$ của cấp số cộng bằng:\par
\hspace*{1.5em}\textbf{A.} 220\quad\quad \textbf{B.} 150\quad\quad \textbf{C.} 165\quad\quad \textbf{D.} 210\par\vspace{2pt}
\textbf{\color{primary} Câu 31.} Một bài kiểm tra định kỳ nội dung Chương II nêu câu hỏi: Cho cấp số cộng ($u_n$) biết u$_2$ = 15 và u$_6$ = 27. Công sai d của cấp số cộng đó bằng. Hãy chọn đáp số đúng:\par
\hspace*{1.5em}\textbf{A.} d = 4\quad\quad \textbf{B.} d = 6\quad\quad \textbf{C.} d = 3\quad\quad \textbf{D.} d = 2\par\vspace{2pt}
\textbf{\color{primary} Câu 32.} Tìm lời giải đúng cho bài toán thông hiểu: Cho cấp số nhân ($u_n$) có u$_1$ = 3 và công bội q = 2. Tổng 6 số hạng đầu tiên S$_6$ của cấp số nhân bằng:\par
\hspace*{1.5em}\textbf{A.} 381\quad\quad \textbf{B.} 192\quad\quad \textbf{C.} 93\quad\quad \textbf{D.} 189\par\vspace{2pt}
\textbf{\color{primary} Câu 33.} Sử dụng mối liên hệ giữa các đại lượng trong Chương II: Cho cấp số cộng ($u_n$) có u$_1$ = 15 và công sai d = 2. Tổng 10 số hạng đầu tiên S$_1$$_0$ của cấp số cộng bằng:\par
\hspace*{1.5em}\textbf{A.} 250\quad\quad \textbf{B.} 240\quad\quad \textbf{C.} 165\quad\quad \textbf{D.} 195\par\vspace{2pt}
\textbf{\color{primary} Câu 34.} Hãy tính toán và chọn phương án đúng đối với bài tập: Cho cấp số cộng ($u_n$) biết u$_2$ = 18 và u$_6$ = 30. Công sai d của cấp số cộng đó bằng:\par
\hspace*{1.5em}\textbf{A.} d = 3\quad\quad \textbf{B.} d = 6\quad\quad \textbf{C.} d = 4\quad\quad \textbf{D.} d = 2\par\vspace{2pt}
\textbf{\color{primary} Câu 35.} Dựa trên các tính chất trọng tâm của Chương II, xét bài toán: Cho cấp số nhân ($u_n$) có u$_1$ = 3 và công bội q = 2. Tổng 6 số hạng đầu tiên S$_6$ của cấp số nhân bằng:\par
\hspace*{1.5em}\textbf{A.} 93\quad\quad \textbf{B.} 189\quad\quad \textbf{C.} 192\quad\quad \textbf{D.} 381\par\vspace{2pt}
\textbf{\color{primary} Câu 36.} Thực hành giải toán (Dãy số. Cấp số cộng và cấp số nhân): Cho cấp số cộng ($u_n$) có u$_1$ = 18 và công sai d = 2. Tổng 10 số hạng đầu tiên S$_1$$_0$ của cấp số cộng bằng. Đáp án nào dưới đây thỏa mãn?\par
\hspace*{1.5em}\textbf{A.} 280\quad\quad \textbf{B.} 225\quad\quad \textbf{C.} 270\quad\quad \textbf{D.} 180\par\vspace{2pt}
\textbf{\color{primary} Câu 37.} Kiểm tra kỹ năng biến đổi và suy luận qua bài toán: Cho cấp số cộng ($u_n$) biết u$_2$ = 21 và u$_6$ = 33. Công sai d của cấp số cộng đó bằng:\par
\hspace*{1.5em}\textbf{A.} d = 2\quad\quad \textbf{B.} d = 6\quad\quad \textbf{C.} d = 3\quad\quad \textbf{D.} d = 4\par\vspace{2pt}
\textbf{\color{primary} Câu 38.} Cho giả thiết thuộc phạm vi Chương II: Cho cấp số nhân ($u_n$) có u$_1$ = 3 và công bội q = 2. Tổng 6 số hạng đầu tiên S$_6$ của cấp số nhân bằng. Kết luận đúng là:\par
\hspace*{1.5em}\textbf{A.} 189\quad\quad \textbf{B.} 93\quad\quad \textbf{C.} 381\quad\quad \textbf{D.} 192\par\vspace{2pt}
\textbf{\color{primary} Câu 39.} Rút gọn hoặc tính giá trị theo yêu cầu sau của Chương II: Cho cấp số cộng ($u_n$) có u$_1$ = 21 và công sai d = 2. Tổng 10 số hạng đầu tiên S$_1$$_0$ của cấp số cộng bằng:\par
\hspace*{1.5em}\textbf{A.} 255\quad\quad \textbf{B.} 195\quad\quad \textbf{C.} 300\quad\quad \textbf{D.} 310\par\vspace{2pt}
\textbf{\color{primary} Câu 40.} Xác định phương án đúng sau khi thực hiện phép tính / suy luận: Cho cấp số cộng ($u_n$) biết u$_2$ = 24 và u$_6$ = 36. Công sai d của cấp số cộng đó bằng:\par
\hspace*{1.5em}\textbf{A.} d = 6\quad\quad \textbf{B.} d = 4\quad\quad \textbf{C.} d = 3\quad\quad \textbf{D.} d = 2\par\vspace{2pt}
\subsection*{\color{accent} Phần III. Mức độ Khó (Vận dụng \& Vận dụng cao -- Câu 41 đến Câu 50)}
\textbf{\color{primary} Câu 41.} Tìm x dương để ba số 2; x; 8 theo thứ tự lập thành một cấp số nhân:\par
\hspace*{1.5em}\textbf{A.} x = 5\quad\quad \textbf{B.} x = 2\quad\quad \textbf{C.} x = 4\quad\quad \textbf{D.} x = 8\par\vspace{2pt}
\textbf{\color{primary} Câu 42.} Bài toán vận dụng (Chương II – Dãy số. Cấp số cộng và cấp số nhân): Một rạp hát có 15 hàng ghế, hàng đầu tiên có 18 ghế, mỗi hàng sau nhiều hơn hàng liền trước 2 ghế. Tổng số ghế của rạp hát là. Kết quả đúng là:\par
\hspace*{1.5em}\textbf{A.} 690 ghế\quad\quad \textbf{B.} 480 ghế\quad\quad \textbf{C.} 375 ghế\quad\quad \textbf{D.} 270 ghế\par\vspace{2pt}
\textbf{\color{primary} Câu 43.} Vận dụng kiến thức nâng cao trong phạm vi Chương II để giải bài toán: Tìm số thực dương x để ba số 4; x; 36 theo thứ tự lập thành một cấp số nhân:\par
\hspace*{1.5em}\textbf{A.} x = 8\quad\quad \textbf{B.} x = 36\quad\quad \textbf{C.} x = 12\quad\quad \textbf{D.} x = 20\par\vspace{2pt}
\textbf{\color{primary} Câu 44.} Phân tích và giải bài toán tổng hợp sau của Chương II: Tìm x dương để ba số 2; x; 8 theo thứ tự lập thành một cấp số nhân:\par
\hspace*{1.5em}\textbf{A.} x = 5\quad\quad \textbf{B.} x = 8\quad\quad \textbf{C.} x = 4\quad\quad \textbf{D.} x = 2\par\vspace{2pt}
\textbf{\color{primary} Câu 45.} Thử thách tư duy trong phạm vi kiến thức Chương II: Một rạp hát có 15 hàng ghế, hàng đầu tiên có 18 ghế, mỗi hàng sau nhiều hơn hàng liền trước 2 ghế. Tổng số ghế của rạp hát là. Đáp số chính xác là:\par
\hspace*{1.5em}\textbf{A.} 375 ghế\quad\quad \textbf{B.} 270 ghế\quad\quad \textbf{C.} 480 ghế\quad\quad \textbf{D.} 690 ghế\par\vspace{2pt}
\textbf{\color{primary} Câu 46.} Sử dụng phương pháp giải toán của Chương II để tìm đáp án cho bài toán: Tìm số thực dương x để ba số 7; x; 63 theo thứ tự lập thành một cấp số nhân:\par
\hspace*{1.5em}\textbf{A.} x = 21\quad\quad \textbf{B.} x = 63\quad\quad \textbf{C.} x = 14\quad\quad \textbf{D.} x = 35\par\vspace{2pt}
\textbf{\color{primary} Câu 47.} Xét bài toán vận dụng thực tiễn / nâng cao thuộc Chương II: Tìm x dương để ba số 2; x; 8 theo thứ tự lập thành một cấp số nhân:\par
\hspace*{1.5em}\textbf{A.} x = 5\quad\quad \textbf{B.} x = 8\quad\quad \textbf{C.} x = 4\quad\quad \textbf{D.} x = 2\par\vspace{2pt}
\textbf{\color{primary} Câu 48.} Tìm kết quả cuối cùng của bài toán phân loại học sinh khá giỏi (Chương II): Một rạp hát có 15 hàng ghế, hàng đầu tiên có 18 ghế, mỗi hàng sau nhiều hơn hàng liền trước 2 ghế. Tổng số ghế của rạp hát là:\par
\hspace*{1.5em}\textbf{A.} 270 ghế\quad\quad \textbf{B.} 690 ghế\quad\quad \textbf{C.} 375 ghế\quad\quad \textbf{D.} 480 ghế\par\vspace{2pt}
\textbf{\color{primary} Câu 49.} Kết hợp các tính chất đã học trong Chương II, hãy giải quyết bài toán: Tìm số thực dương x để ba số 10; x; 90 theo thứ tự lập thành một cấp số nhân:\par
\hspace*{1.5em}\textbf{A.} x = 20\quad\quad \textbf{B.} x = 30\quad\quad \textbf{C.} x = 50\quad\quad \textbf{D.} x = 90\par\vspace{2pt}
\textbf{\color{primary} Câu 50.} Câu hỏi vận dụng cao cuối Chương II: Tìm x dương để ba số 2; x; 8 theo thứ tự lập thành một cấp số nhân. Phương án đúng là:\par
\hspace*{1.5em}\textbf{A.} x = 8\quad\quad \textbf{B.} x = 5\quad\quad \textbf{C.} x = 2\quad\quad \textbf{D.} x = 4\par\vspace{2pt}
\subsection*{\color{primary} Bảng đáp án Chương II: Dãy số. Cấp số cộng và cấp số nhân}
\begin{center}
\renewcommand{\arraystretch}{1.25}
\begin{tabularx}{\textwidth}{*{10}{>{\centering\arraybackslash}X}}
\toprule
\textbf{\color{primary} Cột 1} & \textbf{\color{primary} Cột 2} & \textbf{\color{primary} Cột 3} & \textbf{\color{primary} Cột 4} & \textbf{\color{primary} Cột 5} & \textbf{\color{primary} Cột 6} & \textbf{\color{primary} Cột 7} & \textbf{\color{primary} Cột 8} & \textbf{\color{primary} Cột 9} & \textbf{\color{primary} Cột 10} \\
\midrule
\textbf{1.B} & \textbf{2.A} & \textbf{3.B} & \textbf{4.C} & \textbf{5.B} & \textbf{6.A} & \textbf{7.A} & \textbf{8.B} & \textbf{9.C} & \textbf{10.B} \\
\rowcolor{rowtint} \textbf{11.B} & \textbf{12.C} & \textbf{13.A} & \textbf{14.B} & \textbf{15.B} & \textbf{16.A} & \textbf{17.B} & \textbf{18.B} & \textbf{19.C} & \textbf{20.A} \\
\textbf{21.B} & \textbf{22.D} & \textbf{23.C} & \textbf{24.D} & \textbf{25.C} & \textbf{26.D} & \textbf{27.D} & \textbf{28.D} & \textbf{29.A} & \textbf{30.D} \\
\rowcolor{rowtint} \textbf{31.C} & \textbf{32.D} & \textbf{33.B} & \textbf{34.A} & \textbf{35.B} & \textbf{36.C} & \textbf{37.C} & \textbf{38.A} & \textbf{39.C} & \textbf{40.C} \\
\textbf{41.C} & \textbf{42.B} & \textbf{43.C} & \textbf{44.C} & \textbf{45.C} & \textbf{46.A} & \textbf{47.C} & \textbf{48.D} & \textbf{49.B} & \textbf{50.D} \\
\bottomrule
\end{tabularx}
\end{center}\vspace{10pt}
\section*{\color{primary} CHƯƠNG III: CÁC SỐ ĐẶC TRƯNG ĐO XU THẾ TRUNG TÂM CỦA MẪU SỐ LIỆU GHÉP NHÓM}
\textit{Trọng tâm kiến thức SGK: Mẫu số liệu ghép nhóm, giá trị đại diện của nhóm, số trung bình cộng, trung vị, tứ phân vị và mốt của mẫu số liệu ghép nhóm.}\par\vspace{4pt}
\subsection*{\color{accent} Phần I. Mức độ Dễ (Nhận biết -- Câu 1 đến Câu 20)}
\textbf{\color{primary} Câu 1.} Giá trị đại diện $c_i$ của nhóm số liệu nửa khoảng [10; 18) bằng:\par
\hspace*{1.5em}\textbf{A.} 12\quad\quad \textbf{B.} 14\quad\quad \textbf{C.} 16\quad\quad \textbf{D.} 8\par\vspace{2pt}
\textbf{\color{primary} Câu 2.} Hãy trả lời câu hỏi cơ bản sau: Trong bảng phân bố tần số ghép nhóm, nhóm chứa mốt là nhóm có đặc điểm nào sau đây?\par
\begin{itemize}\setlength{\itemsep}{1pt}\setlength{\parskip}{0pt}
  \item[] \textbf{A.} Nhóm có tần số lớn nhất
  \item[] \textbf{B.} Nhóm đầu tiên của bảng
  \item[] \textbf{C.} Nhóm có độ dài lớn nhất
  \item[] \textbf{D.} Nhóm có giá trị đại diện lớn nhất
\end{itemize}\vspace{2pt}
\textbf{\color{primary} Câu 3.} Trong bài học thuộc Chương III, xét nội dung: Cho mẫu số liệu ghép nhóm có tổng cỡ mẫu n = 56. Nhóm chứa tứ phân vị thứ nhất Q$_1$ là nhóm đầu tiên có tần số tích lũy lớn hơn hoặc bằng. Đáp án đúng là:\par
\hspace*{1.5em}\textbf{A.} 42\quad\quad \textbf{B.} 14\quad\quad \textbf{C.} 56\quad\quad \textbf{D.} 28\par\vspace{2pt}
\textbf{\color{primary} Câu 4.} Dựa vào kiến thức chuẩn SGK của Chương III (Các số đặc trưng đo xu thế trung tâm của mẫu số liệu ghép nhóm), giá trị đại diện $c_i$ của nhóm số liệu nửa khoảng [22; 30) bằng:\par
\hspace*{1.5em}\textbf{A.} 26\quad\quad \textbf{B.} 28\quad\quad \textbf{C.} 8\quad\quad \textbf{D.} 24\par\vspace{2pt}
\textbf{\color{primary} Câu 5.} Chọn phương án trả lời đúng nhất: Trong bảng phân bố tần số ghép nhóm, nhóm chứa mốt là nhóm có đặc điểm nào sau đây?\par
\begin{itemize}\setlength{\itemsep}{1pt}\setlength{\parskip}{0pt}
  \item[] \textbf{A.} Nhóm đầu tiên của bảng
  \item[] \textbf{B.} Nhóm có độ dài lớn nhất
  \item[] \textbf{C.} Nhóm có giá trị đại diện lớn nhất
  \item[] \textbf{D.} Nhóm có tần số lớn nhất
\end{itemize}\vspace{2pt}
\textbf{\color{primary} Câu 6.} Khi làm bài tập cơ bản về các số đặc trưng đo xu thế trung tâm của mẫu số liệu ghép nhóm, một học sinh gặp câu hỏi: Cho mẫu số liệu ghép nhóm có tổng cỡ mẫu n = 80. Nhóm chứa tứ phân vị thứ nhất Q$_1$ là nhóm đầu tiên có tần số tích lũy lớn hơn hoặc bằng. Kết quả nào dưới đây đúng?\par
\hspace*{1.5em}\textbf{A.} 80\quad\quad \textbf{B.} 20\quad\quad \textbf{C.} 40\quad\quad \textbf{D.} 60\par\vspace{2pt}
\textbf{\color{primary} Câu 7.} Xác định giá trị hoặc khẳng định đúng cho bài toán sau: Giá trị đại diện $c_i$ của nhóm số liệu nửa khoảng [34; 42) bằng:\par
\hspace*{1.5em}\textbf{A.} 8\quad\quad \textbf{B.} 38\quad\quad \textbf{C.} 36\quad\quad \textbf{D.} 40\par\vspace{2pt}
\textbf{\color{primary} Câu 8.} Áp dụng trực tiếp định nghĩa và quy tắc trong Chương III: Trong bảng phân bố tần số ghép nhóm, nhóm chứa mốt là nhóm có đặc điểm nào sau đây:\par
\begin{itemize}\setlength{\itemsep}{1pt}\setlength{\parskip}{0pt}
  \item[] \textbf{A.} Nhóm có độ dài lớn nhất
  \item[] \textbf{B.} Nhóm có giá trị đại diện lớn nhất
  \item[] \textbf{C.} Nhóm đầu tiên của bảng
  \item[] \textbf{D.} Nhóm có tần số lớn nhất
\end{itemize}\vspace{2pt}
\textbf{\color{primary} Câu 9.} Kiểm tra nhanh kiến thức nền tảng: Cho mẫu số liệu ghép nhóm có tổng cỡ mẫu n = 104. Nhóm chứa tứ phân vị thứ nhất Q$_1$ là nhóm đầu tiên có tần số tích lũy lớn hơn hoặc bằng:\par
\hspace*{1.5em}\textbf{A.} 52\quad\quad \textbf{B.} 104\quad\quad \textbf{C.} 26\quad\quad \textbf{D.} 78\par\vspace{2pt}
\textbf{\color{primary} Câu 10.} Cho dữ kiện nhận biết sau đây: Giá trị đại diện $c_i$ của nhóm số liệu nửa khoảng [46; 54) bằng. Phương án chính xác là:\par
\hspace*{1.5em}\textbf{A.} 8\quad\quad \textbf{B.} 48\quad\quad \textbf{C.} 50\quad\quad \textbf{D.} 52\par\vspace{2pt}
\textbf{\color{primary} Câu 11.} Hãy chỉ ra kết quả đúng trong bốn lựa chọn A, B, C, D đối với yêu cầu: Trong bảng phân bố tần số ghép nhóm, nhóm chứa mốt là nhóm có đặc điểm nào sau đây:\par
\begin{itemize}\setlength{\itemsep}{1pt}\setlength{\parskip}{0pt}
  \item[] \textbf{A.} Nhóm đầu tiên của bảng
  \item[] \textbf{B.} Nhóm có độ dài lớn nhất
  \item[] \textbf{C.} Nhóm có giá trị đại diện lớn nhất
  \item[] \textbf{D.} Nhóm có tần số lớn nhất
\end{itemize}\vspace{2pt}
\textbf{\color{primary} Câu 12.} Câu hỏi ôn tập lý thuyết và kỹ năng cơ bản (Chương III): Cho mẫu số liệu ghép nhóm có tổng cỡ mẫu n = 128. Nhóm chứa tứ phân vị thứ nhất Q$_1$ là nhóm đầu tiên có tần số tích lũy lớn hơn hoặc bằng:\par
\hspace*{1.5em}\textbf{A.} 96\quad\quad \textbf{B.} 64\quad\quad \textbf{C.} 32\quad\quad \textbf{D.} 128\par\vspace{2pt}
\textbf{\color{primary} Câu 13.} Theo nội dung sách giáo khoa Kết nối tri thức với cuộc sống (Chương III), giá trị đại diện $c_i$ của nhóm số liệu nửa khoảng [58; 66) bằng:\par
\hspace*{1.5em}\textbf{A.} 64\quad\quad \textbf{B.} 62\quad\quad \textbf{C.} 60\quad\quad \textbf{D.} 8\par\vspace{2pt}
\textbf{\color{primary} Câu 14.} Tìm đáp số đúng của bài toán cơ bản sau: Trong bảng phân bố tần số ghép nhóm, nhóm chứa mốt là nhóm có đặc điểm nào sau đây:\par
\begin{itemize}\setlength{\itemsep}{1pt}\setlength{\parskip}{0pt}
  \item[] \textbf{A.} Nhóm có giá trị đại diện lớn nhất
  \item[] \textbf{B.} Nhóm có tần số lớn nhất
  \item[] \textbf{C.} Nhóm đầu tiên của bảng
  \item[] \textbf{D.} Nhóm có độ dài lớn nhất
\end{itemize}\vspace{2pt}
\textbf{\color{primary} Câu 15.} Xét đối tượng toán học trong Chương III: Cho mẫu số liệu ghép nhóm có tổng cỡ mẫu n = 152. Nhóm chứa tứ phân vị thứ nhất Q$_1$ là nhóm đầu tiên có tần số tích lũy lớn hơn hoặc bằng. Giá trị hoặc tính chất tương ứng là:\par
\hspace*{1.5em}\textbf{A.} 38\quad\quad \textbf{B.} 152\quad\quad \textbf{C.} 114\quad\quad \textbf{D.} 76\par\vspace{2pt}
\textbf{\color{primary} Câu 16.} Hoàn thành phép xác định cơ bản sau đây: Giá trị đại diện $c_i$ của nhóm số liệu nửa khoảng [70; 78) bằng:\par
\hspace*{1.5em}\textbf{A.} 76\quad\quad \textbf{B.} 72\quad\quad \textbf{C.} 74\quad\quad \textbf{D.} 8\par\vspace{2pt}
\textbf{\color{primary} Câu 17.} Trong các khẳng định hoặc giá trị dưới đây, đâu là đáp án đúng cho câu hỏi: Trong bảng phân bố tần số ghép nhóm, nhóm chứa mốt là nhóm có đặc điểm nào sau đây?\par
\begin{itemize}\setlength{\itemsep}{1pt}\setlength{\parskip}{0pt}
  \item[] \textbf{A.} Nhóm có tần số lớn nhất
  \item[] \textbf{B.} Nhóm đầu tiên của bảng
  \item[] \textbf{C.} Nhóm có độ dài lớn nhất
  \item[] \textbf{D.} Nhóm có giá trị đại diện lớn nhất
\end{itemize}\vspace{2pt}
\textbf{\color{primary} Câu 18.} Bài tập khởi động Chương III: Cho mẫu số liệu ghép nhóm có tổng cỡ mẫu n = 176. Nhóm chứa tứ phân vị thứ nhất Q$_1$ là nhóm đầu tiên có tần số tích lũy lớn hơn hoặc bằng:\par
\hspace*{1.5em}\textbf{A.} 132\quad\quad \textbf{B.} 44\quad\quad \textbf{C.} 88\quad\quad \textbf{D.} 176\par\vspace{2pt}
\textbf{\color{primary} Câu 19.} Quan sát dữ kiện đề bài: Giá trị đại diện $c_i$ của nhóm số liệu nửa khoảng [82; 90) bằng. Hãy chọn phương án đúng:\par
\hspace*{1.5em}\textbf{A.} 88\quad\quad \textbf{B.} 86\quad\quad \textbf{C.} 8\quad\quad \textbf{D.} 84\par\vspace{2pt}
\textbf{\color{primary} Câu 20.} Sử dụng công thức hoặc tính chất đầu bài của Chương III để giải quyết yêu cầu: Trong bảng phân bố tần số ghép nhóm, nhóm chứa mốt là nhóm có đặc điểm nào sau đây:\par
\begin{itemize}\setlength{\itemsep}{1pt}\setlength{\parskip}{0pt}
  \item[] \textbf{A.} Nhóm đầu tiên của bảng
  \item[] \textbf{B.} Nhóm có tần số lớn nhất
  \item[] \textbf{C.} Nhóm có độ dài lớn nhất
  \item[] \textbf{D.} Nhóm có giá trị đại diện lớn nhất
\end{itemize}\vspace{2pt}
\subsection*{\color{accent} Phần II. Mức độ Trung bình (Thông hiểu -- Câu 21 đến Câu 40)}
\textbf{\color{primary} Câu 21.} Một mẫu số liệu ghép nhóm gồm 2 nhóm: nhóm [0; 10) có tần số 4 và nhóm [10; 20) có tần số 6. Số trung bình $\bar{x}$ (tính theo giá trị đại diện 5 và 15) bằng:\par
\hspace*{1.5em}\textbf{A.} 9,5\quad\quad \textbf{B.} 12\quad\quad \textbf{C.} 11\quad\quad \textbf{D.} 10\par\vspace{2pt}
\textbf{\color{primary} Câu 22.} Vận dụng quy tắc trong Chương III (Các số đặc trưng đo xu thế trung tâm của mẫu số liệu ghép nhóm) để giải bài toán: Khi tính trung vị $M_e$ của mẫu số liệu ghép nhóm có cỡ mẫu n, ta xác định nhóm chứa trung vị dựa trên đại lượng nào:\par
\begin{itemize}\setlength{\itemsep}{1pt}\setlength{\parskip}{0pt}
  \item[] \textbf{A.} Lấy trung bình cộng hai đầu mút của toàn mẫu
  \item[] \textbf{B.} So sánh n/2 với tần số tích lũy của các nhóm
  \item[] \textbf{C.} Chọn nhóm cuối cùng
  \item[] \textbf{D.} Chọn nhóm có tần số nhỏ nhất
\end{itemize}\vspace{2pt}
\textbf{\color{primary} Câu 23.} Biến đổi và tính toán theo kiến thức Chương III: Khảo sát điểm thi của 40 học sinh được ghép thành 4 nhóm: [4; 6) có 6 em, [6; 8) có 20 em, [8; 10) có 14 em. Nhóm chứa trung vị $M_e$ là. Kết quả thu được là:\par
\hspace*{1.5em}\textbf{A.} Không xác định được\quad\quad \textbf{B.} Nhóm [6; 8)\quad\quad \textbf{C.} Nhóm [8; 10)\quad\quad \textbf{D.} Nhóm [4; 6)\par\vspace{2pt}
\textbf{\color{primary} Câu 24.} Xét bài toán mức độ thông hiểu sau: Cho mẫu số liệu ghép nhóm gồm 2 nhóm: [40; 50) có tần số 4 và [50; 60) có tần số 6. Số trung bình cộng $\bar{x}$ của mẫu số liệu ghép nhóm bằng. Đáp án đúng là:\par
\hspace*{1.5em}\textbf{A.} 55\quad\quad \textbf{B.} 56\quad\quad \textbf{C.} 54\quad\quad \textbf{D.} 52\par\vspace{2pt}
\textbf{\color{primary} Câu 25.} Để hoàn thành phiếu học tập Chương III, hãy xác định: Một mẫu số liệu ghép nhóm gồm 2 nhóm: nhóm [0; 10) có tần số 4 và nhóm [10; 20) có tần số 6. Số trung bình $\bar{x}$ (tính theo giá trị đại diện 5 và 15) bằng:\par
\hspace*{1.5em}\textbf{A.} 11\quad\quad \textbf{B.} 10\quad\quad \textbf{C.} 12\quad\quad \textbf{D.} 9,5\par\vspace{2pt}
\textbf{\color{primary} Câu 26.} Phân tích dữ kiện bài toán: Khi tính trung vị $M_e$ của mẫu số liệu ghép nhóm có cỡ mẫu n, ta xác định nhóm chứa trung vị dựa trên đại lượng nào. Giá trị cần tìm bằng:\par
\begin{itemize}\setlength{\itemsep}{1pt}\setlength{\parskip}{0pt}
  \item[] \textbf{A.} Chọn nhóm có tần số nhỏ nhất
  \item[] \textbf{B.} Chọn nhóm cuối cùng
  \item[] \textbf{C.} So sánh n/2 với tần số tích lũy của các nhóm
  \item[] \textbf{D.} Lấy trung bình cộng hai đầu mút của toàn mẫu
\end{itemize}\vspace{2pt}
\textbf{\color{primary} Câu 27.} Áp dụng tính chất đã học ở Chương III vào bài tập sau: Khảo sát điểm thi của 40 học sinh được ghép thành 4 nhóm: [4; 6) có 6 em, [6; 8) có 20 em, [8; 10) có 14 em. Nhóm chứa trung vị $M_e$ là:\par
\hspace*{1.5em}\textbf{A.} Nhóm [6; 8)\quad\quad \textbf{B.} Nhóm [8; 10)\quad\quad \textbf{C.} Nhóm [4; 6)\quad\quad \textbf{D.} Không xác định được\par\vspace{2pt}
\textbf{\color{primary} Câu 28.} Thực hiện các bước suy luận hợp lý để giải quyết yêu cầu: Cho mẫu số liệu ghép nhóm gồm 2 nhóm: [80; 90) có tần số 4 và [90; 100) có tần số 6. Số trung bình cộng $\bar{x}$ của mẫu số liệu ghép nhóm bằng:\par
\hspace*{1.5em}\textbf{A.} 92\quad\quad \textbf{B.} 95\quad\quad \textbf{C.} 96\quad\quad \textbf{D.} 94\par\vspace{2pt}
\textbf{\color{primary} Câu 29.} Cho bài toán tự luyện (Các số đặc trưng đo xu thế trung tâm của mẫu số liệu ghép nhóm): Một mẫu số liệu ghép nhóm gồm 2 nhóm: nhóm [0; 10) có tần số 4 và nhóm [10; 20) có tần số 6. Số trung bình $\bar{x}$ (tính theo giá trị đại diện 5 và 15) bằng. Phương án nào sau đây chính xác?\par
\hspace*{1.5em}\textbf{A.} 10\quad\quad \textbf{B.} 12\quad\quad \textbf{C.} 9,5\quad\quad \textbf{D.} 11\par\vspace{2pt}
\textbf{\color{primary} Câu 30.} Kết hợp các phép biến đổi trong Chương III, hãy tìm kết quả của bài toán: Khi tính trung vị $M_e$ của mẫu số liệu ghép nhóm có cỡ mẫu n, ta xác định nhóm chứa trung vị dựa trên đại lượng nào:\par
\begin{itemize}\setlength{\itemsep}{1pt}\setlength{\parskip}{0pt}
  \item[] \textbf{A.} Lấy trung bình cộng hai đầu mút của toàn mẫu
  \item[] \textbf{B.} So sánh n/2 với tần số tích lũy của các nhóm
  \item[] \textbf{C.} Chọn nhóm cuối cùng
  \item[] \textbf{D.} Chọn nhóm có tần số nhỏ nhất
\end{itemize}\vspace{2pt}
\textbf{\color{primary} Câu 31.} Một bài kiểm tra định kỳ nội dung Chương III nêu câu hỏi: Khảo sát điểm thi của 40 học sinh được ghép thành 4 nhóm: [4; 6) có 6 em, [6; 8) có 20 em, [8; 10) có 14 em. Nhóm chứa trung vị $M_e$ là. Hãy chọn đáp số đúng:\par
\hspace*{1.5em}\textbf{A.} Nhóm [8; 10)\quad\quad \textbf{B.} Nhóm [4; 6)\quad\quad \textbf{C.} Không xác định được\quad\quad \textbf{D.} Nhóm [6; 8)\par\vspace{2pt}
\textbf{\color{primary} Câu 32.} Tìm lời giải đúng cho bài toán thông hiểu: Cho mẫu số liệu ghép nhóm gồm 2 nhóm: [120; 130) có tần số 4 và [130; 140) có tần số 6. Số trung bình cộng $\bar{x}$ của mẫu số liệu ghép nhóm bằng:\par
\hspace*{1.5em}\textbf{A.} 134\quad\quad \textbf{B.} 132\quad\quad \textbf{C.} 135\quad\quad \textbf{D.} 136\par\vspace{2pt}
\textbf{\color{primary} Câu 33.} Sử dụng mối liên hệ giữa các đại lượng trong Chương III: Một mẫu số liệu ghép nhóm gồm 2 nhóm: nhóm [0; 10) có tần số 4 và nhóm [10; 20) có tần số 6. Số trung bình $\bar{x}$ (tính theo giá trị đại diện 5 và 15) bằng:\par
\hspace*{1.5em}\textbf{A.} 11\quad\quad \textbf{B.} 12\quad\quad \textbf{C.} 9,5\quad\quad \textbf{D.} 10\par\vspace{2pt}
\textbf{\color{primary} Câu 34.} Hãy tính toán và chọn phương án đúng đối với bài tập: Khi tính trung vị $M_e$ của mẫu số liệu ghép nhóm có cỡ mẫu n, ta xác định nhóm chứa trung vị dựa trên đại lượng nào:\par
\begin{itemize}\setlength{\itemsep}{1pt}\setlength{\parskip}{0pt}
  \item[] \textbf{A.} Chọn nhóm có tần số nhỏ nhất
  \item[] \textbf{B.} So sánh n/2 với tần số tích lũy của các nhóm
  \item[] \textbf{C.} Lấy trung bình cộng hai đầu mút của toàn mẫu
  \item[] \textbf{D.} Chọn nhóm cuối cùng
\end{itemize}\vspace{2pt}
\textbf{\color{primary} Câu 35.} Dựa trên các tính chất trọng tâm của Chương III, xét bài toán: Khảo sát điểm thi của 40 học sinh được ghép thành 4 nhóm: [4; 6) có 6 em, [6; 8) có 20 em, [8; 10) có 14 em. Nhóm chứa trung vị $M_e$ là:\par
\hspace*{1.5em}\textbf{A.} Nhóm [4; 6)\quad\quad \textbf{B.} Nhóm [6; 8)\quad\quad \textbf{C.} Nhóm [8; 10)\quad\quad \textbf{D.} Không xác định được\par\vspace{2pt}
\textbf{\color{primary} Câu 36.} Thực hành giải toán (Các số đặc trưng đo xu thế trung tâm của mẫu số liệu ghép nhóm): Cho mẫu số liệu ghép nhóm gồm 2 nhóm: [160; 170) có tần số 4 và [170; 180) có tần số 6. Số trung bình cộng $\bar{x}$ của mẫu số liệu ghép nhóm bằng. Đáp án nào dưới đây thỏa mãn?\par
\hspace*{1.5em}\textbf{A.} 175\quad\quad \textbf{B.} 174\quad\quad \textbf{C.} 172\quad\quad \textbf{D.} 176\par\vspace{2pt}
\textbf{\color{primary} Câu 37.} Kiểm tra kỹ năng biến đổi và suy luận qua bài toán: Một mẫu số liệu ghép nhóm gồm 2 nhóm: nhóm [0; 10) có tần số 4 và nhóm [10; 20) có tần số 6. Số trung bình $\bar{x}$ (tính theo giá trị đại diện 5 và 15) bằng:\par
\hspace*{1.5em}\textbf{A.} 11\quad\quad \textbf{B.} 12\quad\quad \textbf{C.} 9,5\quad\quad \textbf{D.} 10\par\vspace{2pt}
\textbf{\color{primary} Câu 38.} Cho giả thiết thuộc phạm vi Chương III: Khi tính trung vị $M_e$ của mẫu số liệu ghép nhóm có cỡ mẫu n, ta xác định nhóm chứa trung vị dựa trên đại lượng nào. Kết luận đúng là:\par
\begin{itemize}\setlength{\itemsep}{1pt}\setlength{\parskip}{0pt}
  \item[] \textbf{A.} Lấy trung bình cộng hai đầu mút của toàn mẫu
  \item[] \textbf{B.} So sánh n/2 với tần số tích lũy của các nhóm
  \item[] \textbf{C.} Chọn nhóm cuối cùng
  \item[] \textbf{D.} Chọn nhóm có tần số nhỏ nhất
\end{itemize}\vspace{2pt}
\textbf{\color{primary} Câu 39.} Rút gọn hoặc tính giá trị theo yêu cầu sau của Chương III: Khảo sát điểm thi của 40 học sinh được ghép thành 4 nhóm: [4; 6) có 6 em, [6; 8) có 20 em, [8; 10) có 14 em. Nhóm chứa trung vị $M_e$ là:\par
\hspace*{1.5em}\textbf{A.} Nhóm [6; 8)\quad\quad \textbf{B.} Nhóm [8; 10)\quad\quad \textbf{C.} Không xác định được\quad\quad \textbf{D.} Nhóm [4; 6)\par\vspace{2pt}
\textbf{\color{primary} Câu 40.} Xác định phương án đúng sau khi thực hiện phép tính / suy luận: Cho mẫu số liệu ghép nhóm gồm 2 nhóm: [200; 210) có tần số 4 và [210; 220) có tần số 6. Số trung bình cộng $\bar{x}$ của mẫu số liệu ghép nhóm bằng:\par
\hspace*{1.5em}\textbf{A.} 215\quad\quad \textbf{B.} 214\quad\quad \textbf{C.} 212\quad\quad \textbf{D.} 216\par\vspace{2pt}
\subsection*{\color{accent} Phần III. Mức độ Khó (Vận dụng \& Vận dụng cao -- Câu 41 đến Câu 50)}
\textbf{\color{primary} Câu 41.} Cho mẫu số liệu ghép nhóm cỡ mẫu n = 40 với 3 nhóm: [10; 20) có 10 số liệu; [20; 30) có 20 số liệu; [30; 40) có 10 số liệu. Trung vị $M_e$ của mẫu số liệu ghép nhóm này bằng:\par
\hspace*{1.5em}\textbf{A.} 22,5\quad\quad \textbf{B.} 25\quad\quad \textbf{C.} 20\quad\quad \textbf{D.} 30\par\vspace{2pt}
\textbf{\color{primary} Câu 42.} Bài toán vận dụng (Chương III – Các số đặc trưng đo xu thế trung tâm của mẫu số liệu ghép nhóm): Cho mẫu số liệu ghép nhóm gồm 3 nhóm có độ dài bằng nhau: [10; 20) tần số 5; [20; 30) tần số 15; [30; 40) tần số 5. Mốt $M_o$ của mẫu số liệu ghép nhóm này bằng. Kết quả đúng là:\par
\hspace*{1.5em}\textbf{A.} 25\quad\quad \textbf{B.} 20\quad\quad \textbf{C.} 22,5\quad\quad \textbf{D.} 27,5\par\vspace{2pt}
\textbf{\color{primary} Câu 43.} Vận dụng kiến thức nâng cao trong phạm vi Chương III để giải bài toán: Cho mẫu số liệu ghép nhóm có 3 nhóm độ dài bằng nhau: [30; 40) tần số 5; [40; 50) tần số 15; [50; 60) tần số 5. Trung vị $M_e$ và Mốt $M_o$ của mẫu số liệu này cùng bằng:\par
\hspace*{1.5em}\textbf{A.} 40\quad\quad \textbf{B.} 45\quad\quad \textbf{C.} (42, 5)\quad\quad \textbf{D.} 50\par\vspace{2pt}
\textbf{\color{primary} Câu 44.} Phân tích và giải bài toán tổng hợp sau của Chương III: Cho mẫu số liệu ghép nhóm cỡ mẫu n = 40 với 3 nhóm: [10; 20) có 10 số liệu; [20; 30) có 20 số liệu; [30; 40) có 10 số liệu. Trung vị $M_e$ của mẫu số liệu ghép nhóm này bằng:\par
\hspace*{1.5em}\textbf{A.} 25\quad\quad \textbf{B.} 30\quad\quad \textbf{C.} 22,5\quad\quad \textbf{D.} 20\par\vspace{2pt}
\textbf{\color{primary} Câu 45.} Thử thách tư duy trong phạm vi kiến thức Chương III: Cho mẫu số liệu ghép nhóm gồm 3 nhóm có độ dài bằng nhau: [10; 20) tần số 5; [20; 30) tần số 15; [30; 40) tần số 5. Mốt $M_o$ của mẫu số liệu ghép nhóm này bằng. Đáp số chính xác là:\par
\hspace*{1.5em}\textbf{A.} 20\quad\quad \textbf{B.} 27,5\quad\quad \textbf{C.} 22,5\quad\quad \textbf{D.} 25\par\vspace{2pt}
\textbf{\color{primary} Câu 46.} Sử dụng phương pháp giải toán của Chương III để tìm đáp án cho bài toán: Cho mẫu số liệu ghép nhóm có 3 nhóm độ dài bằng nhau: [60; 70) tần số 5; [70; 80) tần số 15; [80; 90) tần số 5. Trung vị $M_e$ và Mốt $M_o$ của mẫu số liệu này cùng bằng:\par
\hspace*{1.5em}\textbf{A.} 75\quad\quad \textbf{B.} 80\quad\quad \textbf{C.} (72, 5)\quad\quad \textbf{D.} 70\par\vspace{2pt}
\textbf{\color{primary} Câu 47.} Xét bài toán vận dụng thực tiễn / nâng cao thuộc Chương III: Cho mẫu số liệu ghép nhóm cỡ mẫu n = 40 với 3 nhóm: [10; 20) có 10 số liệu; [20; 30) có 20 số liệu; [30; 40) có 10 số liệu. Trung vị $M_e$ của mẫu số liệu ghép nhóm này bằng:\par
\hspace*{1.5em}\textbf{A.} 25\quad\quad \textbf{B.} 20\quad\quad \textbf{C.} 30\quad\quad \textbf{D.} 22,5\par\vspace{2pt}
\textbf{\color{primary} Câu 48.} Tìm kết quả cuối cùng của bài toán phân loại học sinh khá giỏi (Chương III): Cho mẫu số liệu ghép nhóm gồm 3 nhóm có độ dài bằng nhau: [10; 20) tần số 5; [20; 30) tần số 15; [30; 40) tần số 5. Mốt $M_o$ của mẫu số liệu ghép nhóm này bằng:\par
\hspace*{1.5em}\textbf{A.} 27,5\quad\quad \textbf{B.} 25\quad\quad \textbf{C.} 20\quad\quad \textbf{D.} 22,5\par\vspace{2pt}
\textbf{\color{primary} Câu 49.} Kết hợp các tính chất đã học trong Chương III, hãy giải quyết bài toán: Cho mẫu số liệu ghép nhóm có 3 nhóm độ dài bằng nhau: [90; 100) tần số 5; [100; 110) tần số 15; [110; 120) tần số 5. Trung vị $M_e$ và Mốt $M_o$ của mẫu số liệu này cùng bằng:\par
\hspace*{1.5em}\textbf{A.} 110\quad\quad \textbf{B.} 105\quad\quad \textbf{C.} (102, 5)\quad\quad \textbf{D.} 100\par\vspace{2pt}
\textbf{\color{primary} Câu 50.} Câu hỏi vận dụng cao cuối Chương III: Cho mẫu số liệu ghép nhóm cỡ mẫu n = 40 với 3 nhóm: [10; 20) có 10 số liệu; [20; 30) có 20 số liệu; [30; 40) có 10 số liệu. Trung vị $M_e$ của mẫu số liệu ghép nhóm này bằng. Phương án đúng là:\par
\hspace*{1.5em}\textbf{A.} 22,5\quad\quad \textbf{B.} 20\quad\quad \textbf{C.} 30\quad\quad \textbf{D.} 25\par\vspace{2pt}
\subsection*{\color{primary} Bảng đáp án Chương III: Các số đặc trưng đo xu thế trung tâm của mẫu số liệu ghép nhóm}
\begin{center}
\renewcommand{\arraystretch}{1.25}
\begin{tabularx}{\textwidth}{*{10}{>{\centering\arraybackslash}X}}
\toprule
\textbf{\color{primary} Cột 1} & \textbf{\color{primary} Cột 2} & \textbf{\color{primary} Cột 3} & \textbf{\color{primary} Cột 4} & \textbf{\color{primary} Cột 5} & \textbf{\color{primary} Cột 6} & \textbf{\color{primary} Cột 7} & \textbf{\color{primary} Cột 8} & \textbf{\color{primary} Cột 9} & \textbf{\color{primary} Cột 10} \\
\midrule
\textbf{1.B} & \textbf{2.A} & \textbf{3.B} & \textbf{4.A} & \textbf{5.D} & \textbf{6.B} & \textbf{7.B} & \textbf{8.D} & \textbf{9.C} & \textbf{10.C} \\
\rowcolor{rowtint} \textbf{11.D} & \textbf{12.C} & \textbf{13.B} & \textbf{14.B} & \textbf{15.A} & \textbf{16.C} & \textbf{17.A} & \textbf{18.B} & \textbf{19.B} & \textbf{20.B} \\
\textbf{21.C} & \textbf{22.B} & \textbf{23.B} & \textbf{24.B} & \textbf{25.A} & \textbf{26.C} & \textbf{27.A} & \textbf{28.C} & \textbf{29.D} & \textbf{30.B} \\
\rowcolor{rowtint} \textbf{31.D} & \textbf{32.D} & \textbf{33.A} & \textbf{34.B} & \textbf{35.B} & \textbf{36.D} & \textbf{37.A} & \textbf{38.B} & \textbf{39.A} & \textbf{40.D} \\
\textbf{41.B} & \textbf{42.A} & \textbf{43.B} & \textbf{44.A} & \textbf{45.D} & \textbf{46.A} & \textbf{47.A} & \textbf{48.B} & \textbf{49.B} & \textbf{50.D} \\
\bottomrule
\end{tabularx}
\end{center}\vspace{10pt}
\section*{\color{primary} CHƯƠNG IV: QUAN HỆ SONG SONG TRONG KHÔNG GIAN}
\textit{Trọng tâm kiến thức SGK: Đường thẳng và mặt phẳng trong không gian, hai đường thẳng chéo nhau và hai đường thẳng song song, đường thẳng song song với mặt phẳng, hai mặt phẳng song song, phép chiếu song song.}\par\vspace{4pt}
\subsection*{\color{accent} Phần I. Mức độ Dễ (Nhận biết -- Câu 1 đến Câu 20)}
\textbf{\color{primary} Câu 1.} Qua ba điểm không thẳng hàng trong không gian, có bao nhiêu mặt phẳng đi qua cả ba điểm đó?\par
\begin{itemize}\setlength{\itemsep}{1pt}\setlength{\parskip}{0pt}
  \item[] \textbf{A.} Có vô số mặt phẳng
  \item[] \textbf{B.} Có một và chỉ một mặt phẳng
  \item[] \textbf{C.} Không có mặt phẳng nào
  \item[] \textbf{D.} Có hai mặt phẳng
\end{itemize}\vspace{2pt}
\textbf{\color{primary} Câu 2.} Hãy trả lời câu hỏi cơ bản sau: Hai đường thẳng trong không gian được gọi là chéo nhau khi nào?\par
\begin{itemize}\setlength{\itemsep}{1pt}\setlength{\parskip}{0pt}
  \item[] \textbf{A.} Khi chúng cùng nằm trong một mặt phẳng và không có điểm chung
  \item[] \textbf{B.} Khi chúng cắt nhau tại một điểm
  \item[] \textbf{C.} Khi chúng vuông góc với nhau
  \item[] \textbf{D.} Khi chúng không cùng nằm trong bất kỳ mặt phẳng nào (không đồng phẳng)
\end{itemize}\vspace{2pt}
\textbf{\color{primary} Câu 3.} Trong bài học thuộc Chương IV, xét nội dung: Cho hình chóp S.ABCD có đáy ABCD là hình bình hành. Giao tuyến của hai mặt phẳng (SAB) và (SCD) là. Đáp án đúng là:\par
\begin{itemize}\setlength{\itemsep}{1pt}\setlength{\parskip}{0pt}
  \item[] \textbf{A.} Đường thẳng AB
  \item[] \textbf{B.} Đường thẳng SO (O là tâm đáy)
  \item[] \textbf{C.} Đường thẳng đi qua S và song song với AB, CD
  \item[] \textbf{D.} Đường thẳng SC
\end{itemize}\vspace{2pt}
\textbf{\color{primary} Câu 4.} Dựa vào kiến thức chuẩn SGK của Chương IV (Quan hệ song song trong không gian), cho hình chóp có đáy là đa giác lồi 7 cạnh. Số mặt bên (là các hình tam giác) của hình chóp đó bằng:\par
\hspace*{1.5em}\textbf{A.} 6 mặt bên\quad\quad \textbf{B.} 7 mặt bên\quad\quad \textbf{C.} 8 mặt bên\quad\quad \textbf{D.} 14 mặt bên\par\vspace{2pt}
\textbf{\color{primary} Câu 5.} Chọn phương án trả lời đúng nhất: Qua ba điểm không thẳng hàng trong không gian, có bao nhiêu mặt phẳng đi qua cả ba điểm đó?\par
\begin{itemize}\setlength{\itemsep}{1pt}\setlength{\parskip}{0pt}
  \item[] \textbf{A.} Có hai mặt phẳng
  \item[] \textbf{B.} Có một và chỉ một mặt phẳng
  \item[] \textbf{C.} Có vô số mặt phẳng
  \item[] \textbf{D.} Không có mặt phẳng nào
\end{itemize}\vspace{2pt}
\textbf{\color{primary} Câu 6.} Khi làm bài tập cơ bản về quan hệ song song trong không gian, một học sinh gặp câu hỏi: Hai đường thẳng trong không gian được gọi là chéo nhau khi nào. Kết quả nào dưới đây đúng?\par
\begin{itemize}\setlength{\itemsep}{1pt}\setlength{\parskip}{0pt}
  \item[] \textbf{A.} Khi chúng cắt nhau tại một điểm
  \item[] \textbf{B.} Khi chúng vuông góc với nhau
  \item[] \textbf{C.} Khi chúng không cùng nằm trong bất kỳ mặt phẳng nào (không đồng phẳng)
  \item[] \textbf{D.} Khi chúng cùng nằm trong một mặt phẳng và không có điểm chung
\end{itemize}\vspace{2pt}
\textbf{\color{primary} Câu 7.} Xác định giá trị hoặc khẳng định đúng cho bài toán sau: Cho hình chóp S.ABCD có đáy ABCD là hình bình hành. Giao tuyến của hai mặt phẳng (SAB) và (SCD) là:\par
\begin{itemize}\setlength{\itemsep}{1pt}\setlength{\parskip}{0pt}
  \item[] \textbf{A.} Đường thẳng SO (O là tâm đáy)
  \item[] \textbf{B.} Đường thẳng đi qua S và song song với AB, CD
  \item[] \textbf{C.} Đường thẳng AB
  \item[] \textbf{D.} Đường thẳng SC
\end{itemize}\vspace{2pt}
\textbf{\color{primary} Câu 8.} Áp dụng trực tiếp định nghĩa và quy tắc trong Chương IV: Cho hình chóp có đáy là đa giác lồi 11 cạnh. Số mặt bên (là các hình tam giác) của hình chóp đó bằng:\par
\hspace*{1.5em}\textbf{A.} 11 mặt bên\quad\quad \textbf{B.} 22 mặt bên\quad\quad \textbf{C.} 12 mặt bên\quad\quad \textbf{D.} 10 mặt bên\par\vspace{2pt}
\textbf{\color{primary} Câu 9.} Kiểm tra nhanh kiến thức nền tảng: Qua ba điểm không thẳng hàng trong không gian, có bao nhiêu mặt phẳng đi qua cả ba điểm đó?\par
\begin{itemize}\setlength{\itemsep}{1pt}\setlength{\parskip}{0pt}
  \item[] \textbf{A.} Có một và chỉ một mặt phẳng
  \item[] \textbf{B.} Có hai mặt phẳng
  \item[] \textbf{C.} Không có mặt phẳng nào
  \item[] \textbf{D.} Có vô số mặt phẳng
\end{itemize}\vspace{2pt}
\textbf{\color{primary} Câu 10.} Cho dữ kiện nhận biết sau đây: Hai đường thẳng trong không gian được gọi là chéo nhau khi nào. Phương án chính xác là:\par
\begin{itemize}\setlength{\itemsep}{1pt}\setlength{\parskip}{0pt}
  \item[] \textbf{A.} Khi chúng vuông góc với nhau
  \item[] \textbf{B.} Khi chúng không cùng nằm trong bất kỳ mặt phẳng nào (không đồng phẳng)
  \item[] \textbf{C.} Khi chúng cùng nằm trong một mặt phẳng và không có điểm chung
  \item[] \textbf{D.} Khi chúng cắt nhau tại một điểm
\end{itemize}\vspace{2pt}
\textbf{\color{primary} Câu 11.} Hãy chỉ ra kết quả đúng trong bốn lựa chọn A, B, C, D đối với yêu cầu: Cho hình chóp S.ABCD có đáy ABCD là hình bình hành. Giao tuyến của hai mặt phẳng (SAB) và (SCD) là:\par
\begin{itemize}\setlength{\itemsep}{1pt}\setlength{\parskip}{0pt}
  \item[] \textbf{A.} Đường thẳng SC
  \item[] \textbf{B.} Đường thẳng AB
  \item[] \textbf{C.} Đường thẳng SO (O là tâm đáy)
  \item[] \textbf{D.} Đường thẳng đi qua S và song song với AB, CD
\end{itemize}\vspace{2pt}
\textbf{\color{primary} Câu 12.} Câu hỏi ôn tập lý thuyết và kỹ năng cơ bản (Chương IV): Cho hình chóp có đáy là đa giác lồi 15 cạnh. Số mặt bên (là các hình tam giác) của hình chóp đó bằng:\par
\hspace*{1.5em}\textbf{A.} 30 mặt bên\quad\quad \textbf{B.} 16 mặt bên\quad\quad \textbf{C.} 14 mặt bên\quad\quad \textbf{D.} 15 mặt bên\par\vspace{2pt}
\textbf{\color{primary} Câu 13.} Theo nội dung sách giáo khoa Kết nối tri thức với cuộc sống (Chương IV), qua ba điểm không thẳng hàng trong không gian, có bao nhiêu mặt phẳng đi qua cả ba điểm đó:\par
\begin{itemize}\setlength{\itemsep}{1pt}\setlength{\parskip}{0pt}
  \item[] \textbf{A.} Có vô số mặt phẳng
  \item[] \textbf{B.} Có hai mặt phẳng
  \item[] \textbf{C.} Không có mặt phẳng nào
  \item[] \textbf{D.} Có một và chỉ một mặt phẳng
\end{itemize}\vspace{2pt}
\textbf{\color{primary} Câu 14.} Tìm đáp số đúng của bài toán cơ bản sau: Hai đường thẳng trong không gian được gọi là chéo nhau khi nào:\par
\begin{itemize}\setlength{\itemsep}{1pt}\setlength{\parskip}{0pt}
  \item[] \textbf{A.} Khi chúng không cùng nằm trong bất kỳ mặt phẳng nào (không đồng phẳng)
  \item[] \textbf{B.} Khi chúng cắt nhau tại một điểm
  \item[] \textbf{C.} Khi chúng vuông góc với nhau
  \item[] \textbf{D.} Khi chúng cùng nằm trong một mặt phẳng và không có điểm chung
\end{itemize}\vspace{2pt}
\textbf{\color{primary} Câu 15.} Xét đối tượng toán học trong Chương IV: Cho hình chóp S.ABCD có đáy ABCD là hình bình hành. Giao tuyến của hai mặt phẳng (SAB) và (SCD) là. Giá trị hoặc tính chất tương ứng là:\par
\begin{itemize}\setlength{\itemsep}{1pt}\setlength{\parskip}{0pt}
  \item[] \textbf{A.} Đường thẳng đi qua S và song song với AB, CD
  \item[] \textbf{B.} Đường thẳng AB
  \item[] \textbf{C.} Đường thẳng SO (O là tâm đáy)
  \item[] \textbf{D.} Đường thẳng SC
\end{itemize}\vspace{2pt}
\textbf{\color{primary} Câu 16.} Hoàn thành phép xác định cơ bản sau đây: Cho hình chóp có đáy là đa giác lồi 19 cạnh. Số mặt bên (là các hình tam giác) của hình chóp đó bằng:\par
\hspace*{1.5em}\textbf{A.} 20 mặt bên\quad\quad \textbf{B.} 18 mặt bên\quad\quad \textbf{C.} 19 mặt bên\quad\quad \textbf{D.} 38 mặt bên\par\vspace{2pt}
\textbf{\color{primary} Câu 17.} Trong các khẳng định hoặc giá trị dưới đây, đâu là đáp án đúng cho câu hỏi: Qua ba điểm không thẳng hàng trong không gian, có bao nhiêu mặt phẳng đi qua cả ba điểm đó?\par
\begin{itemize}\setlength{\itemsep}{1pt}\setlength{\parskip}{0pt}
  \item[] \textbf{A.} Có một và chỉ một mặt phẳng
  \item[] \textbf{B.} Có vô số mặt phẳng
  \item[] \textbf{C.} Không có mặt phẳng nào
  \item[] \textbf{D.} Có hai mặt phẳng
\end{itemize}\vspace{2pt}
\textbf{\color{primary} Câu 18.} Bài tập khởi động Chương IV: Hai đường thẳng trong không gian được gọi là chéo nhau khi nào:\par
\begin{itemize}\setlength{\itemsep}{1pt}\setlength{\parskip}{0pt}
  \item[] \textbf{A.} Khi chúng cắt nhau tại một điểm
  \item[] \textbf{B.} Khi chúng không cùng nằm trong bất kỳ mặt phẳng nào (không đồng phẳng)
  \item[] \textbf{C.} Khi chúng cùng nằm trong một mặt phẳng và không có điểm chung
  \item[] \textbf{D.} Khi chúng vuông góc với nhau
\end{itemize}\vspace{2pt}
\textbf{\color{primary} Câu 19.} Quan sát dữ kiện đề bài: Cho hình chóp S.ABCD có đáy ABCD là hình bình hành. Giao tuyến của hai mặt phẳng (SAB) và (SCD) là. Hãy chọn phương án đúng:\par
\begin{itemize}\setlength{\itemsep}{1pt}\setlength{\parskip}{0pt}
  \item[] \textbf{A.} Đường thẳng đi qua S và song song với AB, CD
  \item[] \textbf{B.} Đường thẳng SC
  \item[] \textbf{C.} Đường thẳng SO (O là tâm đáy)
  \item[] \textbf{D.} Đường thẳng AB
\end{itemize}\vspace{2pt}
\textbf{\color{primary} Câu 20.} Sử dụng công thức hoặc tính chất đầu bài của Chương IV để giải quyết yêu cầu: Cho hình chóp có đáy là đa giác lồi 23 cạnh. Số mặt bên (là các hình tam giác) của hình chóp đó bằng:\par
\hspace*{1.5em}\textbf{A.} 24 mặt bên\quad\quad \textbf{B.} 22 mặt bên\quad\quad \textbf{C.} 46 mặt bên\quad\quad \textbf{D.} 23 mặt bên\par\vspace{2pt}
\subsection*{\color{accent} Phần II. Mức độ Trung bình (Thông hiểu -- Câu 21 đến Câu 40)}
\textbf{\color{primary} Câu 21.} Cho tứ diện ABCD. Gọi M, N lần lượt là trung điểm của AB và AC. Khi đó đường thẳng MN song song với mặt phẳng nào dưới đây?\par
\hspace*{1.5em}\textbf{A.} Mặt phẳng (ABD)\quad\quad \textbf{B.} Mặt phẳng (ABC)\quad\quad \textbf{C.} Mặt phẳng (ACD)\quad\quad \textbf{D.} Mặt phẳng (BCD)\par\vspace{2pt}
\textbf{\color{primary} Câu 22.} Vận dụng quy tắc trong Chương IV (Quan hệ song song trong không gian) để giải bài toán: Điều kiện đủ để một mặt phẳng (P) song song với mặt phẳng (Q) là:\par
\begin{itemize}\setlength{\itemsep}{1pt}\setlength{\parskip}{0pt}
  \item[] \textbf{A.} Mặt phẳng (P) chứa một đường thẳng song song với (Q)
  \item[] \textbf{B.} Mặt phẳng (P) chứa hai đường thẳng song song với nhau và song song với (Q)
  \item[] \textbf{C.} Mặt phẳng (P) chứa hai đường thẳng cắt nhau và cùng song song với mặt phẳng (Q)
  \item[] \textbf{D.} Mặt phẳng (P) và (Q) có chung một giao tuyến
\end{itemize}\vspace{2pt}
\textbf{\color{primary} Câu 23.} Biến đổi và tính toán theo kiến thức Chương IV: Tính chất nào sau đây được bảo toàn qua phép chiếu song song (khi đường thẳng không song song với phương chiếu). Kết quả thu được là:\par
\begin{itemize}\setlength{\itemsep}{1pt}\setlength{\parskip}{0pt}
  \item[] \textbf{A.} Độ dài thực của mọi đoạn thẳng
  \item[] \textbf{B.} Diện tích của mọi tam giác
  \item[] \textbf{C.} Số đo của mọi góc
  \item[] \textbf{D.} Tỉ số độ dài của hai đoạn thẳng cùng nằm trên một đường thẳng hoặc trên hai đường thẳng song song
\end{itemize}\vspace{2pt}
\textbf{\color{primary} Câu 24.} Xét bài toán mức độ thông hiểu sau: Cho tứ diện ABCD có cạnh BC = 14 cm. Gọi M, N lần lượt là trung điểm của AB và AC. Đường thẳng MN song song với mặt phẳng (BCD) và có độ dài đoạn MN bằng. Đáp án đúng là:\par
\hspace*{1.5em}\textbf{A.} 7 cm\quad\quad \textbf{B.} 14 cm\quad\quad \textbf{C.} 28 cm\quad\quad \textbf{D.} 5 cm\par\vspace{2pt}
\textbf{\color{primary} Câu 25.} Để hoàn thành phiếu học tập Chương IV, hãy xác định: Cho tứ diện ABCD. Gọi M, N lần lượt là trung điểm của AB và AC. Khi đó đường thẳng MN song song với mặt phẳng nào dưới đây:\par
\hspace*{1.5em}\textbf{A.} Mặt phẳng (ABD)\quad\quad \textbf{B.} Mặt phẳng (BCD)\quad\quad \textbf{C.} Mặt phẳng (ABC)\quad\quad \textbf{D.} Mặt phẳng (ACD)\par\vspace{2pt}
\textbf{\color{primary} Câu 26.} Phân tích dữ kiện bài toán: Điều kiện đủ để một mặt phẳng (P) song song với mặt phẳng (Q) là. Giá trị cần tìm bằng:\par
\begin{itemize}\setlength{\itemsep}{1pt}\setlength{\parskip}{0pt}
  \item[] \textbf{A.} Mặt phẳng (P) chứa hai đường thẳng cắt nhau và cùng song song với mặt phẳng (Q)
  \item[] \textbf{B.} Mặt phẳng (P) chứa một đường thẳng song song với (Q)
  \item[] \textbf{C.} Mặt phẳng (P) và (Q) có chung một giao tuyến
  \item[] \textbf{D.} Mặt phẳng (P) chứa hai đường thẳng song song với nhau và song song với (Q)
\end{itemize}\vspace{2pt}
\textbf{\color{primary} Câu 27.} Áp dụng tính chất đã học ở Chương IV vào bài tập sau: Tính chất nào sau đây được bảo toàn qua phép chiếu song song (khi đường thẳng không song song với phương chiếu):\par
\begin{itemize}\setlength{\itemsep}{1pt}\setlength{\parskip}{0pt}
  \item[] \textbf{A.} Độ dài thực của mọi đoạn thẳng
  \item[] \textbf{B.} Diện tích của mọi tam giác
  \item[] \textbf{C.} Số đo của mọi góc
  \item[] \textbf{D.} Tỉ số độ dài của hai đoạn thẳng cùng nằm trên một đường thẳng hoặc trên hai đường thẳng song song
\end{itemize}\vspace{2pt}
\textbf{\color{primary} Câu 28.} Thực hiện các bước suy luận hợp lý để giải quyết yêu cầu: Cho tứ diện ABCD có cạnh BC = 22 cm. Gọi M, N lần lượt là trung điểm của AB và AC. Đường thẳng MN song song với mặt phẳng (BCD) và có độ dài đoạn MN bằng:\par
\hspace*{1.5em}\textbf{A.} 9 cm\quad\quad \textbf{B.} 22 cm\quad\quad \textbf{C.} 11 cm\quad\quad \textbf{D.} 44 cm\par\vspace{2pt}
\textbf{\color{primary} Câu 29.} Cho bài toán tự luyện (Quan hệ song song trong không gian): Cho tứ diện ABCD. Gọi M, N lần lượt là trung điểm của AB và AC. Khi đó đường thẳng MN song song với mặt phẳng nào dưới đây. Phương án nào sau đây chính xác?\par
\hspace*{1.5em}\textbf{A.} Mặt phẳng (ACD)\quad\quad \textbf{B.} Mặt phẳng (BCD)\quad\quad \textbf{C.} Mặt phẳng (ABC)\quad\quad \textbf{D.} Mặt phẳng (ABD)\par\vspace{2pt}
\textbf{\color{primary} Câu 30.} Kết hợp các phép biến đổi trong Chương IV, hãy tìm kết quả của bài toán: Điều kiện đủ để một mặt phẳng (P) song song với mặt phẳng (Q) là:\par
\begin{itemize}\setlength{\itemsep}{1pt}\setlength{\parskip}{0pt}
  \item[] \textbf{A.} Mặt phẳng (P) chứa hai đường thẳng song song với nhau và song song với (Q)
  \item[] \textbf{B.} Mặt phẳng (P) chứa một đường thẳng song song với (Q)
  \item[] \textbf{C.} Mặt phẳng (P) chứa hai đường thẳng cắt nhau và cùng song song với mặt phẳng (Q)
  \item[] \textbf{D.} Mặt phẳng (P) và (Q) có chung một giao tuyến
\end{itemize}\vspace{2pt}
\textbf{\color{primary} Câu 31.} Một bài kiểm tra định kỳ nội dung Chương IV nêu câu hỏi: Tính chất nào sau đây được bảo toàn qua phép chiếu song song (khi đường thẳng không song song với phương chiếu). Hãy chọn đáp số đúng:\par
\begin{itemize}\setlength{\itemsep}{1pt}\setlength{\parskip}{0pt}
  \item[] \textbf{A.} Độ dài thực của mọi đoạn thẳng
  \item[] \textbf{B.} Diện tích của mọi tam giác
  \item[] \textbf{C.} Số đo của mọi góc
  \item[] \textbf{D.} Tỉ số độ dài của hai đoạn thẳng cùng nằm trên một đường thẳng hoặc trên hai đường thẳng song song
\end{itemize}\vspace{2pt}
\textbf{\color{primary} Câu 32.} Tìm lời giải đúng cho bài toán thông hiểu: Cho tứ diện ABCD có cạnh BC = 30 cm. Gọi M, N lần lượt là trung điểm của AB và AC. Đường thẳng MN song song với mặt phẳng (BCD) và có độ dài đoạn MN bằng:\par
\hspace*{1.5em}\textbf{A.} 30 cm\quad\quad \textbf{B.} 13 cm\quad\quad \textbf{C.} 60 cm\quad\quad \textbf{D.} 15 cm\par\vspace{2pt}
\textbf{\color{primary} Câu 33.} Sử dụng mối liên hệ giữa các đại lượng trong Chương IV: Cho tứ diện ABCD. Gọi M, N lần lượt là trung điểm của AB và AC. Khi đó đường thẳng MN song song với mặt phẳng nào dưới đây:\par
\hspace*{1.5em}\textbf{A.} Mặt phẳng (ABD)\quad\quad \textbf{B.} Mặt phẳng (ACD)\quad\quad \textbf{C.} Mặt phẳng (BCD)\quad\quad \textbf{D.} Mặt phẳng (ABC)\par\vspace{2pt}
\textbf{\color{primary} Câu 34.} Hãy tính toán và chọn phương án đúng đối với bài tập: Điều kiện đủ để một mặt phẳng (P) song song với mặt phẳng (Q) là:\par
\begin{itemize}\setlength{\itemsep}{1pt}\setlength{\parskip}{0pt}
  \item[] \textbf{A.} Mặt phẳng (P) chứa hai đường thẳng song song với nhau và song song với (Q)
  \item[] \textbf{B.} Mặt phẳng (P) chứa hai đường thẳng cắt nhau và cùng song song với mặt phẳng (Q)
  \item[] \textbf{C.} Mặt phẳng (P) và (Q) có chung một giao tuyến
  \item[] \textbf{D.} Mặt phẳng (P) chứa một đường thẳng song song với (Q)
\end{itemize}\vspace{2pt}
\textbf{\color{primary} Câu 35.} Dựa trên các tính chất trọng tâm của Chương IV, xét bài toán: Tính chất nào sau đây được bảo toàn qua phép chiếu song song (khi đường thẳng không song song với phương chiếu):\par
\begin{itemize}\setlength{\itemsep}{1pt}\setlength{\parskip}{0pt}
  \item[] \textbf{A.} Số đo của mọi góc
  \item[] \textbf{B.} Tỉ số độ dài của hai đoạn thẳng cùng nằm trên một đường thẳng hoặc trên hai đường thẳng song song
  \item[] \textbf{C.} Độ dài thực của mọi đoạn thẳng
  \item[] \textbf{D.} Diện tích của mọi tam giác
\end{itemize}\vspace{2pt}
\textbf{\color{primary} Câu 36.} Thực hành giải toán (Quan hệ song song trong không gian): Cho tứ diện ABCD có cạnh BC = 38 cm. Gọi M, N lần lượt là trung điểm của AB và AC. Đường thẳng MN song song với mặt phẳng (BCD) và có độ dài đoạn MN bằng. Đáp án nào dưới đây thỏa mãn?\par
\hspace*{1.5em}\textbf{A.} 38 cm\quad\quad \textbf{B.} 17 cm\quad\quad \textbf{C.} 76 cm\quad\quad \textbf{D.} 19 cm\par\vspace{2pt}
\textbf{\color{primary} Câu 37.} Kiểm tra kỹ năng biến đổi và suy luận qua bài toán: Cho tứ diện ABCD. Gọi M, N lần lượt là trung điểm của AB và AC. Khi đó đường thẳng MN song song với mặt phẳng nào dưới đây:\par
\hspace*{1.5em}\textbf{A.} Mặt phẳng (ABC)\quad\quad \textbf{B.} Mặt phẳng (BCD)\quad\quad \textbf{C.} Mặt phẳng (ABD)\quad\quad \textbf{D.} Mặt phẳng (ACD)\par\vspace{2pt}
\textbf{\color{primary} Câu 38.} Cho giả thiết thuộc phạm vi Chương IV: Điều kiện đủ để một mặt phẳng (P) song song với mặt phẳng (Q) là. Kết luận đúng là:\par
\begin{itemize}\setlength{\itemsep}{1pt}\setlength{\parskip}{0pt}
  \item[] \textbf{A.} Mặt phẳng (P) chứa hai đường thẳng song song với nhau và song song với (Q)
  \item[] \textbf{B.} Mặt phẳng (P) chứa một đường thẳng song song với (Q)
  \item[] \textbf{C.} Mặt phẳng (P) chứa hai đường thẳng cắt nhau và cùng song song với mặt phẳng (Q)
  \item[] \textbf{D.} Mặt phẳng (P) và (Q) có chung một giao tuyến
\end{itemize}\vspace{2pt}
\textbf{\color{primary} Câu 39.} Rút gọn hoặc tính giá trị theo yêu cầu sau của Chương IV: Tính chất nào sau đây được bảo toàn qua phép chiếu song song (khi đường thẳng không song song với phương chiếu):\par
\begin{itemize}\setlength{\itemsep}{1pt}\setlength{\parskip}{0pt}
  \item[] \textbf{A.} Số đo của mọi góc
  \item[] \textbf{B.} Tỉ số độ dài của hai đoạn thẳng cùng nằm trên một đường thẳng hoặc trên hai đường thẳng song song
  \item[] \textbf{C.} Diện tích của mọi tam giác
  \item[] \textbf{D.} Độ dài thực của mọi đoạn thẳng
\end{itemize}\vspace{2pt}
\textbf{\color{primary} Câu 40.} Xác định phương án đúng sau khi thực hiện phép tính / suy luận: Cho tứ diện ABCD có cạnh BC = 46 cm. Gọi M, N lần lượt là trung điểm của AB và AC. Đường thẳng MN song song với mặt phẳng (BCD) và có độ dài đoạn MN bằng:\par
\hspace*{1.5em}\textbf{A.} 92 cm\quad\quad \textbf{B.} 46 cm\quad\quad \textbf{C.} 23 cm\quad\quad \textbf{D.} 21 cm\par\vspace{2pt}
\subsection*{\color{accent} Phần III. Mức độ Khó (Vận dụng \& Vận dụng cao -- Câu 41 đến Câu 50)}
\textbf{\color{primary} Câu 41.} Cho tứ diện đều ABCD có độ dài tất cả các cạnh bằng 6 cm. Gọi M, N, P, Q lần lượt là trung điểm của các cạnh AB, BC, CD, DA. Thiết diện MNPQ là hình bình hành có chu vi bằng:\par
\hspace*{1.5em}\textbf{A.} 6 cm\quad\quad \textbf{B.} 12 cm\quad\quad \textbf{C.} 24 cm\quad\quad \textbf{D.} 18 cm\par\vspace{2pt}
\textbf{\color{primary} Câu 42.} Bài toán vận dụng (Chương IV – Quan hệ song song trong không gian): Cho hình hộp ABCD.A'B'C'D'. Mặt phẳng (AB'D') song song với mặt phẳng nào sau đây. Kết quả đúng là:\par
\hspace*{1.5em}\textbf{A.} Mặt phẳng (BC'D)\quad\quad \textbf{B.} Mặt phẳng (BDD'B')\quad\quad \textbf{C.} Mặt phẳng (A'B'C'D')\quad\quad \textbf{D.} Mặt phẳng (ACC'A')\par\vspace{2pt}
\textbf{\color{primary} Câu 43.} Vận dụng kiến thức nâng cao trong phạm vi Chương IV để giải bài toán: Cho tứ diện đều ABCD có độ dài tất cả các cạnh bằng 10 cm. Gọi M, N, P, Q lần lượt là trung điểm của các cạnh AB, BC, CD, DA. Thiết diện MNPQ là hình bình hành có chu vi bằng:\par
\hspace*{1.5em}\textbf{A.} 20 cm\quad\quad \textbf{B.} 40 cm\quad\quad \textbf{C.} 30 cm\quad\quad \textbf{D.} 10 cm\par\vspace{2pt}
\textbf{\color{primary} Câu 44.} Phân tích và giải bài toán tổng hợp sau của Chương IV: Cho hình hộp ABCD.A'B'C'D'. Mặt phẳng (AB'D') song song với mặt phẳng nào sau đây:\par
\hspace*{1.5em}\textbf{A.} Mặt phẳng (BC'D)\quad\quad \textbf{B.} Mặt phẳng (ACC'A')\quad\quad \textbf{C.} Mặt phẳng (A'B'C'D')\quad\quad \textbf{D.} Mặt phẳng (BDD'B')\par\vspace{2pt}
\textbf{\color{primary} Câu 45.} Thử thách tư duy trong phạm vi kiến thức Chương IV: Cho tứ diện đều ABCD có độ dài tất cả các cạnh bằng 14 cm. Gọi M, N, P, Q lần lượt là trung điểm của các cạnh AB, BC, CD, DA. Thiết diện MNPQ là hình bình hành có chu vi bằng. Đáp số chính xác là:\par
\hspace*{1.5em}\textbf{A.} 28 cm\quad\quad \textbf{B.} 14 cm\quad\quad \textbf{C.} 42 cm\quad\quad \textbf{D.} 56 cm\par\vspace{2pt}
\textbf{\color{primary} Câu 46.} Sử dụng phương pháp giải toán của Chương IV để tìm đáp án cho bài toán: Cho hình hộp ABCD.A'B'C'D'. Mặt phẳng (AB'D') song song với mặt phẳng nào sau đây:\par
\hspace*{1.5em}\textbf{A.} Mặt phẳng (BDD'B')\quad\quad \textbf{B.} Mặt phẳng (ACC'A')\quad\quad \textbf{C.} Mặt phẳng (BC'D)\quad\quad \textbf{D.} Mặt phẳng (A'B'C'D')\par\vspace{2pt}
\textbf{\color{primary} Câu 47.} Xét bài toán vận dụng thực tiễn / nâng cao thuộc Chương IV: Cho tứ diện đều ABCD có độ dài tất cả các cạnh bằng 18 cm. Gọi M, N, P, Q lần lượt là trung điểm của các cạnh AB, BC, CD, DA. Thiết diện MNPQ là hình bình hành có chu vi bằng:\par
\hspace*{1.5em}\textbf{A.} 54 cm\quad\quad \textbf{B.} 72 cm\quad\quad \textbf{C.} 18 cm\quad\quad \textbf{D.} 36 cm\par\vspace{2pt}
\textbf{\color{primary} Câu 48.} Tìm kết quả cuối cùng của bài toán phân loại học sinh khá giỏi (Chương IV): Cho hình hộp ABCD.A'B'C'D'. Mặt phẳng (AB'D') song song với mặt phẳng nào sau đây:\par
\hspace*{1.5em}\textbf{A.} Mặt phẳng (BDD'B')\quad\quad \textbf{B.} Mặt phẳng (ACC'A')\quad\quad \textbf{C.} Mặt phẳng (A'B'C'D')\quad\quad \textbf{D.} Mặt phẳng (BC'D)\par\vspace{2pt}
\textbf{\color{primary} Câu 49.} Kết hợp các tính chất đã học trong Chương IV, hãy giải quyết bài toán: Cho tứ diện đều ABCD có độ dài tất cả các cạnh bằng 22 cm. Gọi M, N, P, Q lần lượt là trung điểm của các cạnh AB, BC, CD, DA. Thiết diện MNPQ là hình bình hành có chu vi bằng:\par
\hspace*{1.5em}\textbf{A.} 88 cm\quad\quad \textbf{B.} 66 cm\quad\quad \textbf{C.} 22 cm\quad\quad \textbf{D.} 44 cm\par\vspace{2pt}
\textbf{\color{primary} Câu 50.} Câu hỏi vận dụng cao cuối Chương IV: Cho hình hộp ABCD.A'B'C'D'. Mặt phẳng (AB'D') song song với mặt phẳng nào sau đây. Phương án đúng là:\par
\hspace*{1.5em}\textbf{A.} Mặt phẳng (BC'D)\quad\quad \textbf{B.} Mặt phẳng (BDD'B')\quad\quad \textbf{C.} Mặt phẳng (A'B'C'D')\quad\quad \textbf{D.} Mặt phẳng (ACC'A')\par\vspace{2pt}
\subsection*{\color{primary} Bảng đáp án Chương IV: Quan hệ song song trong không gian}
\begin{center}
\renewcommand{\arraystretch}{1.25}
\begin{tabularx}{\textwidth}{*{10}{>{\centering\arraybackslash}X}}
\toprule
\textbf{\color{primary} Cột 1} & \textbf{\color{primary} Cột 2} & \textbf{\color{primary} Cột 3} & \textbf{\color{primary} Cột 4} & \textbf{\color{primary} Cột 5} & \textbf{\color{primary} Cột 6} & \textbf{\color{primary} Cột 7} & \textbf{\color{primary} Cột 8} & \textbf{\color{primary} Cột 9} & \textbf{\color{primary} Cột 10} \\
\midrule
\textbf{1.B} & \textbf{2.D} & \textbf{3.C} & \textbf{4.B} & \textbf{5.B} & \textbf{6.C} & \textbf{7.B} & \textbf{8.A} & \textbf{9.A} & \textbf{10.B} \\
\rowcolor{rowtint} \textbf{11.D} & \textbf{12.D} & \textbf{13.D} & \textbf{14.A} & \textbf{15.A} & \textbf{16.C} & \textbf{17.A} & \textbf{18.B} & \textbf{19.A} & \textbf{20.D} \\
\textbf{21.D} & \textbf{22.C} & \textbf{23.D} & \textbf{24.A} & \textbf{25.B} & \textbf{26.A} & \textbf{27.D} & \textbf{28.C} & \textbf{29.B} & \textbf{30.C} \\
\rowcolor{rowtint} \textbf{31.D} & \textbf{32.D} & \textbf{33.C} & \textbf{34.B} & \textbf{35.B} & \textbf{36.D} & \textbf{37.B} & \textbf{38.C} & \textbf{39.B} & \textbf{40.C} \\
\textbf{41.B} & \textbf{42.A} & \textbf{43.A} & \textbf{44.A} & \textbf{45.A} & \textbf{46.C} & \textbf{47.D} & \textbf{48.D} & \textbf{49.D} & \textbf{50.A} \\
\bottomrule
\end{tabularx}
\end{center}\vspace{10pt}
\section*{\color{primary} CHƯƠNG V: GIỚI HẠN. HÀM SỐ LIÊN TỤC}
\textit{Trọng tâm kiến thức SGK: Giới hạn của dãy số, tổng của cấp số nhân lùi vô hạn, giới hạn của hàm số tại một điểm và tại vô cực, hàm số liên tục tại một điểm và trên một khoảng/đoạn.}\par\vspace{4pt}
\subsection*{\color{accent} Phần I. Mức độ Dễ (Nhận biết -- Câu 1 đến Câu 20)}
\textbf{\color{primary} Câu 1.} Giá trị của giới hạn dãy số lim ((2n + 3) / (n + 1)) khi n $\to$ +$\infty$ bằng:\par
\hspace*{1.5em}\textbf{A.} 2\quad\quad \textbf{B.} 3\quad\quad \textbf{C.} 1\quad\quad \textbf{D.} 0\par\vspace{2pt}
\textbf{\color{primary} Câu 2.} Thực hiện yêu cầu nhận biết sau: Giá trị của giới hạn hàm số lim(x $\to$ 2) (2x$^2$ + 3x - 1) bằng.\par
\hspace*{1.5em}\textbf{A.} 3\quad\quad \textbf{B.} 0\quad\quad \textbf{C.} 9\quad\quad \textbf{D.} 13\par\vspace{2pt}
\textbf{\color{primary} Câu 3.} Trong bài học thuộc Chương V, xét nội dung: Tổng của cấp số nhân lùi vô hạn có số hạng đầu u$_1$ = 1 và công bội q = 1/2 bằng. Đáp án đúng là:\par
\hspace*{1.5em}\textbf{A.} S = 1\quad\quad \textbf{B.} S = 4\quad\quad \textbf{C.} S = 2\quad\quad \textbf{D.} S = 3/2\par\vspace{2pt}
\textbf{\color{primary} Câu 4.} Dựa vào kiến thức chuẩn SGK của Chương V (Giới hạn. Hàm số liên tục), giá trị của giới hạn dãy số lim ((5n + 3) / (n + 1)) khi n $\to$ +$\infty$ bằng:\par
\hspace*{1.5em}\textbf{A.} 5\quad\quad \textbf{B.} 1\quad\quad \textbf{C.} 3\quad\quad \textbf{D.} 0\par\vspace{2pt}
\textbf{\color{primary} Câu 5.} Chọn kết quả đúng nhất khi xét: Giá trị của giới hạn hàm số lim(x $\to$ 2) (5x$^2$ + 3x - 1) bằng:\par
\hspace*{1.5em}\textbf{A.} 25\quad\quad \textbf{B.} 25 (1)\quad\quad \textbf{C.} 0\quad\quad \textbf{D.} 15\par\vspace{2pt}
\textbf{\color{primary} Câu 6.} Khi làm bài tập cơ bản về giới hạn. hàm số liên tục, một học sinh gặp câu hỏi: Tổng của cấp số nhân lùi vô hạn có số hạng đầu u$_1$ = 1 và công bội q = 1/2 bằng. Kết quả nào dưới đây đúng?\par
\hspace*{1.5em}\textbf{A.} S = 3/2\quad\quad \textbf{B.} S = 2\quad\quad \textbf{C.} S = 1\quad\quad \textbf{D.} S = 4\par\vspace{2pt}
\textbf{\color{primary} Câu 7.} Xác định giá trị hoặc khẳng định đúng cho bài toán sau: Giá trị của giới hạn dãy số lim ((8n + 3) / (n + 1)) khi n $\to$ +$\infty$ bằng:\par
\hspace*{1.5em}\textbf{A.} 0\quad\quad \textbf{B.} 3\quad\quad \textbf{C.} 8\quad\quad \textbf{D.} 1\par\vspace{2pt}
\textbf{\color{primary} Câu 8.} Áp dụng trực tiếp định nghĩa và quy tắc trong Chương V: Giá trị của giới hạn hàm số lim(x $\to$ 2) (8x$^2$ + 3x - 1) bằng:\par
\hspace*{1.5em}\textbf{A.} 37\quad\quad \textbf{B.} 0\quad\quad \textbf{C.} 21\quad\quad \textbf{D.} 27\par\vspace{2pt}
\textbf{\color{primary} Câu 9.} Kiểm tra nhanh kiến thức nền tảng: Tổng của cấp số nhân lùi vô hạn có số hạng đầu u$_1$ = 1 và công bội q = 1/2 bằng:\par
\hspace*{1.5em}\textbf{A.} S = 3/2\quad\quad \textbf{B.} S = 1\quad\quad \textbf{C.} S = 4\quad\quad \textbf{D.} S = 2\par\vspace{2pt}
\textbf{\color{primary} Câu 10.} Cho dữ kiện nhận biết sau đây: Giá trị của giới hạn dãy số lim ((11n + 3) / (n + 1)) khi n $\to$ +$\infty$ bằng. Phương án chính xác là:\par
\hspace*{1.5em}\textbf{A.} 1\quad\quad \textbf{B.} 11\quad\quad \textbf{C.} 0\quad\quad \textbf{D.} 3\par\vspace{2pt}
\textbf{\color{primary} Câu 11.} Hãy chỉ ra kết quả đúng trong bốn lựa chọn A, B, C, D đối với yêu cầu: Giá trị của giới hạn hàm số lim(x $\to$ 2) (11x$^2$ + 3x - 1) bằng:\par
\hspace*{1.5em}\textbf{A.} 39\quad\quad \textbf{B.} 0\quad\quad \textbf{C.} 27\quad\quad \textbf{D.} 49\par\vspace{2pt}
\textbf{\color{primary} Câu 12.} Câu hỏi ôn tập lý thuyết và kỹ năng cơ bản (Chương V): Tổng của cấp số nhân lùi vô hạn có số hạng đầu u$_1$ = 1 và công bội q = 1/2 bằng:\par
\hspace*{1.5em}\textbf{A.} S = 4\quad\quad \textbf{B.} S = 1\quad\quad \textbf{C.} S = 2\quad\quad \textbf{D.} S = 3/2\par\vspace{2pt}
\textbf{\color{primary} Câu 13.} Theo nội dung sách giáo khoa Kết nối tri thức với cuộc sống (Chương V), giá trị của giới hạn dãy số lim ((14n + 3) / (n + 1)) khi n $\to$ +$\infty$ bằng:\par
\hspace*{1.5em}\textbf{A.} 3\quad\quad \textbf{B.} 14\quad\quad \textbf{C.} 1\quad\quad \textbf{D.} 0\par\vspace{2pt}
\textbf{\color{primary} Câu 14.} Tìm đáp số đúng của bài toán cơ bản sau: Giá trị của giới hạn hàm số lim(x $\to$ 2) (14x$^2$ + 3x - 1) bằng:\par
\hspace*{1.5em}\textbf{A.} 61\quad\quad \textbf{B.} 33\quad\quad \textbf{C.} 51\quad\quad \textbf{D.} 0\par\vspace{2pt}
\textbf{\color{primary} Câu 15.} Xét đối tượng toán học trong Chương V: Tổng của cấp số nhân lùi vô hạn có số hạng đầu u$_1$ = 1 và công bội q = 1/2 bằng. Giá trị hoặc tính chất tương ứng là:\par
\hspace*{1.5em}\textbf{A.} S = 4\quad\quad \textbf{B.} S = 1\quad\quad \textbf{C.} S = 3/2\quad\quad \textbf{D.} S = 2\par\vspace{2pt}
\textbf{\color{primary} Câu 16.} Hoàn thành phép xác định cơ bản sau đây: Giá trị của giới hạn dãy số lim ((17n + 3) / (n + 1)) khi n $\to$ +$\infty$ bằng:\par
\hspace*{1.5em}\textbf{A.} 1\quad\quad \textbf{B.} 17\quad\quad \textbf{C.} 3\quad\quad \textbf{D.} 0\par\vspace{2pt}
\textbf{\color{primary} Câu 17.} Trong các khẳng định hoặc giá trị dưới đây, đâu là đáp án đúng cho câu hỏi: Giá trị của giới hạn hàm số lim(x $\to$ 2) (17x$^2$ + 3x - 1) bằng?\par
\hspace*{1.5em}\textbf{A.} 0\quad\quad \textbf{B.} 73\quad\quad \textbf{C.} 63\quad\quad \textbf{D.} 39\par\vspace{2pt}
\textbf{\color{primary} Câu 18.} Bài tập khởi động Chương V: Tổng của cấp số nhân lùi vô hạn có số hạng đầu u$_1$ = 1 và công bội q = 1/2 bằng:\par
\hspace*{1.5em}\textbf{A.} S = 4\quad\quad \textbf{B.} S = 2\quad\quad \textbf{C.} S = 1\quad\quad \textbf{D.} S = 3/2\par\vspace{2pt}
\textbf{\color{primary} Câu 19.} Quan sát dữ kiện đề bài: Giá trị của giới hạn dãy số lim ((20n + 3) / (n + 1)) khi n $\to$ +$\infty$ bằng. Hãy chọn phương án đúng:\par
\hspace*{1.5em}\textbf{A.} 1\quad\quad \textbf{B.} 0\quad\quad \textbf{C.} 3\quad\quad \textbf{D.} 20\par\vspace{2pt}
\textbf{\color{primary} Câu 20.} Sử dụng công thức hoặc tính chất đầu bài của Chương V để giải quyết yêu cầu: Giá trị của giới hạn hàm số lim(x $\to$ 2) (20x$^2$ + 3x - 1) bằng:\par
\hspace*{1.5em}\textbf{A.} 45\quad\quad \textbf{B.} 85\quad\quad \textbf{C.} 0\quad\quad \textbf{D.} 75\par\vspace{2pt}
\subsection*{\color{accent} Phần II. Mức độ Trung bình (Thông hiểu -- Câu 21 đến Câu 40)}
\textbf{\color{primary} Câu 21.} Giá trị của giới hạn vô định lim(x $\to$ 2) ((x$^2$ - 4) / (x - 2)) bằng:\par
\hspace*{1.5em}\textbf{A.} 4\quad\quad \textbf{B.} 0\quad\quad \textbf{C.} 2\quad\quad \textbf{D.} +$\infty$\par\vspace{2pt}
\textbf{\color{primary} Câu 22.} Vận dụng quy tắc trong Chương V (Giới hạn. Hàm số liên tục) để giải bài toán: Giá trị của giới hạn lim(x $\to$ +$\infty$) ((5x - 5) / (x + 2)) bằng:\par
\hspace*{1.5em}\textbf{A.} 5\quad\quad \textbf{B.} 1\quad\quad \textbf{C.} 0\quad\quad \textbf{D.} -5/2\par\vspace{2pt}
\textbf{\color{primary} Câu 23.} Biến đổi và tính toán theo kiến thức Chương V: Hàm số phân thức y = (x + 1)/(x - 5) gián đoạn tại điểm nào dưới đây. Kết quả thu được là:\par
\hspace*{1.5em}\textbf{A.} Liên tục trên toàn bộ $\mathbb{R}$\quad\quad \textbf{B.} x = -5\quad\quad \textbf{C.} x = 5\quad\quad \textbf{D.} x = -1\par\vspace{2pt}
\textbf{\color{primary} Câu 24.} Xét bài toán mức độ thông hiểu sau: Giá trị của giới hạn vô định lim(x $\to$ 5) ((x$^2$ - 25) / (x - 5)) bằng. Đáp án đúng là:\par
\hspace*{1.5em}\textbf{A.} 10\quad\quad \textbf{B.} 5\quad\quad \textbf{C.} 0\quad\quad \textbf{D.} +$\infty$\par\vspace{2pt}
\textbf{\color{primary} Câu 25.} Để hoàn thành phiếu học tập Chương V, hãy xác định: Giá trị của giới hạn lim(x $\to$ +$\infty$) ((11x - 5) / (x + 2)) bằng:\par
\hspace*{1.5em}\textbf{A.} -5/2\quad\quad \textbf{B.} 1\quad\quad \textbf{C.} 0\quad\quad \textbf{D.} 11\par\vspace{2pt}
\textbf{\color{primary} Câu 26.} Phân tích dữ kiện bài toán: Hàm số phân thức y = (x + 1)/(x - 8) gián đoạn tại điểm nào dưới đây. Giá trị cần tìm bằng:\par
\hspace*{1.5em}\textbf{A.} x = 8\quad\quad \textbf{B.} Liên tục trên toàn bộ $\mathbb{R}$\quad\quad \textbf{C.} x = -1\quad\quad \textbf{D.} x = -8\par\vspace{2pt}
\textbf{\color{primary} Câu 27.} Áp dụng tính chất đã học ở Chương V vào bài tập sau: Giá trị của giới hạn vô định lim(x $\to$ 8) ((x$^2$ - 64) / (x - 8)) bằng:\par
\hspace*{1.5em}\textbf{A.} 8\quad\quad \textbf{B.} +$\infty$\quad\quad \textbf{C.} 0\quad\quad \textbf{D.} 16\par\vspace{2pt}
\textbf{\color{primary} Câu 28.} Thực hiện các bước suy luận hợp lý để giải quyết yêu cầu: Giá trị của giới hạn lim(x $\to$ +$\infty$) ((17x - 5) / (x + 2)) bằng:\par
\hspace*{1.5em}\textbf{A.} -5/2\quad\quad \textbf{B.} 17\quad\quad \textbf{C.} 1\quad\quad \textbf{D.} 0\par\vspace{2pt}
\textbf{\color{primary} Câu 29.} Cho bài toán tự luyện (Giới hạn. Hàm số liên tục): Hàm số phân thức y = (x + 1)/(x - 11) gián đoạn tại điểm nào dưới đây. Phương án nào sau đây chính xác?\par
\hspace*{1.5em}\textbf{A.} x = 11\quad\quad \textbf{B.} Liên tục trên toàn bộ $\mathbb{R}$\quad\quad \textbf{C.} x = -1\quad\quad \textbf{D.} x = -11\par\vspace{2pt}
\textbf{\color{primary} Câu 30.} Kết hợp các phép biến đổi trong Chương V, hãy tìm kết quả của bài toán: Giá trị của giới hạn vô định lim(x $\to$ 11) ((x$^2$ - 121) / (x - 11)) bằng:\par
\hspace*{1.5em}\textbf{A.} 0\quad\quad \textbf{B.} 11\quad\quad \textbf{C.} +$\infty$\quad\quad \textbf{D.} 22\par\vspace{2pt}
\textbf{\color{primary} Câu 31.} Một bài kiểm tra định kỳ nội dung Chương V nêu câu hỏi: Giá trị của giới hạn lim(x $\to$ +$\infty$) ((23x - 5) / (x + 2)) bằng. Hãy chọn đáp số đúng:\par
\hspace*{1.5em}\textbf{A.} 0\quad\quad \textbf{B.} 1\quad\quad \textbf{C.} -5/2\quad\quad \textbf{D.} 23\par\vspace{2pt}
\textbf{\color{primary} Câu 32.} Tìm lời giải đúng cho bài toán thông hiểu: Hàm số phân thức y = (x + 1)/(x - 14) gián đoạn tại điểm nào dưới đây:\par
\hspace*{1.5em}\textbf{A.} x = -1\quad\quad \textbf{B.} Liên tục trên toàn bộ $\mathbb{R}$\quad\quad \textbf{C.} x = 14\quad\quad \textbf{D.} x = -14\par\vspace{2pt}
\textbf{\color{primary} Câu 33.} Sử dụng mối liên hệ giữa các đại lượng trong Chương V: Giá trị của giới hạn vô định lim(x $\to$ 14) ((x$^2$ - 196) / (x - 14)) bằng:\par
\hspace*{1.5em}\textbf{A.} 14\quad\quad \textbf{B.} +$\infty$\quad\quad \textbf{C.} 28\quad\quad \textbf{D.} 0\par\vspace{2pt}
\textbf{\color{primary} Câu 34.} Hãy tính toán và chọn phương án đúng đối với bài tập: Giá trị của giới hạn lim(x $\to$ +$\infty$) ((29x - 5) / (x + 2)) bằng:\par
\hspace*{1.5em}\textbf{A.} -5/2\quad\quad \textbf{B.} 0\quad\quad \textbf{C.} 1\quad\quad \textbf{D.} 29\par\vspace{2pt}
\textbf{\color{primary} Câu 35.} Dựa trên các tính chất trọng tâm của Chương V, xét bài toán: Hàm số phân thức y = (x + 1)/(x - 17) gián đoạn tại điểm nào dưới đây:\par
\hspace*{1.5em}\textbf{A.} x = -1\quad\quad \textbf{B.} x = -17\quad\quad \textbf{C.} Liên tục trên toàn bộ $\mathbb{R}$\quad\quad \textbf{D.} x = 17\par\vspace{2pt}
\textbf{\color{primary} Câu 36.} Thực hành giải toán (Giới hạn. Hàm số liên tục): Giá trị của giới hạn vô định lim(x $\to$ 17) ((x$^2$ - 289) / (x - 17)) bằng. Đáp án nào dưới đây thỏa mãn?\par
\hspace*{1.5em}\textbf{A.} +$\infty$\quad\quad \textbf{B.} 17\quad\quad \textbf{C.} 0\quad\quad \textbf{D.} 34\par\vspace{2pt}
\textbf{\color{primary} Câu 37.} Kiểm tra kỹ năng biến đổi và suy luận qua bài toán: Giá trị của giới hạn lim(x $\to$ +$\infty$) ((35x - 5) / (x + 2)) bằng:\par
\hspace*{1.5em}\textbf{A.} 0\quad\quad \textbf{B.} 1\quad\quad \textbf{C.} 35\quad\quad \textbf{D.} -5/2\par\vspace{2pt}
\textbf{\color{primary} Câu 38.} Cho giả thiết thuộc phạm vi Chương V: Hàm số phân thức y = (x + 1)/(x - 20) gián đoạn tại điểm nào dưới đây. Kết luận đúng là:\par
\hspace*{1.5em}\textbf{A.} x = -1\quad\quad \textbf{B.} Liên tục trên toàn bộ $\mathbb{R}$\quad\quad \textbf{C.} x = 20\quad\quad \textbf{D.} x = -20\par\vspace{2pt}
\textbf{\color{primary} Câu 39.} Rút gọn hoặc tính giá trị theo yêu cầu sau của Chương V: Giá trị của giới hạn vô định lim(x $\to$ 20) ((x$^2$ - 400) / (x - 20)) bằng:\par
\hspace*{1.5em}\textbf{A.} 40\quad\quad \textbf{B.} 20\quad\quad \textbf{C.} +$\infty$\quad\quad \textbf{D.} 0\par\vspace{2pt}
\textbf{\color{primary} Câu 40.} Xác định phương án đúng sau khi thực hiện phép tính / suy luận: Giá trị của giới hạn lim(x $\to$ +$\infty$) ((41x - 5) / (x + 2)) bằng:\par
\hspace*{1.5em}\textbf{A.} 1\quad\quad \textbf{B.} -5/2\quad\quad \textbf{C.} 41\quad\quad \textbf{D.} 0\par\vspace{2pt}
\subsection*{\color{accent} Phần III. Mức độ Khó (Vận dụng \& Vận dụng cao -- Câu 41 đến Câu 50)}
\textbf{\color{primary} Câu 41.} Tìm giá trị của tham số m để hàm số f(x) = (x$^2$ - 9)/(x - 3) khi x $\ne$ 3 và f(x) = m + 1 khi x = 3 liên tục tại điểm x$_0$ = 3:\par
\hspace*{1.5em}\textbf{A.} m = 7\quad\quad \textbf{B.} m = -5\quad\quad \textbf{C.} m = 6\quad\quad \textbf{D.} m = 5\par\vspace{2pt}
\textbf{\color{primary} Câu 42.} Bài toán vận dụng (Chương V – Giới hạn. Hàm số liên tục): Giá trị của giới hạn chứa căn lim(x $\to$ 0) (($\sqrt{x + 4}$ - 2) / x) bằng. Kết quả đúng là:\par
\hspace*{1.5em}\textbf{A.} 1/4\quad\quad \textbf{B.} 0\quad\quad \textbf{C.} 1/2\quad\quad \textbf{D.} 4\par\vspace{2pt}
\textbf{\color{primary} Câu 43.} Vận dụng kiến thức nâng cao trong phạm vi Chương V để giải bài toán: Tìm giá trị của tham số m để hàm số f(x) = (x$^2$ - 9)/(x - 3) khi x $\ne$ 3 và f(x) = m + 3 khi x = 3 liên tục tại điểm x$_0$ = 3:\par
\hspace*{1.5em}\textbf{A.} m = 6\quad\quad \textbf{B.} m = -3\quad\quad \textbf{C.} m = 3\quad\quad \textbf{D.} m = 9\par\vspace{2pt}
\textbf{\color{primary} Câu 44.} Phân tích và giải bài toán tổng hợp sau của Chương V: Giá trị của giới hạn chứa căn lim(x $\to$ 0) (($\sqrt{x + 4}$ - 2) / x) bằng:\par
\hspace*{1.5em}\textbf{A.} 4\quad\quad \textbf{B.} 1/2\quad\quad \textbf{C.} 0\quad\quad \textbf{D.} 1/4\par\vspace{2pt}
\textbf{\color{primary} Câu 45.} Thử thách tư duy trong phạm vi kiến thức Chương V: Tìm giá trị của tham số m để hàm số f(x) = (x$^2$ - 9)/(x - 3) khi x $\ne$ 3 và f(x) = m + 5 khi x = 3 liên tục tại điểm x$_0$ = 3. Đáp số chính xác là:\par
\hspace*{1.5em}\textbf{A.} m = 1\quad\quad \textbf{B.} m = 6\quad\quad \textbf{C.} m = -1\quad\quad \textbf{D.} m = 11\par\vspace{2pt}
\textbf{\color{primary} Câu 46.} Sử dụng phương pháp giải toán của Chương V để tìm đáp án cho bài toán: Giá trị của giới hạn chứa căn lim(x $\to$ 0) (($\sqrt{x + 4}$ - 2) / x) bằng:\par
\hspace*{1.5em}\textbf{A.} 1/2\quad\quad \textbf{B.} 1/4\quad\quad \textbf{C.} 0\quad\quad \textbf{D.} 4\par\vspace{2pt}
\textbf{\color{primary} Câu 47.} Xét bài toán vận dụng thực tiễn / nâng cao thuộc Chương V: Tìm giá trị của tham số m để hàm số f(x) = (x$^2$ - 9)/(x - 3) khi x $\ne$ 3 và f(x) = m + 7 khi x = 3 liên tục tại điểm x$_0$ = 3:\par
\hspace*{1.5em}\textbf{A.} m = --1\quad\quad \textbf{B.} m = 13\quad\quad \textbf{C.} m = 6\quad\quad \textbf{D.} m = -1\par\vspace{2pt}
\textbf{\color{primary} Câu 48.} Tìm kết quả cuối cùng của bài toán phân loại học sinh khá giỏi (Chương V): Giá trị của giới hạn chứa căn lim(x $\to$ 0) (($\sqrt{x + 4}$ - 2) / x) bằng:\par
\hspace*{1.5em}\textbf{A.} 0\quad\quad \textbf{B.} 1/2\quad\quad \textbf{C.} 4\quad\quad \textbf{D.} 1/4\par\vspace{2pt}
\textbf{\color{primary} Câu 49.} Kết hợp các tính chất đã học trong Chương V, hãy giải quyết bài toán: Tìm giá trị của tham số m để hàm số f(x) = (x$^2$ - 9)/(x - 3) khi x $\ne$ 3 và f(x) = m + 9 khi x = 3 liên tục tại điểm x$_0$ = 3:\par
\hspace*{1.5em}\textbf{A.} m = --3\quad\quad \textbf{B.} m = 6\quad\quad \textbf{C.} m = 15\quad\quad \textbf{D.} m = -3\par\vspace{2pt}
\textbf{\color{primary} Câu 50.} Câu hỏi vận dụng cao cuối Chương V: Giá trị của giới hạn chứa căn lim(x $\to$ 0) (($\sqrt{x + 4}$ - 2) / x) bằng. Phương án đúng là:\par
\hspace*{1.5em}\textbf{A.} 0\quad\quad \textbf{B.} 1/4\quad\quad \textbf{C.} 1/2\quad\quad \textbf{D.} 4\par\vspace{2pt}
\subsection*{\color{primary} Bảng đáp án Chương V: Giới hạn. Hàm số liên tục}
\begin{center}
\renewcommand{\arraystretch}{1.25}
\begin{tabularx}{\textwidth}{*{10}{>{\centering\arraybackslash}X}}
\toprule
\textbf{\color{primary} Cột 1} & \textbf{\color{primary} Cột 2} & \textbf{\color{primary} Cột 3} & \textbf{\color{primary} Cột 4} & \textbf{\color{primary} Cột 5} & \textbf{\color{primary} Cột 6} & \textbf{\color{primary} Cột 7} & \textbf{\color{primary} Cột 8} & \textbf{\color{primary} Cột 9} & \textbf{\color{primary} Cột 10} \\
\midrule
\textbf{1.A} & \textbf{2.D} & \textbf{3.C} & \textbf{4.A} & \textbf{5.A} & \textbf{6.B} & \textbf{7.C} & \textbf{8.A} & \textbf{9.D} & \textbf{10.B} \\
\rowcolor{rowtint} \textbf{11.D} & \textbf{12.C} & \textbf{13.B} & \textbf{14.A} & \textbf{15.D} & \textbf{16.B} & \textbf{17.B} & \textbf{18.B} & \textbf{19.D} & \textbf{20.B} \\
\textbf{21.A} & \textbf{22.A} & \textbf{23.C} & \textbf{24.A} & \textbf{25.D} & \textbf{26.A} & \textbf{27.D} & \textbf{28.B} & \textbf{29.A} & \textbf{30.D} \\
\rowcolor{rowtint} \textbf{31.D} & \textbf{32.C} & \textbf{33.C} & \textbf{34.D} & \textbf{35.D} & \textbf{36.D} & \textbf{37.C} & \textbf{38.C} & \textbf{39.A} & \textbf{40.C} \\
\textbf{41.D} & \textbf{42.A} & \textbf{43.C} & \textbf{44.D} & \textbf{45.A} & \textbf{46.B} & \textbf{47.D} & \textbf{48.D} & \textbf{49.D} & \textbf{50.B} \\
\bottomrule
\end{tabularx}
\end{center}\vspace{10pt}
\section*{\color{primary} CHƯƠNG VI: HÀM SỐ MŨ VÀ HÀM SỐ LÔGARIT}
\textit{Trọng tâm kiến thức SGK: Lũy thừa với số mũ thực, lôgarit và các quy tắc tính lôgarit, hàm số mũ và hàm số lôgarit, phương trình và bất phương trình mũ – lôgarit.}\par\vspace{4pt}
\subsection*{\color{accent} Phần I. Mức độ Dễ (Nhận biết -- Câu 1 đến Câu 20)}
\textbf{\color{primary} Câu 1.} Với a > 0, a $\ne$ 1, giá trị của $\log_a$(a$^{3}$) bằng:\par
\hspace*{1.5em}\textbf{A.} 1\quad\quad \textbf{B.} 3\quad\quad \textbf{C.} a$^{3}$\quad\quad \textbf{D.} 0\par\vspace{2pt}
\textbf{\color{primary} Câu 2.} Thực hiện yêu cầu nhận biết sau: Tập xác định D của hàm số lôgarit y = $\log_2$(x - 3) là.\par
\hspace*{1.5em}\textbf{A.} D = $\mathbb{R}$ \textbackslash\{\} \{3\}\quad\quad \textbf{B.} D = (3; +$\infty$)\quad\quad \textbf{C.} D = [3; +$\infty$)\quad\quad \textbf{D.} D = (0; +$\infty$)\par\vspace{2pt}
\textbf{\color{primary} Câu 3.} Trong bài học thuộc Chương VI, xét nội dung: Nghiệm của phương trình mũ cơ bản 2$^{x - 1}$ = 16 là. Đáp án đúng là:\par
\hspace*{1.5em}\textbf{A.} x = 5\quad\quad \textbf{B.} x = 4\quad\quad \textbf{C.} x = 3\quad\quad \textbf{D.} x = 9\par\vspace{2pt}
\textbf{\color{primary} Câu 4.} Dựa vào kiến thức chuẩn SGK của Chương VI (Hàm số mũ và hàm số lôgarit), với a > 0, a $\ne$ 1, giá trị của $\log_a$(a$^{6}$) bằng:\par
\hspace*{1.5em}\textbf{A.} a$^{6}$\quad\quad \textbf{B.} 6\quad\quad \textbf{C.} 0\quad\quad \textbf{D.} 1\par\vspace{2pt}
\textbf{\color{primary} Câu 5.} Chọn kết quả đúng nhất khi xét: Tập xác định D của hàm số lôgarit y = $\log_2$(x - 6) là:\par
\hspace*{1.5em}\textbf{A.} D = (6; +$\infty$)\quad\quad \textbf{B.} D = [6; +$\infty$)\quad\quad \textbf{C.} D = (0; +$\infty$)\quad\quad \textbf{D.} D = $\mathbb{R}$ \textbackslash\{\} \{6\}\par\vspace{2pt}
\textbf{\color{primary} Câu 6.} Khi làm bài tập cơ bản về hàm số mũ và hàm số lôgarit, một học sinh gặp câu hỏi: Nghiệm của phương trình mũ cơ bản 2$^{x - 1}$ = 8 là. Kết quả nào dưới đây đúng?\par
\hspace*{1.5em}\textbf{A.} x = 4\quad\quad \textbf{B.} x = 3\quad\quad \textbf{C.} x = 5\quad\quad \textbf{D.} x = 2\par\vspace{2pt}
\textbf{\color{primary} Câu 7.} Xác định giá trị hoặc khẳng định đúng cho bài toán sau: Với a > 0, a $\ne$ 1, giá trị của $\log_a$(a$^{9}$) bằng:\par
\hspace*{1.5em}\textbf{A.} a$^{9}$\quad\quad \textbf{B.} 0\quad\quad \textbf{C.} 1\quad\quad \textbf{D.} 9\par\vspace{2pt}
\textbf{\color{primary} Câu 8.} Áp dụng trực tiếp định nghĩa và quy tắc trong Chương VI: Tập xác định D của hàm số lôgarit y = $\log_2$(x - 9) là:\par
\hspace*{1.5em}\textbf{A.} D = (9; +$\infty$)\quad\quad \textbf{B.} D = (0; +$\infty$)\quad\quad \textbf{C.} D = $\mathbb{R}$ \textbackslash\{\} \{9\}\quad\quad \textbf{D.} D = [9; +$\infty$)\par\vspace{2pt}
\textbf{\color{primary} Câu 9.} Kiểm tra nhanh kiến thức nền tảng: Nghiệm của phương trình mũ cơ bản 2$^{x - 1}$ = 4 là:\par
\hspace*{1.5em}\textbf{A.} x = 3\quad\quad \textbf{B.} x = 2\quad\quad \textbf{C.} x = 1\quad\quad \textbf{D.} x = 3 (1)\par\vspace{2pt}
\textbf{\color{primary} Câu 10.} Cho dữ kiện nhận biết sau đây: Với a > 0, a $\ne$ 1, giá trị của $\log_a$(a$^{12}$) bằng. Phương án chính xác là:\par
\hspace*{1.5em}\textbf{A.} 1\quad\quad \textbf{B.} a$^{12}$\quad\quad \textbf{C.} 12\quad\quad \textbf{D.} 0\par\vspace{2pt}
\textbf{\color{primary} Câu 11.} Hãy chỉ ra kết quả đúng trong bốn lựa chọn A, B, C, D đối với yêu cầu: Tập xác định D của hàm số lôgarit y = $\log_2$(x - 12) là:\par
\hspace*{1.5em}\textbf{A.} D = $\mathbb{R}$ \textbackslash\{\} \{12\}\quad\quad \textbf{B.} D = [12; +$\infty$)\quad\quad \textbf{C.} D = (12; +$\infty$)\quad\quad \textbf{D.} D = (0; +$\infty$)\par\vspace{2pt}
\textbf{\color{primary} Câu 12.} Câu hỏi ôn tập lý thuyết và kỹ năng cơ bản (Chương VI): Nghiệm của phương trình mũ cơ bản 2$^{x - 1}$ = 32 là:\par
\hspace*{1.5em}\textbf{A.} x = 5\quad\quad \textbf{B.} x = 6\quad\quad \textbf{C.} x = 4\quad\quad \textbf{D.} x = 17\par\vspace{2pt}
\textbf{\color{primary} Câu 13.} Theo nội dung sách giáo khoa Kết nối tri thức với cuộc sống (Chương VI), với a > 0, a $\ne$ 1, giá trị của $\log_a$(a$^{15}$) bằng:\par
\hspace*{1.5em}\textbf{A.} 0\quad\quad \textbf{B.} 15\quad\quad \textbf{C.} a$^{15}$\quad\quad \textbf{D.} 1\par\vspace{2pt}
\textbf{\color{primary} Câu 14.} Tìm đáp số đúng của bài toán cơ bản sau: Tập xác định D của hàm số lôgarit y = $\log_2$(x - 15) là:\par
\hspace*{1.5em}\textbf{A.} D = [15; +$\infty$)\quad\quad \textbf{B.} D = $\mathbb{R}$ \textbackslash\{\} \{15\}\quad\quad \textbf{C.} D = (15; +$\infty$)\quad\quad \textbf{D.} D = (0; +$\infty$)\par\vspace{2pt}
\textbf{\color{primary} Câu 15.} Xét đối tượng toán học trong Chương VI: Nghiệm của phương trình mũ cơ bản 2$^{x - 1}$ = 16 là. Giá trị hoặc tính chất tương ứng là:\par
\hspace*{1.5em}\textbf{A.} x = 3\quad\quad \textbf{B.} x = 4\quad\quad \textbf{C.} x = 9\quad\quad \textbf{D.} x = 5\par\vspace{2pt}
\textbf{\color{primary} Câu 16.} Hoàn thành phép xác định cơ bản sau đây: Với a > 0, a $\ne$ 1, giá trị của $\log_a$(a$^{18}$) bằng:\par
\hspace*{1.5em}\textbf{A.} 0\quad\quad \textbf{B.} 1\quad\quad \textbf{C.} 18\quad\quad \textbf{D.} a$^{18}$\par\vspace{2pt}
\textbf{\color{primary} Câu 17.} Trong các khẳng định hoặc giá trị dưới đây, đâu là đáp án đúng cho câu hỏi: Tập xác định D của hàm số lôgarit y = $\log_2$(x - 18) là?\par
\hspace*{1.5em}\textbf{A.} D = [18; +$\infty$)\quad\quad \textbf{B.} D = (18; +$\infty$)\quad\quad \textbf{C.} D = $\mathbb{R}$ \textbackslash\{\} \{18\}\quad\quad \textbf{D.} D = (0; +$\infty$)\par\vspace{2pt}
\textbf{\color{primary} Câu 18.} Bài tập khởi động Chương VI: Nghiệm của phương trình mũ cơ bản 2$^{x - 1}$ = 8 là:\par
\hspace*{1.5em}\textbf{A.} x = 5\quad\quad \textbf{B.} x = 2\quad\quad \textbf{C.} x = 3\quad\quad \textbf{D.} x = 4\par\vspace{2pt}
\textbf{\color{primary} Câu 19.} Quan sát dữ kiện đề bài: Với a > 0, a $\ne$ 1, giá trị của $\log_a$(a$^{21}$) bằng. Hãy chọn phương án đúng:\par
\hspace*{1.5em}\textbf{A.} 1\quad\quad \textbf{B.} 0\quad\quad \textbf{C.} a$^{21}$\quad\quad \textbf{D.} 21\par\vspace{2pt}
\textbf{\color{primary} Câu 20.} Sử dụng công thức hoặc tính chất đầu bài của Chương VI để giải quyết yêu cầu: Tập xác định D của hàm số lôgarit y = $\log_2$(x - 21) là:\par
\hspace*{1.5em}\textbf{A.} D = (0; +$\infty$)\quad\quad \textbf{B.} D = $\mathbb{R}$ \textbackslash\{\} \{21\}\quad\quad \textbf{C.} D = (21; +$\infty$)\quad\quad \textbf{D.} D = [21; +$\infty$)\par\vspace{2pt}
\subsection*{\color{accent} Phần II. Mức độ Trung bình (Thông hiểu -- Câu 21 đến Câu 40)}
\textbf{\color{primary} Câu 21.} Cho $\log_2$ a = 2 và $\log_2$ b = 3. Giá trị của biểu thức $\log_2$(a$^2$ $\cdot$ b) bằng:\par
\hspace*{1.5em}\textbf{A.} 7 (1)\quad\quad \textbf{B.} 12\quad\quad \textbf{C.} 7\quad\quad \textbf{D.} 1\par\vspace{2pt}
\textbf{\color{primary} Câu 22.} Vận dụng quy tắc trong Chương VI (Hàm số mũ và hàm số lôgarit) để giải bài toán: Nghiệm của phương trình $\log_3$(2x - 1) = 2 là:\par
\hspace*{1.5em}\textbf{A.} x = 4\quad\quad \textbf{B.} x = 9/2\quad\quad \textbf{C.} x = 7/2\quad\quad \textbf{D.} x = 5\par\vspace{2pt}
\textbf{\color{primary} Câu 23.} Biến đổi và tính toán theo kiến thức Chương VI: Tập nghiệm của bất phương trình mũ (1/2)$^{x - 3}$ > (1/2)$^3$ (chú ý cơ số 0 < 1/2 < 1) là. Kết quả thu được là:\par
\hspace*{1.5em}\textbf{A.} S = (0; 6)\quad\quad \textbf{B.} S = (-$\infty$; 6)\quad\quad \textbf{C.} S = (-$\infty$; 4)\quad\quad \textbf{D.} S = (6; +$\infty$)\par\vspace{2pt}
\textbf{\color{primary} Câu 24.} Xét bài toán mức độ thông hiểu sau: Cho $\log_2$ a = 5 và $\log_2$ b = 3. Giá trị của biểu thức $\log_2$(a$^2$ $\cdot$ b) bằng. Đáp án đúng là:\par
\hspace*{1.5em}\textbf{A.} 13\quad\quad \textbf{B.} 7\quad\quad \textbf{C.} 30\quad\quad \textbf{D.} 28\par\vspace{2pt}
\textbf{\color{primary} Câu 25.} Để hoàn thành phiếu học tập Chương VI, hãy xác định: Nghiệm của phương trình $\log_3$(2x - 1) = 2 là:\par
\hspace*{1.5em}\textbf{A.} x = 5\quad\quad \textbf{B.} x = 7/2\quad\quad \textbf{C.} x = 4\quad\quad \textbf{D.} x = 9/2\par\vspace{2pt}
\textbf{\color{primary} Câu 26.} Phân tích dữ kiện bài toán: Tập nghiệm của bất phương trình mũ (1/2)$^{x - 6}$ > (1/2)$^3$ (chú ý cơ số 0 < 1/2 < 1) là. Giá trị cần tìm bằng:\par
\hspace*{1.5em}\textbf{A.} S = (-$\infty$; 7)\quad\quad \textbf{B.} S = (9; +$\infty$)\quad\quad \textbf{C.} S = (0; 9)\quad\quad \textbf{D.} S = (-$\infty$; 9)\par\vspace{2pt}
\textbf{\color{primary} Câu 27.} Áp dụng tính chất đã học ở Chương VI vào bài tập sau: Cho $\log_2$ a = 8 và $\log_2$ b = 3. Giá trị của biểu thức $\log_2$(a$^2$ $\cdot$ b) bằng:\par
\hspace*{1.5em}\textbf{A.} 13\quad\quad \textbf{B.} 48\quad\quad \textbf{C.} 67\quad\quad \textbf{D.} 19\par\vspace{2pt}
\textbf{\color{primary} Câu 28.} Thực hiện các bước suy luận hợp lý để giải quyết yêu cầu: Nghiệm của phương trình $\log_3$(2x - 1) = 2 là:\par
\hspace*{1.5em}\textbf{A.} x = 7/2\quad\quad \textbf{B.} x = 5\quad\quad \textbf{C.} x = 9/2\quad\quad \textbf{D.} x = 4\par\vspace{2pt}
\textbf{\color{primary} Câu 29.} Cho bài toán tự luyện (Hàm số mũ và hàm số lôgarit): Tập nghiệm của bất phương trình mũ (1/2)$^{x - 9}$ > (1/2)$^3$ (chú ý cơ số 0 < 1/2 < 1) là. Phương án nào sau đây chính xác?\par
\hspace*{1.5em}\textbf{A.} S = (0; 12)\quad\quad \textbf{B.} S = (12; +$\infty$)\quad\quad \textbf{C.} S = (-$\infty$; 12)\quad\quad \textbf{D.} S = (-$\infty$; 10)\par\vspace{2pt}
\textbf{\color{primary} Câu 30.} Kết hợp các phép biến đổi trong Chương VI, hãy tìm kết quả của bài toán: Cho $\log_2$ a = 11 và $\log_2$ b = 3. Giá trị của biểu thức $\log_2$(a$^2$ $\cdot$ b) bằng:\par
\hspace*{1.5em}\textbf{A.} 124\quad\quad \textbf{B.} 25\quad\quad \textbf{C.} 19\quad\quad \textbf{D.} 66\par\vspace{2pt}
\textbf{\color{primary} Câu 31.} Một bài kiểm tra định kỳ nội dung Chương VI nêu câu hỏi: Nghiệm của phương trình $\log_3$(2x - 1) = 2 là. Hãy chọn đáp số đúng:\par
\hspace*{1.5em}\textbf{A.} x = 4\quad\quad \textbf{B.} x = 9/2\quad\quad \textbf{C.} x = 5\quad\quad \textbf{D.} x = 7/2\par\vspace{2pt}
\textbf{\color{primary} Câu 32.} Tìm lời giải đúng cho bài toán thông hiểu: Tập nghiệm của bất phương trình mũ (1/2)$^{x - 12}$ > (1/2)$^3$ (chú ý cơ số 0 < 1/2 < 1) là:\par
\hspace*{1.5em}\textbf{A.} S = (-$\infty$; 13)\quad\quad \textbf{B.} S = (-$\infty$; 15)\quad\quad \textbf{C.} S = (0; 15)\quad\quad \textbf{D.} S = (15; +$\infty$)\par\vspace{2pt}
\textbf{\color{primary} Câu 33.} Sử dụng mối liên hệ giữa các đại lượng trong Chương VI: Cho $\log_2$ a = 14 và $\log_2$ b = 3. Giá trị của biểu thức $\log_2$(a$^2$ $\cdot$ b) bằng:\par
\hspace*{1.5em}\textbf{A.} 199\quad\quad \textbf{B.} 31\quad\quad \textbf{C.} 84\quad\quad \textbf{D.} 25\par\vspace{2pt}
\textbf{\color{primary} Câu 34.} Hãy tính toán và chọn phương án đúng đối với bài tập: Nghiệm của phương trình $\log_3$(2x - 1) = 2 là:\par
\hspace*{1.5em}\textbf{A.} x = 7/2\quad\quad \textbf{B.} x = 4\quad\quad \textbf{C.} x = 9/2\quad\quad \textbf{D.} x = 5\par\vspace{2pt}
\textbf{\color{primary} Câu 35.} Dựa trên các tính chất trọng tâm của Chương VI, xét bài toán: Tập nghiệm của bất phương trình mũ (1/2)$^{x - 15}$ > (1/2)$^3$ (chú ý cơ số 0 < 1/2 < 1) là:\par
\hspace*{1.5em}\textbf{A.} S = (18; +$\infty$)\quad\quad \textbf{B.} S = (-$\infty$; 16)\quad\quad \textbf{C.} S = (-$\infty$; 18)\quad\quad \textbf{D.} S = (0; 18)\par\vspace{2pt}
\textbf{\color{primary} Câu 36.} Thực hành giải toán (Hàm số mũ và hàm số lôgarit): Cho $\log_2$ a = 17 và $\log_2$ b = 3. Giá trị của biểu thức $\log_2$(a$^2$ $\cdot$ b) bằng. Đáp án nào dưới đây thỏa mãn?\par
\hspace*{1.5em}\textbf{A.} 292\quad\quad \textbf{B.} 37\quad\quad \textbf{C.} 31\quad\quad \textbf{D.} 102\par\vspace{2pt}
\textbf{\color{primary} Câu 37.} Kiểm tra kỹ năng biến đổi và suy luận qua bài toán: Nghiệm của phương trình $\log_3$(2x - 1) = 2 là:\par
\hspace*{1.5em}\textbf{A.} x = 9/2\quad\quad \textbf{B.} x = 7/2\quad\quad \textbf{C.} x = 5\quad\quad \textbf{D.} x = 4\par\vspace{2pt}
\textbf{\color{primary} Câu 38.} Cho giả thiết thuộc phạm vi Chương VI: Tập nghiệm của bất phương trình mũ (1/2)$^{x - 18}$ > (1/2)$^3$ (chú ý cơ số 0 < 1/2 < 1) là. Kết luận đúng là:\par
\hspace*{1.5em}\textbf{A.} S = (21; +$\infty$)\quad\quad \textbf{B.} S = (0; 21)\quad\quad \textbf{C.} S = (-$\infty$; 21)\quad\quad \textbf{D.} S = (-$\infty$; 19)\par\vspace{2pt}
\textbf{\color{primary} Câu 39.} Rút gọn hoặc tính giá trị theo yêu cầu sau của Chương VI: Cho $\log_2$ a = 20 và $\log_2$ b = 3. Giá trị của biểu thức $\log_2$(a$^2$ $\cdot$ b) bằng:\par
\hspace*{1.5em}\textbf{A.} 120\quad\quad \textbf{B.} 37\quad\quad \textbf{C.} 403\quad\quad \textbf{D.} 43\par\vspace{2pt}
\textbf{\color{primary} Câu 40.} Xác định phương án đúng sau khi thực hiện phép tính / suy luận: Nghiệm của phương trình $\log_3$(2x - 1) = 2 là:\par
\hspace*{1.5em}\textbf{A.} x = 9/2\quad\quad \textbf{B.} x = 7/2\quad\quad \textbf{C.} x = 5\quad\quad \textbf{D.} x = 4\par\vspace{2pt}
\subsection*{\color{accent} Phần III. Mức độ Khó (Vận dụng \& Vận dụng cao -- Câu 41 đến Câu 50)}
\textbf{\color{primary} Câu 41.} Tổng các nghiệm thực của phương trình 4$^{x}$ - 6$\cdot$2$^{x}$ + 8 = 0 bằng:\par
\hspace*{1.5em}\textbf{A.} 3\quad\quad \textbf{B.} 8\quad\quad \textbf{C.} 6\quad\quad \textbf{D.} 2\par\vspace{2pt}
\textbf{\color{primary} Câu 42.} Bài toán vận dụng (Chương VI – Hàm số mũ và hàm số lôgarit): Tập nghiệm của bất phương trình $\log_{1/3}$(x - 2) $\ge$ -1 (chú ý điều kiện xác định x > 2) là. Kết quả đúng là:\par
\hspace*{1.5em}\textbf{A.} S = [5; +$\infty$)\quad\quad \textbf{B.} S = (2; 5]\quad\quad \textbf{C.} S = [2; 5]\quad\quad \textbf{D.} S = (-$\infty$; 5]\par\vspace{2pt}
\textbf{\color{primary} Câu 43.} Vận dụng kiến thức nâng cao trong phạm vi Chương VI để giải bài toán: Tập nghiệm của bất phương trình lôgarit $\log_2$(x - 4) $\le$ 3 (với điều kiện xác định x > 4) là:\par
\hspace*{1.5em}\textbf{A.} S = [4; 12]\quad\quad \textbf{B.} S = (12; +$\infty$)\quad\quad \textbf{C.} S = (4; 12]\quad\quad \textbf{D.} S = (-$\infty$; 12]\par\vspace{2pt}
\textbf{\color{primary} Câu 44.} Phân tích và giải bài toán tổng hợp sau của Chương VI: Tổng các nghiệm thực của phương trình 4$^{x}$ - 6$\cdot$2$^{x}$ + 8 = 0 bằng:\par
\hspace*{1.5em}\textbf{A.} 2\quad\quad \textbf{B.} 3\quad\quad \textbf{C.} 6\quad\quad \textbf{D.} 8\par\vspace{2pt}
\textbf{\color{primary} Câu 45.} Thử thách tư duy trong phạm vi kiến thức Chương VI: Tập nghiệm của bất phương trình $\log_{1/3}$(x - 2) $\ge$ -1 (chú ý điều kiện xác định x > 2) là. Đáp số chính xác là:\par
\hspace*{1.5em}\textbf{A.} S = [2; 5]\quad\quad \textbf{B.} S = [5; +$\infty$)\quad\quad \textbf{C.} S = (2; 5]\quad\quad \textbf{D.} S = (-$\infty$; 5]\par\vspace{2pt}
\textbf{\color{primary} Câu 46.} Sử dụng phương pháp giải toán của Chương VI để tìm đáp án cho bài toán: Tập nghiệm của bất phương trình lôgarit $\log_2$(x - 7) $\le$ 3 (với điều kiện xác định x > 7) là:\par
\hspace*{1.5em}\textbf{A.} S = (15; +$\infty$)\quad\quad \textbf{B.} S = (7; 15]\quad\quad \textbf{C.} S = [7; 15]\quad\quad \textbf{D.} S = (-$\infty$; 15]\par\vspace{2pt}
\textbf{\color{primary} Câu 47.} Xét bài toán vận dụng thực tiễn / nâng cao thuộc Chương VI: Tổng các nghiệm thực của phương trình 4$^{x}$ - 6$\cdot$2$^{x}$ + 8 = 0 bằng:\par
\hspace*{1.5em}\textbf{A.} 3\quad\quad \textbf{B.} 6\quad\quad \textbf{C.} 8\quad\quad \textbf{D.} 2\par\vspace{2pt}
\textbf{\color{primary} Câu 48.} Tìm kết quả cuối cùng của bài toán phân loại học sinh khá giỏi (Chương VI): Tập nghiệm của bất phương trình $\log_{1/3}$(x - 2) $\ge$ -1 (chú ý điều kiện xác định x > 2) là:\par
\hspace*{1.5em}\textbf{A.} S = (2; 5]\quad\quad \textbf{B.} S = [5; +$\infty$)\quad\quad \textbf{C.} S = (-$\infty$; 5]\quad\quad \textbf{D.} S = [2; 5]\par\vspace{2pt}
\textbf{\color{primary} Câu 49.} Kết hợp các tính chất đã học trong Chương VI, hãy giải quyết bài toán: Tập nghiệm của bất phương trình lôgarit $\log_2$(x - 10) $\le$ 3 (với điều kiện xác định x > 10) là:\par
\hspace*{1.5em}\textbf{A.} S = (-$\infty$; 18]\quad\quad \textbf{B.} S = (10; 18]\quad\quad \textbf{C.} S = (18; +$\infty$)\quad\quad \textbf{D.} S = [10; 18]\par\vspace{2pt}
\textbf{\color{primary} Câu 50.} Câu hỏi vận dụng cao cuối Chương VI: Tổng các nghiệm thực của phương trình 4$^{x}$ - 6$\cdot$2$^{x}$ + 8 = 0 bằng. Phương án đúng là:\par
\hspace*{1.5em}\textbf{A.} 3\quad\quad \textbf{B.} 2\quad\quad \textbf{C.} 6\quad\quad \textbf{D.} 8\par\vspace{2pt}
\subsection*{\color{primary} Bảng đáp án Chương VI: Hàm số mũ và hàm số lôgarit}
\begin{center}
\renewcommand{\arraystretch}{1.25}
\begin{tabularx}{\textwidth}{*{10}{>{\centering\arraybackslash}X}}
\toprule
\textbf{\color{primary} Cột 1} & \textbf{\color{primary} Cột 2} & \textbf{\color{primary} Cột 3} & \textbf{\color{primary} Cột 4} & \textbf{\color{primary} Cột 5} & \textbf{\color{primary} Cột 6} & \textbf{\color{primary} Cột 7} & \textbf{\color{primary} Cột 8} & \textbf{\color{primary} Cột 9} & \textbf{\color{primary} Cột 10} \\
\midrule
\textbf{1.B} & \textbf{2.B} & \textbf{3.A} & \textbf{4.B} & \textbf{5.A} & \textbf{6.A} & \textbf{7.D} & \textbf{8.A} & \textbf{9.A} & \textbf{10.C} \\
\rowcolor{rowtint} \textbf{11.C} & \textbf{12.B} & \textbf{13.B} & \textbf{14.C} & \textbf{15.D} & \textbf{16.C} & \textbf{17.B} & \textbf{18.D} & \textbf{19.D} & \textbf{20.C} \\
\textbf{21.C} & \textbf{22.D} & \textbf{23.B} & \textbf{24.A} & \textbf{25.A} & \textbf{26.D} & \textbf{27.D} & \textbf{28.B} & \textbf{29.C} & \textbf{30.B} \\
\rowcolor{rowtint} \textbf{31.C} & \textbf{32.B} & \textbf{33.B} & \textbf{34.D} & \textbf{35.C} & \textbf{36.B} & \textbf{37.C} & \textbf{38.C} & \textbf{39.D} & \textbf{40.C} \\
\textbf{41.A} & \textbf{42.B} & \textbf{43.C} & \textbf{44.B} & \textbf{45.C} & \textbf{46.B} & \textbf{47.A} & \textbf{48.A} & \textbf{49.B} & \textbf{50.A} \\
\bottomrule
\end{tabularx}
\end{center}\vspace{10pt}
\section*{\color{primary} CHƯƠNG VII: QUAN HỆ VUÔNG GÓC TRONG KHÔNG GIAN}
\textit{Trọng tâm kiến thức SGK: Hai đường thẳng vuông góc, đường thẳng vuông góc với mặt phẳng, phép chiếu vuông góc và góc giữa đường thẳng với mặt phẳng, hai mặt phẳng vuông góc, góc nhị diện, khoảng cách và thể tích khối chóp – lăng trụ.}\par\vspace{4pt}
\subsection*{\color{accent} Phần I. Mức độ Dễ (Nhận biết -- Câu 1 đến Câu 20)}
\textbf{\color{primary} Câu 1.} Đường thẳng d vuông góc với mặt phẳng (P) khi và chỉ khi d vuông góc với:\par
\begin{itemize}\setlength{\itemsep}{1pt}\setlength{\parskip}{0pt}
  \item[] \textbf{A.} Hai đường thẳng song song nằm trong (P)
  \item[] \textbf{B.} Một đường thẳng bất kỳ trong (P)
  \item[] \textbf{C.} Hai đường thẳng cắt nhau nằm trong mặt phẳng (P)
  \item[] \textbf{D.} Giao tuyến của (P) với mặt phẳng khác
\end{itemize}\vspace{2pt}
\textbf{\color{primary} Câu 2.} Thực hiện yêu cầu nhận biết sau: Một khối chóp có diện tích đáy B = 18 và chiều cao h = 4. Thể tích V = 1/3 $\cdot$ B$\cdot$h của khối chóp bằng.\par
\hspace*{1.5em}\textbf{A.} 36\quad\quad \textbf{B.} 72\quad\quad \textbf{C.} 12\quad\quad \textbf{D.} 24\par\vspace{2pt}
\textbf{\color{primary} Câu 3.} Trong bài học thuộc Chương VII, xét nội dung: Một khối lăng trụ có diện tích đáy B = 7 và chiều cao h = 6. Thể tích V = B$\cdot$h của khối lăng trụ bằng. Đáp án đúng là:\par
\hspace*{1.5em}\textbf{A.} 42\quad\quad \textbf{B.} 84\quad\quad \textbf{C.} 14.0\quad\quad \textbf{D.} 13\par\vspace{2pt}
\textbf{\color{primary} Câu 4.} Dựa vào kiến thức chuẩn SGK của Chương VII (Quan hệ vuông góc trong không gian), đường thẳng d vuông góc với mặt phẳng (P) khi và chỉ khi d vuông góc với:\par
\begin{itemize}\setlength{\itemsep}{1pt}\setlength{\parskip}{0pt}
  \item[] \textbf{A.} Một đường thẳng bất kỳ trong (P)
  \item[] \textbf{B.} Hai đường thẳng cắt nhau nằm trong mặt phẳng (P)
  \item[] \textbf{C.} Giao tuyến của (P) với mặt phẳng khác
  \item[] \textbf{D.} Hai đường thẳng song song nằm trong (P)
\end{itemize}\vspace{2pt}
\textbf{\color{primary} Câu 5.} Chọn kết quả đúng nhất khi xét: Một khối chóp có diện tích đáy B = 36 và chiều cao h = 7. Thể tích V = 1/3 $\cdot$ B$\cdot$h của khối chóp bằng:\par
\hspace*{1.5em}\textbf{A.} 126\quad\quad \textbf{B.} 252\quad\quad \textbf{C.} 84\quad\quad \textbf{D.} 42\par\vspace{2pt}
\textbf{\color{primary} Câu 6.} Khi làm bài tập cơ bản về quan hệ vuông góc trong không gian, một học sinh gặp câu hỏi: Một khối lăng trụ có diện tích đáy B = 10 và chiều cao h = 9. Thể tích V = B$\cdot$h của khối lăng trụ bằng. Kết quả nào dưới đây đúng?\par
\hspace*{1.5em}\textbf{A.} 19\quad\quad \textbf{B.} 180\quad\quad \textbf{C.} 30.0\quad\quad \textbf{D.} 90\par\vspace{2pt}
\textbf{\color{primary} Câu 7.} Xác định giá trị hoặc khẳng định đúng cho bài toán sau: Đường thẳng d vuông góc với mặt phẳng (P) khi và chỉ khi d vuông góc với:\par
\begin{itemize}\setlength{\itemsep}{1pt}\setlength{\parskip}{0pt}
  \item[] \textbf{A.} Hai đường thẳng song song nằm trong (P)
  \item[] \textbf{B.} Một đường thẳng bất kỳ trong (P)
  \item[] \textbf{C.} Giao tuyến của (P) với mặt phẳng khác
  \item[] \textbf{D.} Hai đường thẳng cắt nhau nằm trong mặt phẳng (P)
\end{itemize}\vspace{2pt}
\textbf{\color{primary} Câu 8.} Áp dụng trực tiếp định nghĩa và quy tắc trong Chương VII: Một khối chóp có diện tích đáy B = 54 và chiều cao h = 10. Thể tích V = 1/3 $\cdot$ B$\cdot$h của khối chóp bằng:\par
\hspace*{1.5em}\textbf{A.} 270\quad\quad \textbf{B.} 180\quad\quad \textbf{C.} 90\quad\quad \textbf{D.} 540\par\vspace{2pt}
\textbf{\color{primary} Câu 9.} Kiểm tra nhanh kiến thức nền tảng: Một khối lăng trụ có diện tích đáy B = 13 và chiều cao h = 12. Thể tích V = B$\cdot$h của khối lăng trụ bằng:\par
\hspace*{1.5em}\textbf{A.} 312\quad\quad \textbf{B.} 52.0\quad\quad \textbf{C.} 156\quad\quad \textbf{D.} 25\par\vspace{2pt}
\textbf{\color{primary} Câu 10.} Cho dữ kiện nhận biết sau đây: Đường thẳng d vuông góc với mặt phẳng (P) khi và chỉ khi d vuông góc với. Phương án chính xác là:\par
\begin{itemize}\setlength{\itemsep}{1pt}\setlength{\parskip}{0pt}
  \item[] \textbf{A.} Giao tuyến của (P) với mặt phẳng khác
  \item[] \textbf{B.} Hai đường thẳng cắt nhau nằm trong mặt phẳng (P)
  \item[] \textbf{C.} Một đường thẳng bất kỳ trong (P)
  \item[] \textbf{D.} Hai đường thẳng song song nằm trong (P)
\end{itemize}\vspace{2pt}
\textbf{\color{primary} Câu 11.} Hãy chỉ ra kết quả đúng trong bốn lựa chọn A, B, C, D đối với yêu cầu: Một khối chóp có diện tích đáy B = 72 và chiều cao h = 13. Thể tích V = 1/3 $\cdot$ B$\cdot$h của khối chóp bằng:\par
\hspace*{1.5em}\textbf{A.} 312\quad\quad \textbf{B.} 156\quad\quad \textbf{C.} 936\quad\quad \textbf{D.} 468\par\vspace{2pt}
\textbf{\color{primary} Câu 12.} Câu hỏi ôn tập lý thuyết và kỹ năng cơ bản (Chương VII): Một khối lăng trụ có diện tích đáy B = 16 và chiều cao h = 15. Thể tích V = B$\cdot$h của khối lăng trụ bằng:\par
\hspace*{1.5em}\textbf{A.} 31\quad\quad \textbf{B.} 80.0\quad\quad \textbf{C.} 480\quad\quad \textbf{D.} 240\par\vspace{2pt}
\textbf{\color{primary} Câu 13.} Theo nội dung sách giáo khoa Kết nối tri thức với cuộc sống (Chương VII), đường thẳng d vuông góc với mặt phẳng (P) khi và chỉ khi d vuông góc với:\par
\begin{itemize}\setlength{\itemsep}{1pt}\setlength{\parskip}{0pt}
  \item[] \textbf{A.} Hai đường thẳng cắt nhau nằm trong mặt phẳng (P)
  \item[] \textbf{B.} Giao tuyến của (P) với mặt phẳng khác
  \item[] \textbf{C.} Một đường thẳng bất kỳ trong (P)
  \item[] \textbf{D.} Hai đường thẳng song song nằm trong (P)
\end{itemize}\vspace{2pt}
\textbf{\color{primary} Câu 14.} Tìm đáp số đúng của bài toán cơ bản sau: Một khối chóp có diện tích đáy B = 90 và chiều cao h = 16. Thể tích V = 1/3 $\cdot$ B$\cdot$h của khối chóp bằng:\par
\hspace*{1.5em}\textbf{A.} 240\quad\quad \textbf{B.} 480\quad\quad \textbf{C.} 1440\quad\quad \textbf{D.} 720\par\vspace{2pt}
\textbf{\color{primary} Câu 15.} Xét đối tượng toán học trong Chương VII: Một khối lăng trụ có diện tích đáy B = 19 và chiều cao h = 18. Thể tích V = B$\cdot$h của khối lăng trụ bằng. Giá trị hoặc tính chất tương ứng là:\par
\hspace*{1.5em}\textbf{A.} 342\quad\quad \textbf{B.} 684\quad\quad \textbf{C.} 114.0\quad\quad \textbf{D.} 37\par\vspace{2pt}
\textbf{\color{primary} Câu 16.} Hoàn thành phép xác định cơ bản sau đây: Đường thẳng d vuông góc với mặt phẳng (P) khi và chỉ khi d vuông góc với:\par
\begin{itemize}\setlength{\itemsep}{1pt}\setlength{\parskip}{0pt}
  \item[] \textbf{A.} Giao tuyến của (P) với mặt phẳng khác
  \item[] \textbf{B.} Một đường thẳng bất kỳ trong (P)
  \item[] \textbf{C.} Hai đường thẳng song song nằm trong (P)
  \item[] \textbf{D.} Hai đường thẳng cắt nhau nằm trong mặt phẳng (P)
\end{itemize}\vspace{2pt}
\textbf{\color{primary} Câu 17.} Trong các khẳng định hoặc giá trị dưới đây, đâu là đáp án đúng cho câu hỏi: Một khối chóp có diện tích đáy B = 108 và chiều cao h = 19. Thể tích V = 1/3 $\cdot$ B$\cdot$h của khối chóp bằng?\par
\hspace*{1.5em}\textbf{A.} 1026\quad\quad \textbf{B.} 684\quad\quad \textbf{C.} 2052\quad\quad \textbf{D.} 342\par\vspace{2pt}
\textbf{\color{primary} Câu 18.} Bài tập khởi động Chương VII: Một khối lăng trụ có diện tích đáy B = 22 và chiều cao h = 21. Thể tích V = B$\cdot$h của khối lăng trụ bằng:\par
\hspace*{1.5em}\textbf{A.} 154.0\quad\quad \textbf{B.} 43\quad\quad \textbf{C.} 924\quad\quad \textbf{D.} 462\par\vspace{2pt}
\textbf{\color{primary} Câu 19.} Quan sát dữ kiện đề bài: Đường thẳng d vuông góc với mặt phẳng (P) khi và chỉ khi d vuông góc với. Hãy chọn phương án đúng:\par
\begin{itemize}\setlength{\itemsep}{1pt}\setlength{\parskip}{0pt}
  \item[] \textbf{A.} Hai đường thẳng cắt nhau nằm trong mặt phẳng (P)
  \item[] \textbf{B.} Giao tuyến của (P) với mặt phẳng khác
  \item[] \textbf{C.} Hai đường thẳng song song nằm trong (P)
  \item[] \textbf{D.} Một đường thẳng bất kỳ trong (P)
\end{itemize}\vspace{2pt}
\textbf{\color{primary} Câu 20.} Sử dụng công thức hoặc tính chất đầu bài của Chương VII để giải quyết yêu cầu: Một khối chóp có diện tích đáy B = 126 và chiều cao h = 22. Thể tích V = 1/3 $\cdot$ B$\cdot$h của khối chóp bằng:\par
\hspace*{1.5em}\textbf{A.} 924\quad\quad \textbf{B.} 1386\quad\quad \textbf{C.} 2772\quad\quad \textbf{D.} 462\par\vspace{2pt}
\subsection*{\color{accent} Phần II. Mức độ Trung bình (Thông hiểu -- Câu 21 đến Câu 40)}
\textbf{\color{primary} Câu 21.} Cho hình chóp S.ABCD có đáy ABCD là hình vuông cạnh a = 2, cạnh bên SA vuông góc với đáy và SA = 2. Góc giữa đường thẳng SB và mặt phẳng đáy (ABCD) bằng:\par
\hspace*{1.5em}\textbf{A.} 45°\quad\quad \textbf{B.} 30°\quad\quad \textbf{C.} 60°\quad\quad \textbf{D.} 90°\par\vspace{2pt}
\textbf{\color{primary} Câu 22.} Vận dụng quy tắc trong Chương VII (Quan hệ vuông góc trong không gian) để giải bài toán: Cho hình chóp S.ABC có SA $\perp$ (ABC), đáy ABC là tam giác vuông tại B có AB = 6, BC = 8 và SA = 12. Thể tích khối chóp S.ABC bằng:\par
\hspace*{1.5em}\textbf{A.} 192\quad\quad \textbf{B.} 48\quad\quad \textbf{C.} 288\quad\quad \textbf{D.} 96\par\vspace{2pt}
\textbf{\color{primary} Câu 23.} Biến đổi và tính toán theo kiến thức Chương VII: Cho hình lập phương ABCD.A'B'C'D' có cạnh bằng 5. Khoảng cách giữa hai mặt phẳng song song (ABCD) và (A'B'C'D') bằng. Kết quả thu được là:\par
\hspace*{1.5em}\textbf{A.} 5$\sqrt{2}$\quad\quad \textbf{B.} 5\quad\quad \textbf{C.} 2.5\quad\quad \textbf{D.} 5$\sqrt{3}$\par\vspace{2pt}
\textbf{\color{primary} Câu 24.} Xét bài toán mức độ thông hiểu sau: Cho hình chóp S.ABCD có đáy ABCD là hình vuông cạnh a = 5, cạnh bên SA vuông góc với đáy và SA = 5. Góc giữa đường thẳng SB và mặt phẳng đáy (ABCD) bằng. Đáp án đúng là:\par
\hspace*{1.5em}\textbf{A.} 45°\quad\quad \textbf{B.} 30°\quad\quad \textbf{C.} 90°\quad\quad \textbf{D.} 60°\par\vspace{2pt}
\textbf{\color{primary} Câu 25.} Để hoàn thành phiếu học tập Chương VII, hãy xác định: Cho hình chóp S.ABC có SA $\perp$ (ABC), đáy ABC là tam giác vuông tại B có AB = 15, BC = 20 và SA = 30. Thể tích khối chóp S.ABC bằng:\par
\hspace*{1.5em}\textbf{A.} 750\quad\quad \textbf{B.} 1500\quad\quad \textbf{C.} 4500\quad\quad \textbf{D.} 3000\par\vspace{2pt}
\textbf{\color{primary} Câu 26.} Phân tích dữ kiện bài toán: Cho hình lập phương ABCD.A'B'C'D' có cạnh bằng 8. Khoảng cách giữa hai mặt phẳng song song (ABCD) và (A'B'C'D') bằng. Giá trị cần tìm bằng:\par
\hspace*{1.5em}\textbf{A.} 8\quad\quad \textbf{B.} 4.0\quad\quad \textbf{C.} 8$\sqrt{3}$\quad\quad \textbf{D.} 8$\sqrt{2}$\par\vspace{2pt}
\textbf{\color{primary} Câu 27.} Áp dụng tính chất đã học ở Chương VII vào bài tập sau: Cho hình chóp S.ABCD có đáy ABCD là hình vuông cạnh a = 8, cạnh bên SA vuông góc với đáy và SA = 8. Góc giữa đường thẳng SB và mặt phẳng đáy (ABCD) bằng:\par
\hspace*{1.5em}\textbf{A.} 60°\quad\quad \textbf{B.} 90°\quad\quad \textbf{C.} 45°\quad\quad \textbf{D.} 30°\par\vspace{2pt}
\textbf{\color{primary} Câu 28.} Thực hiện các bước suy luận hợp lý để giải quyết yêu cầu: Cho hình chóp S.ABC có SA $\perp$ (ABC), đáy ABC là tam giác vuông tại B có AB = 24, BC = 32 và SA = 48. Thể tích khối chóp S.ABC bằng:\par
\hspace*{1.5em}\textbf{A.} 18432\quad\quad \textbf{B.} 6144\quad\quad \textbf{C.} 12288\quad\quad \textbf{D.} 3072\par\vspace{2pt}
\textbf{\color{primary} Câu 29.} Cho bài toán tự luyện (Quan hệ vuông góc trong không gian): Cho hình lập phương ABCD.A'B'C'D' có cạnh bằng 11. Khoảng cách giữa hai mặt phẳng song song (ABCD) và (A'B'C'D') bằng. Phương án nào sau đây chính xác?\par
\hspace*{1.5em}\textbf{A.} 11\quad\quad \textbf{B.} 5.5\quad\quad \textbf{C.} 11$\sqrt{3}$\quad\quad \textbf{D.} 11$\sqrt{2}$\par\vspace{2pt}
\textbf{\color{primary} Câu 30.} Kết hợp các phép biến đổi trong Chương VII, hãy tìm kết quả của bài toán: Cho hình chóp S.ABCD có đáy ABCD là hình vuông cạnh a = 11, cạnh bên SA vuông góc với đáy và SA = 11. Góc giữa đường thẳng SB và mặt phẳng đáy (ABCD) bằng:\par
\hspace*{1.5em}\textbf{A.} 60°\quad\quad \textbf{B.} 90°\quad\quad \textbf{C.} 30°\quad\quad \textbf{D.} 45°\par\vspace{2pt}
\textbf{\color{primary} Câu 31.} Một bài kiểm tra định kỳ nội dung Chương VII nêu câu hỏi: Cho hình chóp S.ABC có SA $\perp$ (ABC), đáy ABC là tam giác vuông tại B có AB = 33, BC = 44 và SA = 66. Thể tích khối chóp S.ABC bằng. Hãy chọn đáp số đúng:\par
\hspace*{1.5em}\textbf{A.} 47916\quad\quad \textbf{B.} 31944\quad\quad \textbf{C.} 7986\quad\quad \textbf{D.} 15972\par\vspace{2pt}
\textbf{\color{primary} Câu 32.} Tìm lời giải đúng cho bài toán thông hiểu: Cho hình lập phương ABCD.A'B'C'D' có cạnh bằng 14. Khoảng cách giữa hai mặt phẳng song song (ABCD) và (A'B'C'D') bằng:\par
\hspace*{1.5em}\textbf{A.} 14\quad\quad \textbf{B.} 7.0\quad\quad \textbf{C.} 14$\sqrt{2}$\quad\quad \textbf{D.} 14$\sqrt{3}$\par\vspace{2pt}
\textbf{\color{primary} Câu 33.} Sử dụng mối liên hệ giữa các đại lượng trong Chương VII: Cho hình chóp S.ABCD có đáy ABCD là hình vuông cạnh a = 14, cạnh bên SA vuông góc với đáy và SA = 14. Góc giữa đường thẳng SB và mặt phẳng đáy (ABCD) bằng:\par
\hspace*{1.5em}\textbf{A.} 30°\quad\quad \textbf{B.} 90°\quad\quad \textbf{C.} 60°\quad\quad \textbf{D.} 45°\par\vspace{2pt}
\textbf{\color{primary} Câu 34.} Hãy tính toán và chọn phương án đúng đối với bài tập: Cho hình chóp S.ABC có SA $\perp$ (ABC), đáy ABC là tam giác vuông tại B có AB = 42, BC = 56 và SA = 84. Thể tích khối chóp S.ABC bằng:\par
\hspace*{1.5em}\textbf{A.} 32928\quad\quad \textbf{B.} 16464\quad\quad \textbf{C.} 98784\quad\quad \textbf{D.} 65856\par\vspace{2pt}
\textbf{\color{primary} Câu 35.} Dựa trên các tính chất trọng tâm của Chương VII, xét bài toán: Cho hình lập phương ABCD.A'B'C'D' có cạnh bằng 17. Khoảng cách giữa hai mặt phẳng song song (ABCD) và (A'B'C'D') bằng:\par
\hspace*{1.5em}\textbf{A.} 17\quad\quad \textbf{B.} 17$\sqrt{3}$\quad\quad \textbf{C.} 8.5\quad\quad \textbf{D.} 17$\sqrt{2}$\par\vspace{2pt}
\textbf{\color{primary} Câu 36.} Thực hành giải toán (Quan hệ vuông góc trong không gian): Cho hình chóp S.ABCD có đáy ABCD là hình vuông cạnh a = 17, cạnh bên SA vuông góc với đáy và SA = 17. Góc giữa đường thẳng SB và mặt phẳng đáy (ABCD) bằng. Đáp án nào dưới đây thỏa mãn?\par
\hspace*{1.5em}\textbf{A.} 30°\quad\quad \textbf{B.} 45°\quad\quad \textbf{C.} 90°\quad\quad \textbf{D.} 60°\par\vspace{2pt}
\textbf{\color{primary} Câu 37.} Kiểm tra kỹ năng biến đổi và suy luận qua bài toán: Cho hình chóp S.ABC có SA $\perp$ (ABC), đáy ABC là tam giác vuông tại B có AB = 51, BC = 68 và SA = 102. Thể tích khối chóp S.ABC bằng:\par
\hspace*{1.5em}\textbf{A.} 58956\quad\quad \textbf{B.} 117912\quad\quad \textbf{C.} 176868\quad\quad \textbf{D.} 29478\par\vspace{2pt}
\textbf{\color{primary} Câu 38.} Cho giả thiết thuộc phạm vi Chương VII: Cho hình lập phương ABCD.A'B'C'D' có cạnh bằng 20. Khoảng cách giữa hai mặt phẳng song song (ABCD) và (A'B'C'D') bằng. Kết luận đúng là:\par
\hspace*{1.5em}\textbf{A.} 20\quad\quad \textbf{B.} 10.0\quad\quad \textbf{C.} 20$\sqrt{3}$\quad\quad \textbf{D.} 20$\sqrt{2}$\par\vspace{2pt}
\textbf{\color{primary} Câu 39.} Rút gọn hoặc tính giá trị theo yêu cầu sau của Chương VII: Cho hình chóp S.ABCD có đáy ABCD là hình vuông cạnh a = 20, cạnh bên SA vuông góc với đáy và SA = 20. Góc giữa đường thẳng SB và mặt phẳng đáy (ABCD) bằng:\par
\hspace*{1.5em}\textbf{A.} 90°\quad\quad \textbf{B.} 45°\quad\quad \textbf{C.} 60°\quad\quad \textbf{D.} 30°\par\vspace{2pt}
\textbf{\color{primary} Câu 40.} Xác định phương án đúng sau khi thực hiện phép tính / suy luận: Cho hình chóp S.ABC có SA $\perp$ (ABC), đáy ABC là tam giác vuông tại B có AB = 60, BC = 80 và SA = 120. Thể tích khối chóp S.ABC bằng:\par
\hspace*{1.5em}\textbf{A.} 288000\quad\quad \textbf{B.} 192000\quad\quad \textbf{C.} 48000\quad\quad \textbf{D.} 96000\par\vspace{2pt}
\subsection*{\color{accent} Phần III. Mức độ Khó (Vận dụng \& Vận dụng cao -- Câu 41 đến Câu 50)}
\textbf{\color{primary} Câu 41.} Cho hình chóp S.ABCD có đáy ABCD là hình vuông cạnh a = 3, SA $\perp$ (ABCD) và SA = 3. Khoảng cách từ điểm A đến mặt phẳng (SBC) bằng:\par
\hspace*{1.5em}\textbf{A.} 3$\sqrt{3}$ / 2\quad\quad \textbf{B.} 3\quad\quad \textbf{C.} 3/2\quad\quad \textbf{D.} 3$\sqrt{2}$ / 2\par\vspace{2pt}
\textbf{\color{primary} Câu 42.} Bài toán vận dụng (Chương VII – Quan hệ vuông góc trong không gian): Cho hình chóp đều S.ABCD có tất cả các cạnh đáy và cạnh bên đều bằng a. Giá trị cosin của góc giữa cạnh bên SA và mặt phẳng đáy (ABCD) bằng. Kết quả đúng là:\par
\hspace*{1.5em}\textbf{A.} $\sqrt{3}$ / 2\quad\quad \textbf{B.} 1/2\quad\quad \textbf{C.} $\sqrt{2}$ / 2 (tức góc 45°)\quad\quad \textbf{D.} $\sqrt{3}$ / 3\par\vspace{2pt}
\textbf{\color{primary} Câu 43.} Vận dụng kiến thức nâng cao trong phạm vi Chương VII để giải bài toán: Cho hình chóp S.ABCD có đáy ABCD là hình vuông cạnh a = 8, cạnh bên SA $\perp$ (ABCD) và SA = 8. Khoảng cách từ điểm A đến mặt phẳng (SBC) bằng:\par
\hspace*{1.5em}\textbf{A.} 8\quad\quad \textbf{B.} 4$\sqrt{3}$\quad\quad \textbf{C.} 4\quad\quad \textbf{D.} 4$\sqrt{2}$\par\vspace{2pt}
\textbf{\color{primary} Câu 44.} Phân tích và giải bài toán tổng hợp sau của Chương VII: Cho hình chóp S.ABCD có đáy ABCD là hình vuông cạnh a = 3, SA $\perp$ (ABCD) và SA = 3. Khoảng cách từ điểm A đến mặt phẳng (SBC) bằng:\par
\hspace*{1.5em}\textbf{A.} 3\quad\quad \textbf{B.} 3$\sqrt{2}$ / 2\quad\quad \textbf{C.} 3/2\quad\quad \textbf{D.} 3$\sqrt{3}$ / 2\par\vspace{2pt}
\textbf{\color{primary} Câu 45.} Thử thách tư duy trong phạm vi kiến thức Chương VII: Cho hình chóp đều S.ABCD có tất cả các cạnh đáy và cạnh bên đều bằng a. Giá trị cosin của góc giữa cạnh bên SA và mặt phẳng đáy (ABCD) bằng. Đáp số chính xác là:\par
\hspace*{1.5em}\textbf{A.} 1/2\quad\quad \textbf{B.} $\sqrt{2}$ / 2 (tức góc 45°)\quad\quad \textbf{C.} $\sqrt{3}$ / 3\quad\quad \textbf{D.} $\sqrt{3}$ / 2\par\vspace{2pt}
\textbf{\color{primary} Câu 46.} Sử dụng phương pháp giải toán của Chương VII để tìm đáp án cho bài toán: Cho hình chóp S.ABCD có đáy ABCD là hình vuông cạnh a = 14, cạnh bên SA $\perp$ (ABCD) và SA = 14. Khoảng cách từ điểm A đến mặt phẳng (SBC) bằng:\par
\hspace*{1.5em}\textbf{A.} 7$\sqrt{3}$\quad\quad \textbf{B.} 14\quad\quad \textbf{C.} 7$\sqrt{2}$\quad\quad \textbf{D.} 7\par\vspace{2pt}
\textbf{\color{primary} Câu 47.} Xét bài toán vận dụng thực tiễn / nâng cao thuộc Chương VII: Cho hình chóp S.ABCD có đáy ABCD là hình vuông cạnh a = 3, SA $\perp$ (ABCD) và SA = 3. Khoảng cách từ điểm A đến mặt phẳng (SBC) bằng:\par
\hspace*{1.5em}\textbf{A.} 3/2\quad\quad \textbf{B.} 3\quad\quad \textbf{C.} 3$\sqrt{3}$ / 2\quad\quad \textbf{D.} 3$\sqrt{2}$ / 2\par\vspace{2pt}
\textbf{\color{primary} Câu 48.} Tìm kết quả cuối cùng của bài toán phân loại học sinh khá giỏi (Chương VII): Cho hình chóp đều S.ABCD có tất cả các cạnh đáy và cạnh bên đều bằng a. Giá trị cosin của góc giữa cạnh bên SA và mặt phẳng đáy (ABCD) bằng:\par
\hspace*{1.5em}\textbf{A.} $\sqrt{2}$ / 2 (tức góc 45°)\quad\quad \textbf{B.} 1/2\quad\quad \textbf{C.} $\sqrt{3}$ / 3\quad\quad \textbf{D.} $\sqrt{3}$ / 2\par\vspace{2pt}
\textbf{\color{primary} Câu 49.} Kết hợp các tính chất đã học trong Chương VII, hãy giải quyết bài toán: Cho hình chóp S.ABCD có đáy ABCD là hình vuông cạnh a = 20, cạnh bên SA $\perp$ (ABCD) và SA = 20. Khoảng cách từ điểm A đến mặt phẳng (SBC) bằng:\par
\hspace*{1.5em}\textbf{A.} 10$\sqrt{2}$\quad\quad \textbf{B.} 10$\sqrt{3}$\quad\quad \textbf{C.} 20\quad\quad \textbf{D.} 10\par\vspace{2pt}
\textbf{\color{primary} Câu 50.} Câu hỏi vận dụng cao cuối Chương VII: Cho hình chóp S.ABCD có đáy ABCD là hình vuông cạnh a = 3, SA $\perp$ (ABCD) và SA = 3. Khoảng cách từ điểm A đến mặt phẳng (SBC) bằng. Phương án đúng là:\par
\hspace*{1.5em}\textbf{A.} 3$\sqrt{3}$ / 2\quad\quad \textbf{B.} 3/2\quad\quad \textbf{C.} 3$\sqrt{2}$ / 2\quad\quad \textbf{D.} 3\par\vspace{2pt}
\subsection*{\color{primary} Bảng đáp án Chương VII: Quan hệ vuông góc trong không gian}
\begin{center}
\renewcommand{\arraystretch}{1.25}
\begin{tabularx}{\textwidth}{*{10}{>{\centering\arraybackslash}X}}
\toprule
\textbf{\color{primary} Cột 1} & \textbf{\color{primary} Cột 2} & \textbf{\color{primary} Cột 3} & \textbf{\color{primary} Cột 4} & \textbf{\color{primary} Cột 5} & \textbf{\color{primary} Cột 6} & \textbf{\color{primary} Cột 7} & \textbf{\color{primary} Cột 8} & \textbf{\color{primary} Cột 9} & \textbf{\color{primary} Cột 10} \\
\midrule
\textbf{1.C} & \textbf{2.D} & \textbf{3.A} & \textbf{4.B} & \textbf{5.C} & \textbf{6.D} & \textbf{7.D} & \textbf{8.B} & \textbf{9.C} & \textbf{10.B} \\
\rowcolor{rowtint} \textbf{11.A} & \textbf{12.D} & \textbf{13.A} & \textbf{14.B} & \textbf{15.A} & \textbf{16.D} & \textbf{17.B} & \textbf{18.D} & \textbf{19.A} & \textbf{20.A} \\
\textbf{21.A} & \textbf{22.D} & \textbf{23.B} & \textbf{24.A} & \textbf{25.B} & \textbf{26.A} & \textbf{27.C} & \textbf{28.B} & \textbf{29.A} & \textbf{30.D} \\
\rowcolor{rowtint} \textbf{31.D} & \textbf{32.A} & \textbf{33.D} & \textbf{34.A} & \textbf{35.A} & \textbf{36.B} & \textbf{37.A} & \textbf{38.A} & \textbf{39.B} & \textbf{40.D} \\
\textbf{41.D} & \textbf{42.C} & \textbf{43.D} & \textbf{44.B} & \textbf{45.B} & \textbf{46.C} & \textbf{47.D} & \textbf{48.A} & \textbf{49.A} & \textbf{50.C} \\
\bottomrule
\end{tabularx}
\end{center}\vspace{10pt}
\section*{\color{primary} CHƯƠNG VIII: CÁC QUY TẮC TÍNH XÁC SUẤT}
\textit{Trọng tâm kiến thức SGK: Biến cố hợp, biến cố giao, biến cố xung khắc, quy tắc cộng xác suất, biến cố độc lập và quy tắc nhân xác suất.}\par\vspace{4pt}
\subsection*{\color{accent} Phần I. Mức độ Dễ (Nhận biết -- Câu 1 đến Câu 20)}
\textbf{\color{primary} Câu 1.} Nếu hai biến cố A và B xung khắc (A $\cap$ B = $\emptyset$) thì quy tắc cộng xác suất được viết là:\par
\hspace*{1.5em}\textbf{A.} P(A $\cup$ B) = P(A) $\cdot$ P(B)\quad\quad \textbf{B.} P(A $\cup$ B) = P(A) + P(B)\quad\quad \textbf{C.} P(A $\cup$ B) = P(A) - P(B)\quad\quad \textbf{D.} P(A $\cap$ B) = P(A) + P(B)\par\vspace{2pt}
\textbf{\color{primary} Câu 2.} Thực hiện yêu cầu nhận biết sau: Nếu hai biến cố A và B độc lập với nhau thì xác suất của biến cố giao P(AB) bằng.\par
\hspace*{1.5em}\textbf{A.} P(AB) = P(A) + P(B)\quad\quad \textbf{B.} P(AB) = P(A) / P(B)\quad\quad \textbf{C.} P(AB) = 0\quad\quad \textbf{D.} P(AB) = P(A) $\cdot$ P(B)\par\vspace{2pt}
\textbf{\color{primary} Câu 3.} Trong bài học thuộc Chương VIII, xét nội dung: Cho hai biến cố A, B bất kỳ. Công thức cộng xác suất tổng quát là. Đáp án đúng là:\par
\begin{itemize}\setlength{\itemsep}{1pt}\setlength{\parskip}{0pt}
  \item[] \textbf{A.} P(A $\cup$ B) = 1 - P(AB)
  \item[] \textbf{B.} P(A $\cup$ B) = P(A) $\cdot$ P(B)
  \item[] \textbf{C.} P(A $\cup$ B) = P(A) + P(B) - P(AB)
  \item[] \textbf{D.} P(A $\cup$ B) = P(A) + P(B) + P(AB)
\end{itemize}\vspace{2pt}
\textbf{\color{primary} Câu 4.} Dựa vào kiến thức chuẩn SGK của Chương VIII (Các quy tắc tính xác suất), cho hai biến cố xung khắc A và B có P(A) = 4/19 và P(B) = 6/19. Theo quy tắc cộng xác suất, P(A $\cup$ B) bằng:\par
\hspace*{1.5em}\textbf{A.} 7/19\quad\quad \textbf{B.} 24/361\quad\quad \textbf{C.} 2/19\quad\quad \textbf{D.} 10/19\par\vspace{2pt}
\textbf{\color{primary} Câu 5.} Chọn kết quả đúng nhất khi xét: Nếu hai biến cố A và B xung khắc (A $\cap$ B = $\emptyset$) thì quy tắc cộng xác suất được viết là:\par
\hspace*{1.5em}\textbf{A.} P(A $\cup$ B) = P(A) + P(B)\quad\quad \textbf{B.} P(A $\cup$ B) = P(A) $\cdot$ P(B)\quad\quad \textbf{C.} P(A $\cap$ B) = P(A) + P(B)\quad\quad \textbf{D.} P(A $\cup$ B) = P(A) - P(B)\par\vspace{2pt}
\textbf{\color{primary} Câu 6.} Khi làm bài tập cơ bản về các quy tắc tính xác suất, một học sinh gặp câu hỏi: Nếu hai biến cố A và B độc lập với nhau thì xác suất của biến cố giao P(AB) bằng. Kết quả nào dưới đây đúng?\par
\hspace*{1.5em}\textbf{A.} P(AB) = P(A) + P(B)\quad\quad \textbf{B.} P(AB) = 0\quad\quad \textbf{C.} P(AB) = P(A) / P(B)\quad\quad \textbf{D.} P(AB) = P(A) $\cdot$ P(B)\par\vspace{2pt}
\textbf{\color{primary} Câu 7.} Xác định giá trị hoặc khẳng định đúng cho bài toán sau: Cho hai biến cố A, B bất kỳ. Công thức cộng xác suất tổng quát là:\par
\begin{itemize}\setlength{\itemsep}{1pt}\setlength{\parskip}{0pt}
  \item[] \textbf{A.} P(A $\cup$ B) = P(A) + P(B) - P(AB)
  \item[] \textbf{B.} P(A $\cup$ B) = P(A) + P(B) + P(AB)
  \item[] \textbf{C.} P(A $\cup$ B) = 1 - P(AB)
  \item[] \textbf{D.} P(A $\cup$ B) = P(A) $\cdot$ P(B)
\end{itemize}\vspace{2pt}
\textbf{\color{primary} Câu 8.} Áp dụng trực tiếp định nghĩa và quy tắc trong Chương VIII: Cho hai biến cố xung khắc A và B có P(A) = 8/31 và P(B) = 10/31. Theo quy tắc cộng xác suất, P(A $\cup$ B) bằng:\par
\hspace*{1.5em}\textbf{A.} 2/31\quad\quad \textbf{B.} 11/31\quad\quad \textbf{C.} 18/31\quad\quad \textbf{D.} 80/961\par\vspace{2pt}
\textbf{\color{primary} Câu 9.} Kiểm tra nhanh kiến thức nền tảng: Nếu hai biến cố A và B xung khắc (A $\cap$ B = $\emptyset$) thì quy tắc cộng xác suất được viết là:\par
\hspace*{1.5em}\textbf{A.} P(A $\cup$ B) = P(A) $\cdot$ P(B)\quad\quad \textbf{B.} P(A $\cap$ B) = P(A) + P(B)\quad\quad \textbf{C.} P(A $\cup$ B) = P(A) + P(B)\quad\quad \textbf{D.} P(A $\cup$ B) = P(A) - P(B)\par\vspace{2pt}
\textbf{\color{primary} Câu 10.} Cho dữ kiện nhận biết sau đây: Nếu hai biến cố A và B độc lập với nhau thì xác suất của biến cố giao P(AB) bằng. Phương án chính xác là:\par
\hspace*{1.5em}\textbf{A.} P(AB) = P(A) $\cdot$ P(B)\quad\quad \textbf{B.} P(AB) = P(A) / P(B)\quad\quad \textbf{C.} P(AB) = 0\quad\quad \textbf{D.} P(AB) = P(A) + P(B)\par\vspace{2pt}
\textbf{\color{primary} Câu 11.} Hãy chỉ ra kết quả đúng trong bốn lựa chọn A, B, C, D đối với yêu cầu: Cho hai biến cố A, B bất kỳ. Công thức cộng xác suất tổng quát là:\par
\begin{itemize}\setlength{\itemsep}{1pt}\setlength{\parskip}{0pt}
  \item[] \textbf{A.} P(A $\cup$ B) = P(A) + P(B) - P(AB)
  \item[] \textbf{B.} P(A $\cup$ B) = 1 - P(AB)
  \item[] \textbf{C.} P(A $\cup$ B) = P(A) + P(B) + P(AB)
  \item[] \textbf{D.} P(A $\cup$ B) = P(A) $\cdot$ P(B)
\end{itemize}\vspace{2pt}
\textbf{\color{primary} Câu 12.} Câu hỏi ôn tập lý thuyết và kỹ năng cơ bản (Chương VIII): Cho hai biến cố xung khắc A và B có P(A) = 12/43 và P(B) = 14/43. Theo quy tắc cộng xác suất, P(A $\cup$ B) bằng:\par
\hspace*{1.5em}\textbf{A.} 26/43\quad\quad \textbf{B.} 2/43\quad\quad \textbf{C.} 15/43\quad\quad \textbf{D.} 168/1849\par\vspace{2pt}
\textbf{\color{primary} Câu 13.} Theo nội dung sách giáo khoa Kết nối tri thức với cuộc sống (Chương VIII), nếu hai biến cố A và B xung khắc (A $\cap$ B = $\emptyset$) thì quy tắc cộng xác suất được viết là:\par
\hspace*{1.5em}\textbf{A.} P(A $\cup$ B) = P(A) + P(B)\quad\quad \textbf{B.} P(A $\cap$ B) = P(A) + P(B)\quad\quad \textbf{C.} P(A $\cup$ B) = P(A) $\cdot$ P(B)\quad\quad \textbf{D.} P(A $\cup$ B) = P(A) - P(B)\par\vspace{2pt}
\textbf{\color{primary} Câu 14.} Tìm đáp số đúng của bài toán cơ bản sau: Nếu hai biến cố A và B độc lập với nhau thì xác suất của biến cố giao P(AB) bằng:\par
\hspace*{1.5em}\textbf{A.} P(AB) = P(A) $\cdot$ P(B)\quad\quad \textbf{B.} P(AB) = P(A) / P(B)\quad\quad \textbf{C.} P(AB) = 0\quad\quad \textbf{D.} P(AB) = P(A) + P(B)\par\vspace{2pt}
\textbf{\color{primary} Câu 15.} Xét đối tượng toán học trong Chương VIII: Cho hai biến cố A, B bất kỳ. Công thức cộng xác suất tổng quát là. Giá trị hoặc tính chất tương ứng là:\par
\begin{itemize}\setlength{\itemsep}{1pt}\setlength{\parskip}{0pt}
  \item[] \textbf{A.} P(A $\cup$ B) = 1 - P(AB)
  \item[] \textbf{B.} P(A $\cup$ B) = P(A) $\cdot$ P(B)
  \item[] \textbf{C.} P(A $\cup$ B) = P(A) + P(B) + P(AB)
  \item[] \textbf{D.} P(A $\cup$ B) = P(A) + P(B) - P(AB)
\end{itemize}\vspace{2pt}
\textbf{\color{primary} Câu 16.} Hoàn thành phép xác định cơ bản sau đây: Cho hai biến cố xung khắc A và B có P(A) = 16/55 và P(B) = 18/55. Theo quy tắc cộng xác suất, P(A $\cup$ B) bằng:\par
\hspace*{1.5em}\textbf{A.} 2/55\quad\quad \textbf{B.} 19/55\quad\quad \textbf{C.} 34/55\quad\quad \textbf{D.} 288/3025\par\vspace{2pt}
\textbf{\color{primary} Câu 17.} Trong các khẳng định hoặc giá trị dưới đây, đâu là đáp án đúng cho câu hỏi: Nếu hai biến cố A và B xung khắc (A $\cap$ B = $\emptyset$) thì quy tắc cộng xác suất được viết là?\par
\hspace*{1.5em}\textbf{A.} P(A $\cap$ B) = P(A) + P(B)\quad\quad \textbf{B.} P(A $\cup$ B) = P(A) - P(B)\quad\quad \textbf{C.} P(A $\cup$ B) = P(A) $\cdot$ P(B)\quad\quad \textbf{D.} P(A $\cup$ B) = P(A) + P(B)\par\vspace{2pt}
\textbf{\color{primary} Câu 18.} Bài tập khởi động Chương VIII: Nếu hai biến cố A và B độc lập với nhau thì xác suất của biến cố giao P(AB) bằng:\par
\hspace*{1.5em}\textbf{A.} P(AB) = P(A) / P(B)\quad\quad \textbf{B.} P(AB) = 0\quad\quad \textbf{C.} P(AB) = P(A) $\cdot$ P(B)\quad\quad \textbf{D.} P(AB) = P(A) + P(B)\par\vspace{2pt}
\textbf{\color{primary} Câu 19.} Quan sát dữ kiện đề bài: Cho hai biến cố A, B bất kỳ. Công thức cộng xác suất tổng quát là. Hãy chọn phương án đúng:\par
\begin{itemize}\setlength{\itemsep}{1pt}\setlength{\parskip}{0pt}
  \item[] \textbf{A.} P(A $\cup$ B) = P(A) $\cdot$ P(B)
  \item[] \textbf{B.} P(A $\cup$ B) = P(A) + P(B) - P(AB)
  \item[] \textbf{C.} P(A $\cup$ B) = 1 - P(AB)
  \item[] \textbf{D.} P(A $\cup$ B) = P(A) + P(B) + P(AB)
\end{itemize}\vspace{2pt}
\textbf{\color{primary} Câu 20.} Sử dụng công thức hoặc tính chất đầu bài của Chương VIII để giải quyết yêu cầu: Cho hai biến cố xung khắc A và B có P(A) = 20/67 và P(B) = 22/67. Theo quy tắc cộng xác suất, P(A $\cup$ B) bằng:\par
\hspace*{1.5em}\textbf{A.} 2/67\quad\quad \textbf{B.} 23/67\quad\quad \textbf{C.} 440/4489\quad\quad \textbf{D.} 42/67\par\vspace{2pt}
\subsection*{\color{accent} Phần II. Mức độ Trung bình (Thông hiểu -- Câu 21 đến Câu 40)}
\textbf{\color{primary} Câu 21.} Cho hai biến cố độc lập A và B có P(A) = 0,4 và P(B) = 0,5. Xác suất của biến cố giao P(AB) và biến cố hợp P(A $\cup$ B) lần lượt bằng:\par
\begin{itemize}\setlength{\itemsep}{1pt}\setlength{\parskip}{0pt}
  \item[] \textbf{A.} P(AB) = 0,20 và P(A $\cup$ B) = 0,90
  \item[] \textbf{B.} P(AB) = 0,90 và P(A $\cup$ B) = 0,20
  \item[] \textbf{C.} P(AB) = 0,20 và P(A $\cup$ B) = 0,70
  \item[] \textbf{D.} P(AB) = 0,10 và P(A $\cup$ B) = 0,80
\end{itemize}\vspace{2pt}
\textbf{\color{primary} Câu 22.} Vận dụng quy tắc trong Chương VIII (Các quy tắc tính xác suất) để giải bài toán: Hai xạ thủ độc lập cùng bắn vào một bia. Xác suất bắn trúng của người thứ nhất là 0,8 và của người thứ hai là 0,7. Xác suất để cả hai người đều bắn trượt bằng:\par
\hspace*{1.5em}\textbf{A.} 0,06\quad\quad \textbf{B.} 0,15\quad\quad \textbf{C.} 0,94\quad\quad \textbf{D.} 0,56\par\vspace{2pt}
\textbf{\color{primary} Câu 23.} Biến đổi và tính toán theo kiến thức Chương VIII: Cho hai biến cố A, B có P(A) = 0,5; P(B) = 0,6 và P(A $\cup$ B) = 0,8. Giá trị của P(AB) bằng. Kết quả thu được là:\par
\hspace*{1.5em}\textbf{A.} 0,1\quad\quad \textbf{B.} 0,4\quad\quad \textbf{C.} 0,2\quad\quad \textbf{D.} 0,3\par\vspace{2pt}
\textbf{\color{primary} Câu 24.} Xét bài toán mức độ thông hiểu sau: Cho hai biến cố độc lập A và B có P(A) = 1/5 và P(B) = 1/3. Xác suất của biến cố giao P(AB) bằng. Đáp án đúng là:\par
\hspace*{1.5em}\textbf{A.} 1/15\quad\quad \textbf{B.} 1/8\quad\quad \textbf{C.} 2/15\quad\quad \textbf{D.} 8/15\par\vspace{2pt}
\textbf{\color{primary} Câu 25.} Để hoàn thành phiếu học tập Chương VIII, hãy xác định: Cho hai biến cố độc lập A và B có P(A) = 0,4 và P(B) = 0,5. Xác suất của biến cố giao P(AB) và biến cố hợp P(A $\cup$ B) lần lượt bằng:\par
\begin{itemize}\setlength{\itemsep}{1pt}\setlength{\parskip}{0pt}
  \item[] \textbf{A.} P(AB) = 0,20 và P(A $\cup$ B) = 0,90
  \item[] \textbf{B.} P(AB) = 0,10 và P(A $\cup$ B) = 0,80
  \item[] \textbf{C.} P(AB) = 0,20 và P(A $\cup$ B) = 0,70
  \item[] \textbf{D.} P(AB) = 0,90 và P(A $\cup$ B) = 0,20
\end{itemize}\vspace{2pt}
\textbf{\color{primary} Câu 26.} Phân tích dữ kiện bài toán: Hai xạ thủ độc lập cùng bắn vào một bia. Xác suất bắn trúng của người thứ nhất là 0,8 và của người thứ hai là 0,7. Xác suất để cả hai người đều bắn trượt bằng. Giá trị cần tìm bằng:\par
\hspace*{1.5em}\textbf{A.} 0,94\quad\quad \textbf{B.} 0,15\quad\quad \textbf{C.} 0,56\quad\quad \textbf{D.} 0,06\par\vspace{2pt}
\textbf{\color{primary} Câu 27.} Áp dụng tính chất đã học ở Chương VIII vào bài tập sau: Cho hai biến cố A, B có P(A) = 0,5; P(B) = 0,6 và P(A $\cup$ B) = 0,8. Giá trị của P(AB) bằng:\par
\hspace*{1.5em}\textbf{A.} 0,4\quad\quad \textbf{B.} 0,1\quad\quad \textbf{C.} 0,2\quad\quad \textbf{D.} 0,3\par\vspace{2pt}
\textbf{\color{primary} Câu 28.} Thực hiện các bước suy luận hợp lý để giải quyết yêu cầu: Cho hai biến cố độc lập A và B có P(A) = 1/9 và P(B) = 1/3. Xác suất của biến cố giao P(AB) bằng:\par
\hspace*{1.5em}\textbf{A.} 1/27\quad\quad \textbf{B.} 12/27\quad\quad \textbf{C.} 2/27\quad\quad \textbf{D.} 1/12\par\vspace{2pt}
\textbf{\color{primary} Câu 29.} Cho bài toán tự luyện (Các quy tắc tính xác suất): Cho hai biến cố độc lập A và B có P(A) = 0,4 và P(B) = 0,5. Xác suất của biến cố giao P(AB) và biến cố hợp P(A $\cup$ B) lần lượt bằng. Phương án nào sau đây chính xác?\par
\begin{itemize}\setlength{\itemsep}{1pt}\setlength{\parskip}{0pt}
  \item[] \textbf{A.} P(AB) = 0,10 và P(A $\cup$ B) = 0,80
  \item[] \textbf{B.} P(AB) = 0,20 và P(A $\cup$ B) = 0,70
  \item[] \textbf{C.} P(AB) = 0,20 và P(A $\cup$ B) = 0,90
  \item[] \textbf{D.} P(AB) = 0,90 và P(A $\cup$ B) = 0,20
\end{itemize}\vspace{2pt}
\textbf{\color{primary} Câu 30.} Kết hợp các phép biến đổi trong Chương VIII, hãy tìm kết quả của bài toán: Hai xạ thủ độc lập cùng bắn vào một bia. Xác suất bắn trúng của người thứ nhất là 0,8 và của người thứ hai là 0,7. Xác suất để cả hai người đều bắn trượt bằng:\par
\hspace*{1.5em}\textbf{A.} 0,06\quad\quad \textbf{B.} 0,56\quad\quad \textbf{C.} 0,94\quad\quad \textbf{D.} 0,15\par\vspace{2pt}
\textbf{\color{primary} Câu 31.} Một bài kiểm tra định kỳ nội dung Chương VIII nêu câu hỏi: Cho hai biến cố A, B có P(A) = 0,5; P(B) = 0,6 và P(A $\cup$ B) = 0,8. Giá trị của P(AB) bằng. Hãy chọn đáp số đúng:\par
\hspace*{1.5em}\textbf{A.} 0,3\quad\quad \textbf{B.} 0,4\quad\quad \textbf{C.} 0,2\quad\quad \textbf{D.} 0,1\par\vspace{2pt}
\textbf{\color{primary} Câu 32.} Tìm lời giải đúng cho bài toán thông hiểu: Cho hai biến cố độc lập A và B có P(A) = 1/13 và P(B) = 1/3. Xác suất của biến cố giao P(AB) bằng:\par
\hspace*{1.5em}\textbf{A.} 1/16\quad\quad \textbf{B.} 16/39\quad\quad \textbf{C.} 2/39\quad\quad \textbf{D.} 1/39\par\vspace{2pt}
\textbf{\color{primary} Câu 33.} Sử dụng mối liên hệ giữa các đại lượng trong Chương VIII: Cho hai biến cố độc lập A và B có P(A) = 0,4 và P(B) = 0,5. Xác suất của biến cố giao P(AB) và biến cố hợp P(A $\cup$ B) lần lượt bằng:\par
\begin{itemize}\setlength{\itemsep}{1pt}\setlength{\parskip}{0pt}
  \item[] \textbf{A.} P(AB) = 0,10 và P(A $\cup$ B) = 0,80
  \item[] \textbf{B.} P(AB) = 0,20 và P(A $\cup$ B) = 0,90
  \item[] \textbf{C.} P(AB) = 0,90 và P(A $\cup$ B) = 0,20
  \item[] \textbf{D.} P(AB) = 0,20 và P(A $\cup$ B) = 0,70
\end{itemize}\vspace{2pt}
\textbf{\color{primary} Câu 34.} Hãy tính toán và chọn phương án đúng đối với bài tập: Hai xạ thủ độc lập cùng bắn vào một bia. Xác suất bắn trúng của người thứ nhất là 0,8 và của người thứ hai là 0,7. Xác suất để cả hai người đều bắn trượt bằng:\par
\hspace*{1.5em}\textbf{A.} 0,56\quad\quad \textbf{B.} 0,15\quad\quad \textbf{C.} 0,06\quad\quad \textbf{D.} 0,94\par\vspace{2pt}
\textbf{\color{primary} Câu 35.} Dựa trên các tính chất trọng tâm của Chương VIII, xét bài toán: Cho hai biến cố A, B có P(A) = 0,5; P(B) = 0,6 và P(A $\cup$ B) = 0,8. Giá trị của P(AB) bằng:\par
\hspace*{1.5em}\textbf{A.} 0,2\quad\quad \textbf{B.} 0,1\quad\quad \textbf{C.} 0,3\quad\quad \textbf{D.} 0,4\par\vspace{2pt}
\textbf{\color{primary} Câu 36.} Thực hành giải toán (Các quy tắc tính xác suất): Cho hai biến cố độc lập A và B có P(A) = 1/17 và P(B) = 1/3. Xác suất của biến cố giao P(AB) bằng. Đáp án nào dưới đây thỏa mãn?\par
\hspace*{1.5em}\textbf{A.} 1/51\quad\quad \textbf{B.} 20/51\quad\quad \textbf{C.} 2/51\quad\quad \textbf{D.} 1/20\par\vspace{2pt}
\textbf{\color{primary} Câu 37.} Kiểm tra kỹ năng biến đổi và suy luận qua bài toán: Cho hai biến cố độc lập A và B có P(A) = 0,4 và P(B) = 0,5. Xác suất của biến cố giao P(AB) và biến cố hợp P(A $\cup$ B) lần lượt bằng:\par
\begin{itemize}\setlength{\itemsep}{1pt}\setlength{\parskip}{0pt}
  \item[] \textbf{A.} P(AB) = 0,20 và P(A $\cup$ B) = 0,90
  \item[] \textbf{B.} P(AB) = 0,10 và P(A $\cup$ B) = 0,80
  \item[] \textbf{C.} P(AB) = 0,20 và P(A $\cup$ B) = 0,70
  \item[] \textbf{D.} P(AB) = 0,90 và P(A $\cup$ B) = 0,20
\end{itemize}\vspace{2pt}
\textbf{\color{primary} Câu 38.} Cho giả thiết thuộc phạm vi Chương VIII: Hai xạ thủ độc lập cùng bắn vào một bia. Xác suất bắn trúng của người thứ nhất là 0,8 và của người thứ hai là 0,7. Xác suất để cả hai người đều bắn trượt bằng. Kết luận đúng là:\par
\hspace*{1.5em}\textbf{A.} 0,06\quad\quad \textbf{B.} 0,94\quad\quad \textbf{C.} 0,56\quad\quad \textbf{D.} 0,15\par\vspace{2pt}
\textbf{\color{primary} Câu 39.} Rút gọn hoặc tính giá trị theo yêu cầu sau của Chương VIII: Cho hai biến cố A, B có P(A) = 0,5; P(B) = 0,6 và P(A $\cup$ B) = 0,8. Giá trị của P(AB) bằng:\par
\hspace*{1.5em}\textbf{A.} 0,4\quad\quad \textbf{B.} 0,1\quad\quad \textbf{C.} 0,3\quad\quad \textbf{D.} 0,2\par\vspace{2pt}
\textbf{\color{primary} Câu 40.} Xác định phương án đúng sau khi thực hiện phép tính / suy luận: Cho hai biến cố độc lập A và B có P(A) = 1/21 và P(B) = 1/3. Xác suất của biến cố giao P(AB) bằng:\par
\hspace*{1.5em}\textbf{A.} 1/24\quad\quad \textbf{B.} 2/63\quad\quad \textbf{C.} 1/63\quad\quad \textbf{D.} 24/63\par\vspace{2pt}
\subsection*{\color{accent} Phần III. Mức độ Khó (Vận dụng \& Vận dụng cao -- Câu 41 đến Câu 50)}
\textbf{\color{primary} Câu 41.} Hai động cơ hoạt động độc lập có xác suất hoạt động tốt trong ngày lần lượt là 0,9 và 0,8. Xác suất để có ít nhất một động cơ hoạt động tốt trong ngày bằng:\par
\hspace*{1.5em}\textbf{A.} 0,90\quad\quad \textbf{B.} 0,72\quad\quad \textbf{C.} 0,98\quad\quad \textbf{D.} 0,26\par\vspace{2pt}
\textbf{\color{primary} Câu 42.} Bài toán vận dụng (Chương VIII – Các quy tắc tính xác suất): Một đề trắc nghiệm gồm 5 câu hỏi độc lập, mỗi câu có 4 phương án trả lời trong đó chỉ có 1 phương án đúng. Một học sinh khoanh ngẫu nhiên cả 5 câu. Xác suất để học sinh đó trả lời đúng ít nhất 1 câu là. Kết quả đúng là:\par
\hspace*{1.5em}\textbf{A.} (1/4)$^5$ = 1/1024\quad\quad \textbf{B.} 1 - (1/4)$^5$ = 1023/1024\quad\quad \textbf{C.} 1 - (3/4)$^5$ = 781/1024\quad\quad \textbf{D.} (3/4)$^5$ = 243/1024\par\vspace{2pt}
\textbf{\color{primary} Câu 43.} Vận dụng kiến thức nâng cao trong phạm vi Chương VIII để giải bài toán: Hai thiết bị hoạt động độc lập với nhau, xác suất hoạt động tốt của thiết bị I là 5/6 và của thiết bị II là 2/3. Xác suất để cả hai thiết bị cùng hoạt động tốt là:\par
\hspace*{1.5em}\textbf{A.} 5/18\quad\quad \textbf{B.} 1/18\quad\quad \textbf{C.} 10/18\quad\quad \textbf{D.} 17/18\par\vspace{2pt}
\textbf{\color{primary} Câu 44.} Phân tích và giải bài toán tổng hợp sau của Chương VIII: Hai động cơ hoạt động độc lập có xác suất hoạt động tốt trong ngày lần lượt là 0,9 và 0,8. Xác suất để có ít nhất một động cơ hoạt động tốt trong ngày bằng:\par
\hspace*{1.5em}\textbf{A.} 0,90\quad\quad \textbf{B.} 0,98\quad\quad \textbf{C.} 0,72\quad\quad \textbf{D.} 0,26\par\vspace{2pt}
\textbf{\color{primary} Câu 45.} Thử thách tư duy trong phạm vi kiến thức Chương VIII: Một đề trắc nghiệm gồm 5 câu hỏi độc lập, mỗi câu có 4 phương án trả lời trong đó chỉ có 1 phương án đúng. Một học sinh khoanh ngẫu nhiên cả 5 câu. Xác suất để học sinh đó trả lời đúng ít nhất 1 câu là. Đáp số chính xác là:\par
\hspace*{1.5em}\textbf{A.} 1 - (1/4)$^5$ = 1023/1024\quad\quad \textbf{B.} 1 - (3/4)$^5$ = 781/1024\quad\quad \textbf{C.} (3/4)$^5$ = 243/1024\quad\quad \textbf{D.} (1/4)$^5$ = 1/1024\par\vspace{2pt}
\textbf{\color{primary} Câu 46.} Sử dụng phương pháp giải toán của Chương VIII để tìm đáp án cho bài toán: Hai thiết bị hoạt động độc lập với nhau, xác suất hoạt động tốt của thiết bị I là 8/9 và của thiết bị II là 2/3. Xác suất để cả hai thiết bị cùng hoạt động tốt là:\par
\hspace*{1.5em}\textbf{A.} 1/27\quad\quad \textbf{B.} 16/27\quad\quad \textbf{C.} 8/27\quad\quad \textbf{D.} 26/27\par\vspace{2pt}
\textbf{\color{primary} Câu 47.} Xét bài toán vận dụng thực tiễn / nâng cao thuộc Chương VIII: Hai động cơ hoạt động độc lập có xác suất hoạt động tốt trong ngày lần lượt là 0,9 và 0,8. Xác suất để có ít nhất một động cơ hoạt động tốt trong ngày bằng:\par
\hspace*{1.5em}\textbf{A.} 0,98\quad\quad \textbf{B.} 0,90\quad\quad \textbf{C.} 0,26\quad\quad \textbf{D.} 0,72\par\vspace{2pt}
\textbf{\color{primary} Câu 48.} Tìm kết quả cuối cùng của bài toán phân loại học sinh khá giỏi (Chương VIII): Một đề trắc nghiệm gồm 5 câu hỏi độc lập, mỗi câu có 4 phương án trả lời trong đó chỉ có 1 phương án đúng. Một học sinh khoanh ngẫu nhiên cả 5 câu. Xác suất để học sinh đó trả lời đúng ít nhất 1 câu là:\par
\hspace*{1.5em}\textbf{A.} (3/4)$^5$ = 243/1024\quad\quad \textbf{B.} 1 - (3/4)$^5$ = 781/1024\quad\quad \textbf{C.} (1/4)$^5$ = 1/1024\quad\quad \textbf{D.} 1 - (1/4)$^5$ = 1023/1024\par\vspace{2pt}
\textbf{\color{primary} Câu 49.} Kết hợp các tính chất đã học trong Chương VIII, hãy giải quyết bài toán: Hai thiết bị hoạt động độc lập với nhau, xác suất hoạt động tốt của thiết bị I là 11/12 và của thiết bị II là 2/3. Xác suất để cả hai thiết bị cùng hoạt động tốt là:\par
\hspace*{1.5em}\textbf{A.} 11/36\quad\quad \textbf{B.} 35/36\quad\quad \textbf{C.} 1/36\quad\quad \textbf{D.} 22/36\par\vspace{2pt}
\textbf{\color{primary} Câu 50.} Câu hỏi vận dụng cao cuối Chương VIII: Hai động cơ hoạt động độc lập có xác suất hoạt động tốt trong ngày lần lượt là 0,9 và 0,8. Xác suất để có ít nhất một động cơ hoạt động tốt trong ngày bằng. Phương án đúng là:\par
\hspace*{1.5em}\textbf{A.} 0,72\quad\quad \textbf{B.} 0,26\quad\quad \textbf{C.} 0,98\quad\quad \textbf{D.} 0,90\par\vspace{2pt}
\subsection*{\color{primary} Bảng đáp án Chương VIII: Các quy tắc tính xác suất}
\begin{center}
\renewcommand{\arraystretch}{1.25}
\begin{tabularx}{\textwidth}{*{10}{>{\centering\arraybackslash}X}}
\toprule
\textbf{\color{primary} Cột 1} & \textbf{\color{primary} Cột 2} & \textbf{\color{primary} Cột 3} & \textbf{\color{primary} Cột 4} & \textbf{\color{primary} Cột 5} & \textbf{\color{primary} Cột 6} & \textbf{\color{primary} Cột 7} & \textbf{\color{primary} Cột 8} & \textbf{\color{primary} Cột 9} & \textbf{\color{primary} Cột 10} \\
\midrule
\textbf{1.B} & \textbf{2.D} & \textbf{3.C} & \textbf{4.D} & \textbf{5.A} & \textbf{6.D} & \textbf{7.A} & \textbf{8.C} & \textbf{9.C} & \textbf{10.A} \\
\rowcolor{rowtint} \textbf{11.A} & \textbf{12.A} & \textbf{13.A} & \textbf{14.A} & \textbf{15.D} & \textbf{16.C} & \textbf{17.D} & \textbf{18.C} & \textbf{19.B} & \textbf{20.D} \\
\textbf{21.C} & \textbf{22.A} & \textbf{23.D} & \textbf{24.A} & \textbf{25.C} & \textbf{26.D} & \textbf{27.D} & \textbf{28.A} & \textbf{29.B} & \textbf{30.A} \\
\rowcolor{rowtint} \textbf{31.A} & \textbf{32.D} & \textbf{33.D} & \textbf{34.C} & \textbf{35.C} & \textbf{36.A} & \textbf{37.C} & \textbf{38.A} & \textbf{39.C} & \textbf{40.C} \\
\textbf{41.C} & \textbf{42.C} & \textbf{43.C} & \textbf{44.B} & \textbf{45.B} & \textbf{46.B} & \textbf{47.A} & \textbf{48.B} & \textbf{49.D} & \textbf{50.C} \\
\bottomrule
\end{tabularx}
\end{center}\vspace{10pt}
\section*{\color{primary} CHƯƠNG IX: ĐẠO HÀM}
\textit{Trọng tâm kiến thức SGK: Định nghĩa và ý nghĩa hình học/vật lí của đạo hàm, các quy tắc tính đạo hàm, đạo hàm của hàm hợp và đạo hàm cấp hai.}\par\vspace{4pt}
\subsection*{\color{accent} Phần I. Mức độ Dễ (Nhận biết -- Câu 1 đến Câu 20)}
\textbf{\color{primary} Câu 1.} Đạo hàm của hàm số đa thức y = x$^3$ - 2x$^2$ + 5x - 3 là:\par
\hspace*{1.5em}\textbf{A.} y' = 3x$^2$ - 4x + 5\quad\quad \textbf{B.} y' = 3x$^2$ - 2x + 5\quad\quad \textbf{C.} y' = x$^2$ - 4x + 5\quad\quad \textbf{D.} y' = 3x$^2$ - 4x + 2\par\vspace{2pt}
\textbf{\color{primary} Câu 2.} Thực hiện yêu cầu nhận biết sau: Đạo hàm của hàm số lượng giác y = sin x + cos x là.\par
\hspace*{1.5em}\textbf{A.} y' = cos x - sin x\quad\quad \textbf{B.} y' = sin x - cos x\quad\quad \textbf{C.} y' = cos x + sin x\quad\quad \textbf{D.} y' = -cos x - sin x\par\vspace{2pt}
\textbf{\color{primary} Câu 3.} Trong bài học thuộc Chương IX, xét nội dung: Cho hàm số f(x) = x$^2$ + 6x. Hệ số góc của tiếp tuyến của đồ thị hàm số tại điểm có hoành độ x$_0$ = 1 bằng. Đáp án đúng là:\par
\hspace*{1.5em}\textbf{A.} 6\quad\quad \textbf{B.} 8\quad\quad \textbf{C.} 7\quad\quad \textbf{D.} 2\par\vspace{2pt}
\textbf{\color{primary} Câu 4.} Dựa vào kiến thức chuẩn SGK của Chương IX (Đạo hàm), đạo hàm của hàm số đa thức y = x$^3$ - 5x$^2$ + 5x - 3 là:\par
\hspace*{1.5em}\textbf{A.} y' = 3x$^2$ - 10x + 5\quad\quad \textbf{B.} y' = 3x$^2$ - 5x + 5\quad\quad \textbf{C.} y' = 3x$^2$ - 10x + 2\quad\quad \textbf{D.} y' = x$^2$ - 10x + 5\par\vspace{2pt}
\textbf{\color{primary} Câu 5.} Chọn kết quả đúng nhất khi xét: Đạo hàm của hàm số lượng giác y = sin x + cos x là:\par
\hspace*{1.5em}\textbf{A.} y' = -cos x - sin x\quad\quad \textbf{B.} y' = cos x - sin x\quad\quad \textbf{C.} y' = cos x + sin x\quad\quad \textbf{D.} y' = sin x - cos x\par\vspace{2pt}
\textbf{\color{primary} Câu 6.} Khi làm bài tập cơ bản về đạo hàm, một học sinh gặp câu hỏi: Cho hàm số f(x) = x$^2$ + 12x. Hệ số góc của tiếp tuyến của đồ thị hàm số tại điểm có hoành độ x$_0$ = 1 bằng. Kết quả nào dưới đây đúng?\par
\hspace*{1.5em}\textbf{A.} 2\quad\quad \textbf{B.} 12\quad\quad \textbf{C.} 14\quad\quad \textbf{D.} 13\par\vspace{2pt}
\textbf{\color{primary} Câu 7.} Xác định giá trị hoặc khẳng định đúng cho bài toán sau: Đạo hàm của hàm số đa thức y = x$^3$ - 8x$^2$ + 5x - 3 là:\par
\hspace*{1.5em}\textbf{A.} y' = 3x$^2$ - 16x + 5\quad\quad \textbf{B.} y' = 3x$^2$ - 8x + 5\quad\quad \textbf{C.} y' = 3x$^2$ - 16x + 2\quad\quad \textbf{D.} y' = x$^2$ - 16x + 5\par\vspace{2pt}
\textbf{\color{primary} Câu 8.} Áp dụng trực tiếp định nghĩa và quy tắc trong Chương IX: Đạo hàm của hàm số lượng giác y = sin x + cos x là:\par
\hspace*{1.5em}\textbf{A.} y' = sin x - cos x\quad\quad \textbf{B.} y' = -cos x - sin x\quad\quad \textbf{C.} y' = cos x + sin x\quad\quad \textbf{D.} y' = cos x - sin x\par\vspace{2pt}
\textbf{\color{primary} Câu 9.} Kiểm tra nhanh kiến thức nền tảng: Cho hàm số f(x) = x$^2$ + 18x. Hệ số góc của tiếp tuyến của đồ thị hàm số tại điểm có hoành độ x$_0$ = 1 bằng:\par
\hspace*{1.5em}\textbf{A.} 18\quad\quad \textbf{B.} 20\quad\quad \textbf{C.} 2\quad\quad \textbf{D.} 19\par\vspace{2pt}
\textbf{\color{primary} Câu 10.} Cho dữ kiện nhận biết sau đây: Đạo hàm của hàm số đa thức y = x$^3$ - 11x$^2$ + 5x - 3 là. Phương án chính xác là:\par
\hspace*{1.5em}\textbf{A.} y' = x$^2$ - 22x + 5\quad\quad \textbf{B.} y' = 3x$^2$ - 22x + 5\quad\quad \textbf{C.} y' = 3x$^2$ - 11x + 5\quad\quad \textbf{D.} y' = 3x$^2$ - 22x + 2\par\vspace{2pt}
\textbf{\color{primary} Câu 11.} Hãy chỉ ra kết quả đúng trong bốn lựa chọn A, B, C, D đối với yêu cầu: Đạo hàm của hàm số lượng giác y = sin x + cos x là:\par
\hspace*{1.5em}\textbf{A.} y' = sin x - cos x\quad\quad \textbf{B.} y' = -cos x - sin x\quad\quad \textbf{C.} y' = cos x + sin x\quad\quad \textbf{D.} y' = cos x - sin x\par\vspace{2pt}
\textbf{\color{primary} Câu 12.} Câu hỏi ôn tập lý thuyết và kỹ năng cơ bản (Chương IX): Cho hàm số f(x) = x$^2$ + 24x. Hệ số góc của tiếp tuyến của đồ thị hàm số tại điểm có hoành độ x$_0$ = 1 bằng:\par
\hspace*{1.5em}\textbf{A.} 2\quad\quad \textbf{B.} 26\quad\quad \textbf{C.} 25\quad\quad \textbf{D.} 24\par\vspace{2pt}
\textbf{\color{primary} Câu 13.} Theo nội dung sách giáo khoa Kết nối tri thức với cuộc sống (Chương IX), đạo hàm của hàm số đa thức y = x$^3$ - 14x$^2$ + 5x - 3 là:\par
\hspace*{1.5em}\textbf{A.} y' = 3x$^2$ - 28x + 5\quad\quad \textbf{B.} y' = 3x$^2$ - 28x + 2\quad\quad \textbf{C.} y' = 3x$^2$ - 14x + 5\quad\quad \textbf{D.} y' = x$^2$ - 28x + 5\par\vspace{2pt}
\textbf{\color{primary} Câu 14.} Tìm đáp số đúng của bài toán cơ bản sau: Đạo hàm của hàm số lượng giác y = sin x + cos x là:\par
\hspace*{1.5em}\textbf{A.} y' = -cos x - sin x\quad\quad \textbf{B.} y' = cos x + sin x\quad\quad \textbf{C.} y' = sin x - cos x\quad\quad \textbf{D.} y' = cos x - sin x\par\vspace{2pt}
\textbf{\color{primary} Câu 15.} Xét đối tượng toán học trong Chương IX: Cho hàm số f(x) = x$^2$ + 30x. Hệ số góc của tiếp tuyến của đồ thị hàm số tại điểm có hoành độ x$_0$ = 1 bằng. Giá trị hoặc tính chất tương ứng là:\par
\hspace*{1.5em}\textbf{A.} 2\quad\quad \textbf{B.} 31\quad\quad \textbf{C.} 30\quad\quad \textbf{D.} 32\par\vspace{2pt}
\textbf{\color{primary} Câu 16.} Hoàn thành phép xác định cơ bản sau đây: Đạo hàm của hàm số đa thức y = x$^3$ - 17x$^2$ + 5x - 3 là:\par
\hspace*{1.5em}\textbf{A.} y' = 3x$^2$ - 34x + 2\quad\quad \textbf{B.} y' = 3x$^2$ - 34x + 5\quad\quad \textbf{C.} y' = x$^2$ - 34x + 5\quad\quad \textbf{D.} y' = 3x$^2$ - 17x + 5\par\vspace{2pt}
\textbf{\color{primary} Câu 17.} Trong các khẳng định hoặc giá trị dưới đây, đâu là đáp án đúng cho câu hỏi: Đạo hàm của hàm số lượng giác y = sin x + cos x là?\par
\hspace*{1.5em}\textbf{A.} y' = cos x - sin x\quad\quad \textbf{B.} y' = cos x + sin x\quad\quad \textbf{C.} y' = sin x - cos x\quad\quad \textbf{D.} y' = -cos x - sin x\par\vspace{2pt}
\textbf{\color{primary} Câu 18.} Bài tập khởi động Chương IX: Cho hàm số f(x) = x$^2$ + 36x. Hệ số góc của tiếp tuyến của đồ thị hàm số tại điểm có hoành độ x$_0$ = 1 bằng:\par
\hspace*{1.5em}\textbf{A.} 36\quad\quad \textbf{B.} 37\quad\quad \textbf{C.} 2\quad\quad \textbf{D.} 38\par\vspace{2pt}
\textbf{\color{primary} Câu 19.} Quan sát dữ kiện đề bài: Đạo hàm của hàm số đa thức y = x$^3$ - 20x$^2$ + 5x - 3 là. Hãy chọn phương án đúng:\par
\hspace*{1.5em}\textbf{A.} y' = x$^2$ - 40x + 5\quad\quad \textbf{B.} y' = 3x$^2$ - 40x + 5\quad\quad \textbf{C.} y' = 3x$^2$ - 40x + 2\quad\quad \textbf{D.} y' = 3x$^2$ - 20x + 5\par\vspace{2pt}
\textbf{\color{primary} Câu 20.} Sử dụng công thức hoặc tính chất đầu bài của Chương IX để giải quyết yêu cầu: Đạo hàm của hàm số lượng giác y = sin x + cos x là:\par
\hspace*{1.5em}\textbf{A.} y' = -cos x - sin x\quad\quad \textbf{B.} y' = cos x + sin x\quad\quad \textbf{C.} y' = sin x - cos x\quad\quad \textbf{D.} y' = cos x - sin x\par\vspace{2pt}
\subsection*{\color{accent} Phần II. Mức độ Trung bình (Thông hiểu -- Câu 21 đến Câu 40)}
\textbf{\color{primary} Câu 21.} Đạo hàm cấp hai f''(x) của hàm số f(x) = x$^4$ - 2x$^2$ + 7x tại x = 1 bằng:\par
\hspace*{1.5em}\textbf{A.} 8\quad\quad \textbf{B.} 12\quad\quad \textbf{C.} 7\quad\quad \textbf{D.} 10\par\vspace{2pt}
\textbf{\color{primary} Câu 22.} Vận dụng quy tắc trong Chương IX (Đạo hàm) để giải bài toán: Đạo hàm của hàm phân thức y = (2x + 1)/(x - 3) (với x $\ne$ 3) là:\par
\hspace*{1.5em}\textbf{A.} y' = 7 / (x - 3)$^2$\quad\quad \textbf{B.} y' = -7 / (x - 3)$^2$\quad\quad \textbf{C.} y' = 2\quad\quad \textbf{D.} y' = -1 / (x - 3)$^2$\par\vspace{2pt}
\textbf{\color{primary} Câu 23.} Biến đổi và tính toán theo kiến thức Chương IX: Một chất điểm chuyển động thẳng có phương trình quãng đường s(t) = t$^3$ - 3t$^2$ + 18t (m), với t tính bằng giây. Vận tốc tức thời v(t) = s'(t) của chất điểm tại thời điểm t = 2 (s) bằng. Kết quả thu được là:\par
\hspace*{1.5em}\textbf{A.} 24 m/s\quad\quad \textbf{B.} 32 m/s\quad\quad \textbf{C.} 6 m/s\quad\quad \textbf{D.} 18 m/s\par\vspace{2pt}
\textbf{\color{primary} Câu 24.} Xét bài toán mức độ thông hiểu sau: Đạo hàm cấp hai f''(x) của hàm số f(x) = x$^4$ - 5x$^2$ + 7x tại x = 1 bằng. Đáp án đúng là:\par
\hspace*{1.5em}\textbf{A.} 12\quad\quad \textbf{B.} 7\quad\quad \textbf{C.} 1\quad\quad \textbf{D.} 2\par\vspace{2pt}
\textbf{\color{primary} Câu 25.} Để hoàn thành phiếu học tập Chương IX, hãy xác định: Đạo hàm của hàm phân thức y = (2x + 1)/(x - 6) (với x $\ne$ 6) là:\par
\hspace*{1.5em}\textbf{A.} y' = 2\quad\quad \textbf{B.} y' = -1 / (x - 6)$^2$\quad\quad \textbf{C.} y' = -13 / (x - 6)$^2$\quad\quad \textbf{D.} y' = 13 / (x - 6)$^2$\par\vspace{2pt}
\textbf{\color{primary} Câu 26.} Phân tích dữ kiện bài toán: Một chất điểm chuyển động thẳng có phương trình quãng đường s(t) = t$^3$ - 3t$^2$ + 30t (m), với t tính bằng giây. Vận tốc tức thời v(t) = s'(t) của chất điểm tại thời điểm t = 2 (s) bằng. Giá trị cần tìm bằng:\par
\hspace*{1.5em}\textbf{A.} 18 m/s\quad\quad \textbf{B.} 30 m/s\quad\quad \textbf{C.} 56 m/s\quad\quad \textbf{D.} 36 m/s\par\vspace{2pt}
\textbf{\color{primary} Câu 27.} Áp dụng tính chất đã học ở Chương IX vào bài tập sau: Đạo hàm cấp hai f''(x) của hàm số f(x) = x$^4$ - 8x$^2$ + 7x tại x = 1 bằng:\par
\hspace*{1.5em}\textbf{A.} -5\quad\quad \textbf{B.} -4\quad\quad \textbf{C.} 12\quad\quad \textbf{D.} 4\par\vspace{2pt}
\textbf{\color{primary} Câu 28.} Thực hiện các bước suy luận hợp lý để giải quyết yêu cầu: Đạo hàm của hàm phân thức y = (2x + 1)/(x - 9) (với x $\ne$ 9) là:\par
\hspace*{1.5em}\textbf{A.} y' = -1 / (x - 9)$^2$\quad\quad \textbf{B.} y' = 19 / (x - 9)$^2$\quad\quad \textbf{C.} y' = -19 / (x - 9)$^2$\quad\quad \textbf{D.} y' = 2\par\vspace{2pt}
\textbf{\color{primary} Câu 29.} Cho bài toán tự luyện (Đạo hàm): Một chất điểm chuyển động thẳng có phương trình quãng đường s(t) = t$^3$ - 3t$^2$ + 42t (m), với t tính bằng giây. Vận tốc tức thời v(t) = s'(t) của chất điểm tại thời điểm t = 2 (s) bằng. Phương án nào sau đây chính xác?\par
\hspace*{1.5em}\textbf{A.} 80 m/s\quad\quad \textbf{B.} 30 m/s\quad\quad \textbf{C.} 42 m/s\quad\quad \textbf{D.} 48 m/s\par\vspace{2pt}
\textbf{\color{primary} Câu 30.} Kết hợp các phép biến đổi trong Chương IX, hãy tìm kết quả của bài toán: Đạo hàm cấp hai f''(x) của hàm số f(x) = x$^4$ - 11x$^2$ + 7x tại x = 1 bằng:\par
\hspace*{1.5em}\textbf{A.} 1\quad\quad \textbf{B.} -11\quad\quad \textbf{C.} -10\quad\quad \textbf{D.} 12\par\vspace{2pt}
\textbf{\color{primary} Câu 31.} Một bài kiểm tra định kỳ nội dung Chương IX nêu câu hỏi: Đạo hàm của hàm phân thức y = (2x + 1)/(x - 12) (với x $\ne$ 12) là. Hãy chọn đáp số đúng:\par
\hspace*{1.5em}\textbf{A.} y' = 2\quad\quad \textbf{B.} y' = -25 / (x - 12)$^2$\quad\quad \textbf{C.} y' = 25 / (x - 12)$^2$\quad\quad \textbf{D.} y' = -1 / (x - 12)$^2$\par\vspace{2pt}
\textbf{\color{primary} Câu 32.} Tìm lời giải đúng cho bài toán thông hiểu: Một chất điểm chuyển động thẳng có phương trình quãng đường s(t) = t$^3$ - 3t$^2$ + 54t (m), với t tính bằng giây. Vận tốc tức thời v(t) = s'(t) của chất điểm tại thời điểm t = 2 (s) bằng:\par
\hspace*{1.5em}\textbf{A.} 60 m/s\quad\quad \textbf{B.} 54 m/s\quad\quad \textbf{C.} 104 m/s\quad\quad \textbf{D.} 42 m/s\par\vspace{2pt}
\textbf{\color{primary} Câu 33.} Sử dụng mối liên hệ giữa các đại lượng trong Chương IX: Đạo hàm cấp hai f''(x) của hàm số f(x) = x$^4$ - 14x$^2$ + 7x tại x = 1 bằng:\par
\hspace*{1.5em}\textbf{A.} 12\quad\quad \textbf{B.} -2\quad\quad \textbf{C.} -16\quad\quad \textbf{D.} -17\par\vspace{2pt}
\textbf{\color{primary} Câu 34.} Hãy tính toán và chọn phương án đúng đối với bài tập: Đạo hàm của hàm phân thức y = (2x + 1)/(x - 15) (với x $\ne$ 15) là:\par
\hspace*{1.5em}\textbf{A.} y' = -31 / (x - 15)$^2$\quad\quad \textbf{B.} y' = -1 / (x - 15)$^2$\quad\quad \textbf{C.} y' = 2\quad\quad \textbf{D.} y' = 31 / (x - 15)$^2$\par\vspace{2pt}
\textbf{\color{primary} Câu 35.} Dựa trên các tính chất trọng tâm của Chương IX, xét bài toán: Một chất điểm chuyển động thẳng có phương trình quãng đường s(t) = t$^3$ - 3t$^2$ + 66t (m), với t tính bằng giây. Vận tốc tức thời v(t) = s'(t) của chất điểm tại thời điểm t = 2 (s) bằng:\par
\hspace*{1.5em}\textbf{A.} 72 m/s\quad\quad \textbf{B.} 54 m/s\quad\quad \textbf{C.} 66 m/s\quad\quad \textbf{D.} 128 m/s\par\vspace{2pt}
\textbf{\color{primary} Câu 36.} Thực hành giải toán (Đạo hàm): Đạo hàm cấp hai f''(x) của hàm số f(x) = x$^4$ - 17x$^2$ + 7x tại x = 1 bằng. Đáp án nào dưới đây thỏa mãn?\par
\hspace*{1.5em}\textbf{A.} -5\quad\quad \textbf{B.} 12\quad\quad \textbf{C.} -23\quad\quad \textbf{D.} -22\par\vspace{2pt}
\textbf{\color{primary} Câu 37.} Kiểm tra kỹ năng biến đổi và suy luận qua bài toán: Đạo hàm của hàm phân thức y = (2x + 1)/(x - 18) (với x $\ne$ 18) là:\par
\hspace*{1.5em}\textbf{A.} y' = -1 / (x - 18)$^2$\quad\quad \textbf{B.} y' = 37 / (x - 18)$^2$\quad\quad \textbf{C.} y' = 2\quad\quad \textbf{D.} y' = -37 / (x - 18)$^2$\par\vspace{2pt}
\textbf{\color{primary} Câu 38.} Cho giả thiết thuộc phạm vi Chương IX: Một chất điểm chuyển động thẳng có phương trình quãng đường s(t) = t$^3$ - 3t$^2$ + 78t (m), với t tính bằng giây. Vận tốc tức thời v(t) = s'(t) của chất điểm tại thời điểm t = 2 (s) bằng. Kết luận đúng là:\par
\hspace*{1.5em}\textbf{A.} 152 m/s\quad\quad \textbf{B.} 78 m/s\quad\quad \textbf{C.} 66 m/s\quad\quad \textbf{D.} 84 m/s\par\vspace{2pt}
\textbf{\color{primary} Câu 39.} Rút gọn hoặc tính giá trị theo yêu cầu sau của Chương IX: Đạo hàm cấp hai f''(x) của hàm số f(x) = x$^4$ - 20x$^2$ + 7x tại x = 1 bằng:\par
\hspace*{1.5em}\textbf{A.} -28\quad\quad \textbf{B.} -29\quad\quad \textbf{C.} 12\quad\quad \textbf{D.} -8\par\vspace{2pt}
\textbf{\color{primary} Câu 40.} Xác định phương án đúng sau khi thực hiện phép tính / suy luận: Đạo hàm của hàm phân thức y = (2x + 1)/(x - 21) (với x $\ne$ 21) là:\par
\hspace*{1.5em}\textbf{A.} y' = -1 / (x - 21)$^2$\quad\quad \textbf{B.} y' = -43 / (x - 21)$^2$\quad\quad \textbf{C.} y' = 43 / (x - 21)$^2$\quad\quad \textbf{D.} y' = 2\par\vspace{2pt}
\subsection*{\color{accent} Phần III. Mức độ Khó (Vận dụng \& Vận dụng cao -- Câu 41 đến Câu 50)}
\textbf{\color{primary} Câu 41.} Phương trình tiếp tuyến của đồ thị hàm số y = x$^3$ - 3x + 2 tại điểm có hoành độ x$_0$ = 2 là:\par
\hspace*{1.5em}\textbf{A.} y = 9x + 20\quad\quad \textbf{B.} y = 9x - 18\quad\quad \textbf{C.} y = 9x - 16\quad\quad \textbf{D.} y = 3x - 4\par\vspace{2pt}
\textbf{\color{primary} Câu 42.} Bài toán vận dụng (Chương IX – Đạo hàm): Một vật chuyển động theo quy luật s(t) = -t$^3$ + 6t$^2$ + 9t (m). Vận tốc lớn nhất mà vật đạt được trong quá trình chuyển động bằng. Kết quả đúng là:\par
\hspace*{1.5em}\textbf{A.} 27 m/s\quad\quad \textbf{B.} 21 m/s (tại t = 2 s)\quad\quad \textbf{C.} 9 m/s\quad\quad \textbf{D.} 15 m/s\par\vspace{2pt}
\textbf{\color{primary} Câu 43.} Vận dụng kiến thức nâng cao trong phạm vi Chương IX để giải bài toán: Phương trình tiếp tuyến của đồ thị hàm số y = x$^3$ - 3x + 4 tại điểm có hoành độ x$_0$ = 2 là:\par
\hspace*{1.5em}\textbf{A.} y = 9x - 16\quad\quad \textbf{B.} y = 9x + 22\quad\quad \textbf{C.} y = 3x - 2\quad\quad \textbf{D.} y = 9x - 14\par\vspace{2pt}
\textbf{\color{primary} Câu 44.} Phân tích và giải bài toán tổng hợp sau của Chương IX: Một vật chuyển động theo quy luật s(t) = -t$^3$ + 6t$^2$ + 9t (m). Vận tốc lớn nhất mà vật đạt được trong quá trình chuyển động bằng:\par
\hspace*{1.5em}\textbf{A.} 27 m/s\quad\quad \textbf{B.} 9 m/s\quad\quad \textbf{C.} 21 m/s (tại t = 2 s)\quad\quad \textbf{D.} 15 m/s\par\vspace{2pt}
\textbf{\color{primary} Câu 45.} Thử thách tư duy trong phạm vi kiến thức Chương IX: Phương trình tiếp tuyến của đồ thị hàm số y = x$^3$ - 3x + 6 tại điểm có hoành độ x$_0$ = 2 là. Đáp số chính xác là:\par
\hspace*{1.5em}\textbf{A.} y = 9x - 14\quad\quad \textbf{B.} y = 3x - 0\quad\quad \textbf{C.} y = 9x - 12\quad\quad \textbf{D.} y = 9x + 24\par\vspace{2pt}
\textbf{\color{primary} Câu 46.} Sử dụng phương pháp giải toán của Chương IX để tìm đáp án cho bài toán: Một vật chuyển động theo quy luật s(t) = -t$^3$ + 6t$^2$ + 9t (m). Vận tốc lớn nhất mà vật đạt được trong quá trình chuyển động bằng:\par
\hspace*{1.5em}\textbf{A.} 9 m/s\quad\quad \textbf{B.} 15 m/s\quad\quad \textbf{C.} 21 m/s (tại t = 2 s)\quad\quad \textbf{D.} 27 m/s\par\vspace{2pt}
\textbf{\color{primary} Câu 47.} Xét bài toán vận dụng thực tiễn / nâng cao thuộc Chương IX: Phương trình tiếp tuyến của đồ thị hàm số y = x$^3$ - 3x + 8 tại điểm có hoành độ x$_0$ = 2 là:\par
\hspace*{1.5em}\textbf{A.} y = 3x - -2\quad\quad \textbf{B.} y = 9x - 12\quad\quad \textbf{C.} y = 9x - 10\quad\quad \textbf{D.} y = 9x + 26\par\vspace{2pt}
\textbf{\color{primary} Câu 48.} Tìm kết quả cuối cùng của bài toán phân loại học sinh khá giỏi (Chương IX): Một vật chuyển động theo quy luật s(t) = -t$^3$ + 6t$^2$ + 9t (m). Vận tốc lớn nhất mà vật đạt được trong quá trình chuyển động bằng:\par
\hspace*{1.5em}\textbf{A.} 21 m/s (tại t = 2 s)\quad\quad \textbf{B.} 27 m/s\quad\quad \textbf{C.} 15 m/s\quad\quad \textbf{D.} 9 m/s\par\vspace{2pt}
\textbf{\color{primary} Câu 49.} Kết hợp các tính chất đã học trong Chương IX, hãy giải quyết bài toán: Phương trình tiếp tuyến của đồ thị hàm số y = x$^3$ - 3x + 10 tại điểm có hoành độ x$_0$ = 2 là:\par
\hspace*{1.5em}\textbf{A.} y = 9x - 10\quad\quad \textbf{B.} y = 9x + 28\quad\quad \textbf{C.} y = 9x - 8\quad\quad \textbf{D.} y = 3x - -4\par\vspace{2pt}
\textbf{\color{primary} Câu 50.} Câu hỏi vận dụng cao cuối Chương IX: Một vật chuyển động theo quy luật s(t) = -t$^3$ + 6t$^2$ + 9t (m). Vận tốc lớn nhất mà vật đạt được trong quá trình chuyển động bằng. Phương án đúng là:\par
\hspace*{1.5em}\textbf{A.} 15 m/s\quad\quad \textbf{B.} 21 m/s (tại t = 2 s)\quad\quad \textbf{C.} 9 m/s\quad\quad \textbf{D.} 27 m/s\par\vspace{2pt}
\subsection*{\color{primary} Bảng đáp án Chương IX: Đạo hàm}
\begin{center}
\renewcommand{\arraystretch}{1.25}
\begin{tabularx}{\textwidth}{*{10}{>{\centering\arraybackslash}X}}
\toprule
\textbf{\color{primary} Cột 1} & \textbf{\color{primary} Cột 2} & \textbf{\color{primary} Cột 3} & \textbf{\color{primary} Cột 4} & \textbf{\color{primary} Cột 5} & \textbf{\color{primary} Cột 6} & \textbf{\color{primary} Cột 7} & \textbf{\color{primary} Cột 8} & \textbf{\color{primary} Cột 9} & \textbf{\color{primary} Cột 10} \\
\midrule
\textbf{1.A} & \textbf{2.A} & \textbf{3.B} & \textbf{4.A} & \textbf{5.B} & \textbf{6.C} & \textbf{7.A} & \textbf{8.D} & \textbf{9.B} & \textbf{10.B} \\
\rowcolor{rowtint} \textbf{11.D} & \textbf{12.B} & \textbf{13.A} & \textbf{14.D} & \textbf{15.D} & \textbf{16.B} & \textbf{17.A} & \textbf{18.D} & \textbf{19.B} & \textbf{20.D} \\
\textbf{21.A} & \textbf{22.B} & \textbf{23.D} & \textbf{24.D} & \textbf{25.C} & \textbf{26.B} & \textbf{27.B} & \textbf{28.C} & \textbf{29.C} & \textbf{30.C} \\
\rowcolor{rowtint} \textbf{31.B} & \textbf{32.B} & \textbf{33.C} & \textbf{34.A} & \textbf{35.C} & \textbf{36.D} & \textbf{37.D} & \textbf{38.B} & \textbf{39.A} & \textbf{40.B} \\
\textbf{41.C} & \textbf{42.B} & \textbf{43.D} & \textbf{44.C} & \textbf{45.C} & \textbf{46.C} & \textbf{47.C} & \textbf{48.A} & \textbf{49.C} & \textbf{50.B} \\
\bottomrule
\end{tabularx}
\end{center}\vspace{10pt}
\end{document}
